\documentclass[withindex,glossary,times]{cam-thesis}

% Citations using numbers
\usepackage[numbers]{natbib}
\usepackage{amsmath,amssymb,amsthm,amsfonts}
\usepackage{zref-clever}
\zcsetup{cap=true}
\usepackage{stmaryrd}
\usepackage{tikz-cd}
\usepackage{quiver}
\usepackage{CJKutf8}
\usepackage{mathtools}
\usepackage{enumitem}
\setlist[enumerate,1]{label={(\arabic*)}}
\usepackage{stix2}
\usepackage{anyfontsize}
\usepackage{bbold}
\usepackage{ebproof}
\usepackage{etoolbox}
\usepackage{stix}

\tikzcdset{pullback/.style = {"\lrcorner"{anchor = center, pos = 0.125}, draw = none}}

\newcommand{\mc}[1]{\mathcal{#1}}
\newcommand{\mb}[1]{\mathbf{#1}}
\newcommand{\mbb}[1]{\mathbb{#1}}
\newcommand{\gn}[1]{\ulcorner #1 \urcorner}
\newcommand{\mr}[1]{\mathrm{#1}}
\newcommand{\mf}[1]{\mathfrak{#1}}
\newcommand{\ms}[1]{\mathsf{#1}}
\newcommand{\Ind}{\mathbf{Ind}}
\newcommand{\Pro}{\mathbf{Pro}}
\newcommand{\Top}{\mf{Top}}
\newcommand{\Set}{\mb{Set}}
\newcommand{\Prop}{\mb{Prop}}
\newcommand{\sSet}{\mb{sSet}}
\newcommand{\sSeto}{\mb{sSet}_{\le 1}}
\newcommand{\End}{\operatorname{End}}
\newcommand{\Hom}{\operatorname{Hom}}
\newcommand{\Fin}{\mb{Fin}}
\newcommand{\Cnt}{\mb{Cnt}}
\newcommand{\Enm}{\mb{Enm}}
\newcommand{\Grp}{\mb{Grp}}
\newcommand{\sub}{\mr{Sub}}
\newcommand{\Lex}{\mf{Lex}}
\newcommand{\Coh}{\mf{Coh}}
\newcommand{\Pos}{\mb{Pos}}
\newcommand{\alg}{\text{-}\mb{Alg}}
\newcommand{\Var}{\mb{Var}}
\newcommand{\Aff}{\mb{Aff}}
\newcommand{\Vect}{\mb{Vect}}
\newcommand{\CRing}{\mb{CRing}}
\newcommand{\DL}{\mb{DL}}
\newcommand{\BA}{\mb{BA}}
\newcommand{\HA}{\mb{HA}}
\newcommand{\JL}{\mb{JL}}
\newcommand{\ML}{\mb{ML}}
\newcommand{\ProMon}{\mb{ProMod}}
\newcommand{\CTopoi}{\mb{CTopoi}}
\newcommand{\PTc}{\mb{PT}_c}
\newcommand{\Topoi}{\mb{Topoi}}
\newcommand{\sh}{\mb{Sh}}
\newcommand{\psh}{\mb{Psh}}
\newcommand{\Cont}{\mb{Cont}}
\newcommand{\op}{^{\mathrm{op}}}
\newcommand{\inv}{^{\mathrm{-1}}}
\newcommand{\pf}[1]{\widehat{#1}}
\newcommand{\qsi}[1]{\tilde{#1}}
\newcommand{\cob}{\vartriangleleft}
\newcommand{\other}{\mathrm{otherwise}}
\newcommand{\geo}[1]{|#1|}
\newcommand{\ov}[1]{\overline{#1}}
\newcommand{\set}[1]{\{\,#1\,\}}
\newcommand{\eff}{\Leftrightarrow}
\newcommand{\conjt}{\;\&\;}
\newcommand{\pair}[1]{\left\langle#1\right\rangle}
\newcommand{\id}{\mathrm{id}}
\newcommand{\ev}{\mathrm{ev}}
\newcommand{\elem}{\int\!\!}
\newcommand{\nt}{\Rightarrow}
\newcommand{\scomp}[2]{\{\,#1\mid#2\,\}}
\newcommand{\surj}{\twoheadrightarrow}
\newcommand{\inj}{\rightarrowtail}
\newcommand{\hook}{\hookrightarrow}
\newcommand{\prth}[1]{\left(#1\right)}
\newcommand{\fn}{_\omega}
\newcommand{\fp}{_\omega}
\newcommand{\fpp}{_{\omega,+}}
\newcommand{\cp}{_{\omega_1}}
\newcommand{\cpp}{_{\omega_1,+}}
\newcommand{\dv}{\operatorname{\uparrow}}
\newcommand{\cv}{\operatorname{\downarrow}}
\newcommand{\et}{_{\text{\'et}}}
\newcommand{\N}{\mbb N}
\newcommand{\Q}{\mbb Q}
\newcommand{\Z}{\mbb Z}
\newcommand{\T}{\mbb T}
\newcommand{\A}{\mbb A}
\newcommand{\gr}[2]{[#1|#2]}
\newcommand{\fa}[2]{\forall #1\!\in\!#2.\ }
\newcommand{\ex}[2]{\exists #1\!\in\!#2.\ }
\newcommand{\exu}[2]{\exists_! #1\!\in\!#2.\ }
\newcommand{\ld}[2]{\lambda #1\!\in\!\!#2.\ }
\newcommand{\emp}{\emptyset}
\newcommand{\eq}{\leftrightarrow}
\newcommand{\ass}[1]{\llbracket#1\rrbracket} 
\newcommand{\pss}[1]{\left\|#1\right\|}
\newcommand{\mmod}[1]{#1\text{-}\mathbf{Mod}}
\newcommand{\tm}[1]{#1\text{-}\mathrm{Term}}
\newcommand{\eqn}[1]{#1\text{-}\mathrm{Eqn}}
\newcommand{\horn}[1]{#1\text{-}\mathrm{Horn}}
\newcommand{\upp}{_{\ms U}}
\newcommand{\func}{\mb{Func}}
\newcommand{\spec}{\operatorname{Spec}}
\newcommand{\El}{\mr{El}}
\newcommand{\lan}{\ms{lan}}
\newcommand{\ran}{\ms{ran}}
\newcommand{\VT}{V_\T\hspace{-.4ex}\ltimes\hspace{-.4ex}\T}
\newcommand{\Nb}{\mathbf N}
\newcommand{\NA}{\Nb\hspace{-.4ex}\ltimes\hspace{-.4ex}\A}
\newcommand{\uv}[1]{\underline{#1}}
\newcommand{\gp}{\ms{Grpd}}
\newcommand{\hp}{\text{-}}
\newcommand{\st}{\ms{Set}}
\newcommand{\dom}{\operatorname{dom}}
\newcommand{\cod}{\operatorname{cod}}
\newcommand{\bx}{\mdwhtsquare}
\renewcommand{\models}{\vDash}

\newcommand{\buq}[2]{\forall #1\!\!<\!\!#2.\,}
\newcommand{\beq}[2]{\exists #1\!\!<\!\!#2.\,}
\newcommand{\muq}[2]{\mu #1\!\!<\!\!#2.\,}

\newcommand\istsym{\ms{t}}
\newcommand\isfsym{\ms{f}}
\newcommand\ist[1]{\istsym(#1)}
\newcommand\isf[1]{\isfsym(#1)}
\newcommand{\pt}{\ms{pt}}
\newcommand{\tp}{\ms{Type}}
\newcommand{\pp}{\mb{Prop}}
\newcommand{\stt}{\ms{Set}}
\newcommand{\cnt}{\ms{Cnt}}
\newcommand{\pcat}{\ms{PCat}}
\newcommand{\cat}{\ms{Cat}}
\newcommand{\pcatt}{\ms{PCAT}}
\newcommand{\catt}{\ms{CAT}}
\newcommand{\PCat}{\mb{PCat}}
\newcommand{\Cat}{\mb{Cat}}
\newcommand{\sFrm}{\sigma\mb{Frm}}
\newcommand{\Frm}{\mb{Frm}}
\newcommand{\Loc}{\mb{Loc}}
\newcommand{\PCAT}{\mb{PCAT}}
\newcommand{\CAT}{\mb{CAT}}
\newcommand{\Catt}{\mf{Cat}}
\newcommand{\CATT}{\mf{CAT}}
\newcommand{\Topp}{\mb{Esp}}
\newcommand{\Tp}{\ms{TYPE}}
\newcommand{\Pp}{\ms{PROP}}
\newcommand{\St}{\ms{SET}}
\newcommand{\Gp}{\ms{GRPD}}
\newcommand{\fst}{\ms{Fin}}
\newcommand{\Mod}{\mb{Mod}}
\newcommand{\Spec}{\mb{Spec}}
\newcommand{\dneg}{\neg\neg}
\newcommand{\wTop}{\omega\mb{Esp}}
\newcommand{\wCPO}{\omega\mb{CPO}}
\newcommand{\Ob}{\mr{Ob}}
\newcommand{\subst}{\mb{subst}}
\newcommand{\PA}{\mbb{PA}}
\newcommand{\NPA}{\Nb\hspace{-.4ex}\ltimes\hspace{-.4ex}\PA}

\NewDocumentCommand\obsle{o}{\sqsubseteq_{\ms{obs}}\IfValueT{#1}{^{#1}}}
\NewDocumentCommand\satle{o}{\le_{\ms{sat}}\IfValueT{#1}{^{#1}}}
\NewDocumentCommand\behle{o}{\sqsubseteq_{\ms{beh}}\IfValueT{#1}{^{#1}}}
\NewDocumentCommand\opens{}{\mc{O}}


\DeclareFontFamily{U}{dmjhira}{}
\DeclareFontShape{U}{dmjhira}{m}{n}{ <-> dmjhira }{}
\DeclareRobustCommand{\yon}{\text{\usefont{U}{dmjhira}{m}{n}\symbol{"48}}}
\DeclareRobustCommand{\noy}{\text{\reflectbox{\yon}}\!}

\makeatletter
\newcommand{\ct@}[2]{%
  \vtop{\m@th\ialign{##\cr
    \hfil$#1\operator@font lim$\hfil\cr
    \noalign{\nointerlineskip\kern1.5\ex@}#2\cr
    \noalign{\nointerlineskip\kern-\ex@}\cr}}%
}
\newcommand{\ct}{%
  \mathop{\mathpalette\ct@{\rightarrowfill@\textstyle}}\nmlimits@
}
\makeatother
\makeatletter
\newcommand{\lt@}[2]{%
  \vtop{\m@th\ialign{##\cr
    \hfil$#1\operator@font lim$\hfil\cr
    \noalign{\nointerlineskip\kern1.5\ex@}#2\cr
    \noalign{\nointerlineskip\kern-\ex@}\cr}}%
}
\newcommand{\lt}{%
  \mathop{\mathpalette\lt@{\leftarrowfill@\textstyle}}\nmlimits@
}
\makeatother

\newtheorem{theorem}{Theorem}[section]
\newtheorem{lemma}[theorem]{Lemma}
\newtheorem{corollary}[theorem]{Corollary}
\newtheorem{proposition}[theorem]{Proposition}

\theoremstyle{definition}
\newtheorem{example}[theorem]{Example}
\newtheorem{definition}[theorem]{Definition}
\newtheorem{remark}[theorem]{Remark}
\newtheorem{notation}[theorem]{Notation}
\newtheorem{construction}[theorem]{Construction}
\newtheorem{observation}[theorem]{Observation}
\newtheorem*{axiom}{Axiom}

% \AtBeginEnvironment{definition}{%
%   \pushQED{\qed}%
%   \renewcommand{\qedsymbol}{$\vartriangleleft$}% Optional: change symbol
% }
% \AtEndEnvironment{definition}{\popQED}
% \AtBeginEnvironment{example}{%
%   \pushQED{\qed}%
%   \renewcommand{\qedsymbol}{$\vartriangleleft$}% Optional: change symbol
% }
% \AtEndEnvironment{example}{\popQED}
% \AtBeginEnvironment{remark}{%
%   \pushQED{\qed}%
%   \renewcommand{\qedsymbol}{$\vartriangleleft$}% Optional: change symbol
% }
% \AtEndEnvironment{remark}{\popQED}
% \AtBeginEnvironment{notation}{%
%   \pushQED{\qed}%
%   \renewcommand{\qedsymbol}{$\vartriangleleft$}% Optional: change symbol
% }
% \AtEndEnvironment{notation}{\popQED}
% \AtBeginEnvironment{observation}{%
%   \pushQED{\qed}%
%   \renewcommand{\qedsymbol}{$\vartriangleleft$}% Optional: change symbol
% }
% \AtEndEnvironment{observation}{\popQED}
% \AtBeginEnvironment{axiom}{%
%   \pushQED{\qed}%
%   \renewcommand{\qedsymbol}{$\vartriangleleft$}% Optional: change symbol
% }
% \AtEndEnvironment{axiom}{\popQED}

\zcLanguageSetup{english}{
  type = observation,
    Name-sg = Observation,
    name-sg = observation,
    Name-pl = Observations,
    name-pl = observations,
}

\AddToHook{env/lemma/begin}{\zcsetup{countertype={theorem=lemma}}}
\AddToHook{env/corollary/begin}{\zcsetup{countertype={theorem=corollary}}}
\AddToHook{env/proposition/begin}{\zcsetup{countertype={theorem=proposition}}}
\AddToHook{env/definition/begin}{\zcsetup{countertype={theorem=definition}}}
\AddToHook{env/example/begin}{\zcsetup{countertype={theorem=example}}}
\AddToHook{env/remark/begin}{\zcsetup{countertype={theorem=remark}}}
\AddToHook{env/notation/begin}{\zcsetup{countertype={theorem=notation}}}
\AddToHook{env/axiom/begin}{\zcsetup{countertype={theorem=axiom}}}
\AddToHook{env/construction/begin}{\zcsetup{countertype={theorem=construction}}}
\AddToHook{env/observation/begin}{\zcsetup{countertype={theorem=observation}}}

\allowdisplaybreaks

%%%%%%%%%%%%%%%%%%%%%%%%%%%%%%%%%%%%%%%%%%%%%%%%%%%%%%%%%%%%%%%%%%%%%%%%%%%%%%%%
%% Thesis meta-information
%%

%% The title of the thesis:
\title{The internal language of classifying topoi and applications in synthetic mathematics}

%% The full name of the author (e.g.: James Smith):
\author{Lingyuan Ye}

%% College affiliation:
\college{Lucy Cavendish College}

\collegeshield{CollegeShields/LucyCavendish}

%% Submission date [optional]:
\submissiondate{June, 2027}

%% You can redefine the submission notice [optional]:
% \submissionnotice{A badass thesis submitted on time for the Degree of PhD}

%% Declaration date:
\date{06, 2027}

%% PDF meta-info:
\subjectline{Computer Science}
\keywords{topoi, internal language, synthetic mathematics}

%%%%%%%%%%%%%%%%%%%%%%%%%%%%%%%%%%%%%%%%%%%%%%%%%%%%%%%%%%%%%%%%%%%%%%%%%%%%%%%%
%% Abstract:
%%
\abstract{%
  This thesis contributes to the understanding of the internal language of topoi, through the \emph{first-order theory it classifies}. The PhD thesis of~\citet{blechschmidt:2017} with the unpublished but wildly circulated note~\cite{blechschmidt_general_2020} has introduced a \emph{quasi-coherence (\gls{QC}) principle} for classifying topoi of Horn theories, connecting the higher-order internal logic of the classifying topos with the first-order Horn theory it classifies. We have generalised this principle to the case of \emph{$\kappa$-Horn theories}, which allows for operations with \emph{higher arity}, and also to the case of \emph{coherent theories}, which allows for theories with coherent axioms. 

  These generalisations greatly improve the applicability of the quasi-coherence principle, and the second part of this thesis showcases this in various contexts, drawing new connections between different \emph{first-order theories} and branches of \emph{synthetic mathematics}: 
  \begin{itemize}
    \item The classifying topoi for \emph{distributive lattices} and \emph{$\sigma$-frames}: the \gls{QC} for which provides a full account for the axioms of \emph{synthetic domain theory (\gls{SDT})} as specified by~\citet{hyland_first_2006}. These techniques also draw new connections between the properties of an \emph{interval object} and that of the theory of distributive lattices, with further applications in synthetic higher category theory.
    \item The classifying topos for \emph{coherent arithmetic} as in~\citet{ye_categorical_2026}: the \gls{QC} for which induces the structure of \emph{synthetic computability} similar to that specified in~\citet{bauer_first_2006}. In particular, it produces a \emph{sheaf topos} validating a form of \emph{Church's thesis}, and has a dominance that becomes a Lawvere's fixed-point object. This unveils a striking parallelism between the structure of arithmetic and computability.
    % \item The classifying topos for \emph{Heyting algebras}: the \gls{QC} for which provides a new environment for doing \emph{synthetic topology}. In particular, there is a striking new connection between \emph{uniform interpolation} for intuitionistic propositional logic established in~\citet{pitts_interpretation_1992}, and the topological structures of open subspaces.
    \item The classifying topos for \emph{decidable atomless Boolean algebras}: the \gls{QC} for which provides a striking connection with \emph{synthetic probability theory} in the style of~\citet{SimpsonAlex2024EaCI}, and the structure of \emph{partitions} of a Boolean algebra. The \emph{quantifier elimination} for this coherent theory also provides new proof strategies for establishing properties of independence relations of relative probability distribution in a synthetic way.
  \end{itemize}
}



%%%%%%%%%%%%%%%%%%%%%%%%%%%%%%%%%%%%%%%%%%%%%%%%%%%%%%%%%%%%%%%%%%%%%%%%%%%%%%%%
%% Acknowledgements:
%%
\acknowledgements{%

% First and foremost, I would like to thank my supervisor, Jon Sterling.

% At last, I want to thank \begin{CJK}{UTF8}{gbsn}越\end{CJK}. In the current time, both of us seeking a career in academia while at the same time being in love and staying as a couple seems like impossible. Maybe the love we share is both a gift and a curse in this respect. However, you stood by me no matter how difficult the situation was -- for this I'm eternally grateful. I will always love you.
}



%%%%%%%%%%%%%%%%%%%%%%%%%%%%%%%%%%%%%%%%%%%%%%%%%%%%%%%%%%%%%%%%%%%%%%%%%%%%%%%%
%% Glossary [optional]:
% %%
\newglossaryentry{HoTT}{
    name=HoTT,
    description={homotopy type theory}
}
\newglossaryentry{IPC}{
    name=IPC,
    description={intuitionistic propositional calculus}
}
\newglossaryentry{LEM}{
    name=LEM,
    description={the Law of Excluded Middle}
}
\newglossaryentry{PNNO}{
    name=PNNO,
    description={parametrised natural numbers object}
}
\newglossaryentry{SDT}{
    name=SDT,
    description={synthetic domain theory}
}
\newglossaryentry{SGDT}{
    name=SGDT,
    description={synthetic guarded domain theory}
}
\newglossaryentry{QC}{
    name=QC,
    description={quasi-coherence}
}
\newglossaryentry{ZFC}{
    name=ZFC,
    description={Zermelo-Fraenkel set theory with choice}
}


%%%%%%%%%%%%%%%%%%%%%%%%%%%%%%%%%%%%%%%%%%%%%%%%%%%%%%%%%%%%%%%%%%%%%%%%%%%%%%%%
%% Contents:
%%
\begin{document}

%%%%%%%%%%%%%%%%%%%%%%%%%%%%%%%%%%%%%%%%%%%%%%%%%%%%%%%%%%%%%%%%%%%%%%%%%%%%%%%%
%% Title page, abstract, declaration etc.:
%% -    the title page (is automatically omitted in the technical report mode).
\frontmatter{}

\setcounter{page}{1}

%%%%%%%%%%%%%%%%%%%%%%%%%%%%%%%%%%%%%%%%%%%%%%%%%%%%%%%%%%%%%%%%%%%%%%%%%%%%%%%%
%% Thesis body:
%%
\chapter{Introduction}


\section{Notation convention}

\section{Foundation choice}
% \chapter{Preliminaries}\label{ch:preliminary}

In this chapter we set up basic notation convention for this document, and collect various notions and constructions that will be useful later in the file.

\section{Notation and terminology convention}\label{sec:notation}

For any $n\in\N$, we use $\mb n$ to denote the finite cardinal containing exactly $n$ elements. In particular, $\mb 0$ and $\mb 1$ will be identified with the proposition \emph{falsum} and \emph{verum}, respectively.

For any set $X$, we write $\pss{X}$ for the proposition that $X$ is inhabited, viz.\ $\ex{x}{X}\top$.

We reserve symbols $\mc A,\mc B,\mc C$ for small categories. Fixed large categories are often denoted by bold face symbols, like $\Set$, $\Fin$, etc. We will reserve the sf font for 2-categories. For instance, $\Catt$ will denote the \emph{2-category} of categories, while $\Cat$ will be the \emph{1-category} of categories and functors. Other common 2-categories include $\Lex$, the 2-category of small lex categories, and $\Top$, the 2-category of topoi.

We use $\yon$ to denote the Yoneda embedding. More explicitly, we use the symbol for both the covariant and contravariant Yoneda embedding, but distinguishing them by subscription or superscription,
\[ \yon_- \colon \mc C \hook \psh(\mc C), \quad \yon^- \colon \mc C\op \hook [\mc C,\Set]. \]

In our constructive setting, by an \emph{equivalence} between categories we mean a functor $F \colon \mc A \to \mc B$ that is both \emph{fully faithful} and \emph{essentially surjective}. This is sometimes also referred to as a \emph{weak equivalence} between categories. For most practical usage, the notion of weak equivalence suffices, since one can show that all relevant categorical properties are preserved under weak equivalence. 


\begin{remark}[Comparison with homotopy type theory]
  This document has adopted the language of the old-fashioned Bishop-style constructive set theory, instead of currently more trending homotopy type theory (\gls{HoTT}), as described in the book~\cite{hottbook}. In particular, using the notion of \emph{univalent categories} developed in~\citet{AHRENS_KAPULKIN_SHULMAN_2015}, one can in fact show every weak equivalence between univalent categories admits an inverse. However, \ldots
\end{remark}

\section{Dominance}\label{sec:prelim-dominance}

We recall the notion of \emph{dominance} in the sense of~\citet{rosolini1986continuity}:

\begin{definition}[Dominance]\label{def:dominance}
  A dominance $\Sigma$ is a subuniverse of $\pp$ that is closed under dependent sums, i.e.\ $\mb 1\in\Sigma$, and for any $\varphi\in\Sigma$ together with a family $\psi \colon \varphi \to \Sigma$, the dependent sum $\sum_{x\in\varphi}\psi(x)$ also belongs to $\Sigma$.
\end{definition}

\section{Classical propositions}

Since we work in general in a constructive environment, propositions satisfying certain classical properties deserve additional attention:

\begin{definition}[Decidable and classical propositions]
  We say a proposition $p \in \pp$ is \emph{decidable}, if $p = 0 \vee p = 1$. We say $p$ is \emph{classical}, if $\neg\neg p = p$. 
\end{definition}

\begin{remark}[The subuniverse of decidable and classical propositions]
  Note that the set of decidable propositions is given by $\mb 2 = \set{\mb 0,\mb 1}$. We also denote the set of classical propositions as $\pp_{\dneg}$, and it is evident that there are inclusions $\mb 2 \subseteq \pp_{\dneg} \subseteq \pp$. These are both \emph{dominances} as in~\zcref{def:dominance}.
\end{remark}

\begin{definition}[Decidable and classical sets]\label{def:decidableset}
  For a set $X$, we say it is \emph{decidable} (resp.\ \emph{classical}) if its equalities are decidable (resp.\ classical), i.e.\ for any $x,y\in X$, the proposition $x = y$ is decidable (resp.\ classical).
\end{definition}

\begin{remark}[Caveat on complemented subtypes]\label{rem:complementedsubtype}
  Given a set $X$, it is common in the literature to refer to $S$ as a \emph{decidable subset} of $X$ if, as a function $S \colon X \to \pp$, it factors through the decidable propositions $2 \subseteq \pp$. However, this might be misunderstood as the set $S$ itself is decidable as in~\zcref{def:decidableset}. Hence, we will always refer to such a subset as a \emph{complemented} subset.
\end{remark}

\section{Stack semantics}\label{sec:stacksem}

Another important property for a set is whether it satisfies \emph{choice}. We introduce the following notion of \emph{projectivity} of a set:

\begin{definition}[Projectivity]
  A set $X$ is \emph{projective} if for any surjection $f \colon Y \surj X$ there exists a section.
\end{definition}

\begin{example}[Finite sets are projective]\label{exa:finiteprojective}
  For any $n\in\N$, the finite cardinal $\mb n$ is projective.
\end{example}

\begin{example}[Countable choice]\label{exa:countablechoice}
  The projectivity of $\N$ is simply \emph{choice principle} as discussed in~\zcref{sec:foundation}.
\end{example}

\begin{lemma}\label{lem:projdecsubset}
  If $X$ is projective, then any complemented subset $I$ of $X$ is also projective.
\end{lemma}
\begin{proof}
  This holds since if $I$ is complemented with complement $J$, then for any surjection $Y \surj I$ over $I$, one can construct a surjection $Y + J \surj I + J \cong X$, where now projectivity of $X$ induces a section. Restricting that section to $I$ provides a section of $Y \surj I$. 
\end{proof}

Recall the well-established notion of \emph{internal projectivity} in a topos; see e.g.~\citet[Exer. IV.16]{maclane1992sheaves}. An object $X\in\mc E$ is internally projective if and only if the functor $(-)^X \colon \mc E \to \mc E$ preserves epis.

\begin{proposition}[Projectivity in topoi]\label{prop:projectiveintopos}
  An object $X\in\mc E$ is projective in the stack semantics if and only if it is internally projective.
\end{proposition}
\begin{proof}
  For a proof see~\citet[Thm. 3.3]{nlab:internally_projective_object}.
\end{proof}

\begin{theorem}[Presheaf categories have countable choice]\label{thm:liftingofcountablechoice}
  For any small category $\mc C$, the natural numbers object $\N$ in $\psh(\mc C)$ is (internally) projective.
\end{theorem}
\begin{proof}
  Recall that the natural numbers object $\N$ in $\psh(\mc C)$ is the constant presheaf on $\N$, viz.\ $\N \cong \Delta\N \cong \coprod_{n\in\N}\mb 1$. For any $X\in\psh(\mc C)$ and $c\in\mc C$, the exponential is thus computed by 
  \begin{align*}
    &X^\N(c) \\
    &\quad \cong \mc E(\yon_c \times \N, X) \\
    &\quad \cong \mc E(\coprod_{n\in\N}\yon_c,X) \\
    &\quad \cong \prod_{n\in\N}\mc E(\yon_c,X) \\
    &\quad \cong X(c)^\N
  \end{align*}
  In particular, the exponential $(-)^\N$ is in fact computed \emph{pointwise}. Since colimits in $\psh(\mc C)$ are also computed pointwise, an epi in $\psh(\mc C)$ is simply a pointwise surjection. Then the claim holds if $(-)^\N$ preserves epis in $\Set$, which follows from countable choice (in $\Set$).
\end{proof}
\part{The theories of synthetic quasi-coherence}\label{prt:theory}
\chapter{A constructive recall of locally presentable categories}\label{ch:lp}

In this chapter we recall some basic theory of locally presentable categories following~\citet{adamek1994locally}. The reason we are redeveloping some of these materials is that in this document we have adopted a constructive meta-theory, and the treatment of locally presentable categories in \emph{loc.\ cit.}\ is not always evidently constructive. 

Our choice for a constructive meta-theory is \emph{not} philosophical. One technique that will be useful for our development of \gls{QC} for Horn theories with higher arity is to interpret an externally proven fact in the internal logic of a suitable topos. Such an internalisation process is valid only if the internal logic of the topos supports all the logical resources used in the original proof, and in general a topos only supports constructive logic. Thus, a constructive development of the theory of locally presentable categories, despite of interest in itself, is also quite useful for developing the subject of this document.

Though the theory of locally presentable categories is well-established, one of the innovative parts of our constructive treatment is to find a suitable notion of \emph{size}, that can replace the notion of \emph{regular cardinal} in classical set theory.

\section{A notion of size}\label{sec:size}

\begin{definition}[Size]\label{def:size}
  A \emph{notion of size} $\kappa$ is a collection of sets such that:
  \begin{itemize}
    \item $\kappa$ is \emph{replete}, i.e.\ it is closed under isomorphisms of sets;
    \item $\kappa$ contains all \emph{finite cardinals} $\mb n$ for $n\in\N$;
    \item $\kappa$ is closed under \emph{dependent sums}: for any $A \in \kappa$ and a family $B \colon A \to \kappa$, the disjoint union $\sum_{a\in A}B(a)$ belongs to $\kappa$;
    \item sets in $\kappa$ are \emph{projective};
  \end{itemize}
  We say $\kappa$ is \emph{separated}, if it is also closed under identities:
  \begin{itemize}
    \item for any $X\in\kappa$ and $x,y\in X$, the proposition $x = y$ also belongs to $\kappa$.
  \end{itemize}
  We say $\kappa$ is \emph{essentially small}, if it is classified by a small set:
  \begin{itemize}
    \item there exists a small set $K$ and a map $\El \colon K \to \kappa$, such that for any $\kappa$-small set $I$, there exists $k\in K$ such that $I \cong \El(k)$.
  \end{itemize}
  A set is \emph{$\kappa$-small} if it belongs to $\kappa$, and $\kappa\hp\Set$ denotes the full subcategory of $\kappa$-small sets.
\end{definition}

\begin{remark}[Comparison of the notion of size]\label{rem:comparethenotionofsize}
  Classically in \gls{ZFC} the above notion of an essentially small (separated) size corresponds exactly to an infinite \emph{regular cardinal}: under \gls{LEM}, being separated is automatic since every set is decidable; under choice, every set is also projective; thus the remaining condition of containing all finite sets and being closed under dependent sums exactly corresponds to an infinite cardinal. Constructively, a very similar notion of size was also independently proposed in a recent work~\citet{coquand_et_al:LIPIcs.LICS.2026.31}, though they do not require a size to be separated; this condition is to make sure $\kappa$-small sets induce $\kappa$-small discrete groupoids (cf.~\zcref{exa:discgroupoidkappasmall}).
\end{remark}

\begin{remark}[Closure properties of $\kappa$-small sets]\label{rem:closurepropofkappa}
  For a notion of size $\kappa$, being closed under dependent sums, together with the other requirements specified in~\zcref{def:size}, has many useful consequences for closure properties of the category $\kappa\hp\Set$ of $\kappa$-small sets. For instance, $\kappa\hp\Set$ will be closed under \emph{$\kappa$-small coproducts}, in particular it will be closed under \emph{cartesian products}. If $\kappa$ is \emph{separated}, then it is also closed under \emph{equalisers}, since given $f,g \colon X \to Y$, the equaliser is given by the subset 
  \[ \mr{Eq}(f,g) \cong \scomp{x\in X}{fx = gx} \cong \sum_{x\in X}fx = gx. \]
  Hence, $\kappa\hp\Set$ is closed under $\kappa$-small coproducts and finite products in $\Set$, and if $\kappa$ is separated then it is in fact closed under all finite limits. However, in general it is \emph{not} closed under coequalisers (cf.~\zcref{rem:coequalisercntandmarkov}).
\end{remark}

For our purpose, we are exclusively interested in the following two notions of size:

\begin{example}[Finiteness as a notion of size]\label{exa:finitesize}
  We say a set $X$ is \emph{finite} if it is isomorphic to a finite cardinal,
  \[ \ms{isFin}(X) \coloneq \ex n{\N} X \cong \mb n. \]
  Finite sets are \emph{decidable}, thus they are closed under identities. Finite sets are also closed under dependent sums, which can be shown by induction. By construction, finite sets are also projective. Furthermore, finite sets by definition has a classifier $\N$. Thus, finite sets form a separated and essentially small size in the sense of~\zcref{def:size}, and we denote this size as $\omega$. The full subcategory $\omega\hp\Set$ of finite sets will also be denoted as $\Fin$. 
\end{example}

\begin{example}[Countability as a notion of size]\label{exm:countablesize}
  A set $X$ is \emph{countable} if it is isomorphic to a complemented subset of the natural numbers,
  \[ \ms{isCnt}(X) \coloneq \ex{I}{\mb 2^\N} X \cong \scomp{n\in\N}{I(n)}. \]
  The collection of countable sets will be denoted as $\omega_1$. By construction, all countable sets are (isomorphic to) complemented subsets of $\N$, thus they are all decidable. This in particular means $\omega_1$ is closed under identities. By countable choice, $\N$ is projective, and thus by~\zcref{lem:projdecsubset} all countable sets are projective. Suppose $X$ is countable and $Y \colon X \to \omega_1$ is a countable family of countable sets. To show $\sum_{x:X}Y_x$ is countable, we may assume $X$ is isomorphic to $I$ for some complemented subset $I$ of $\N$. Then by assumption, 
  \[ \fa{i}{I}\ex{J}{\mb 2^\N} Y_i \cong \scomp{n\in\N}{J(n)}. \] 
  By $I$ being projective, there exists a family $J_i$ of complemented subsets of $\N$ such that $Y_i \cong J_i$. Thus, the dependent sum $\sum_{x:X}Y_x$ will be equivalent to $\sum_{i:I}J_i$, which is a complemented subset of $\N \times \N$. Via a bijection $\N \times \N \cong \N$, this shows $\sum_{x:X}Y_x$ is countable. Furthermore, again by construction $\omega_1$ has a classifier $\mb 2^\N$. This shows $\omega_1$ is also a separated essentially small size. We denote the category $\omega_1\hp\Set$ of countable sets also as $\Cnt$.
\end{example}

\begin{remark}[Coequalisers of countable sets]\label{rem:coequalisercntandmarkov}
  $\Fin$ can be shown to be a Boolean elementary topos, see e.g.~\citet[D5.2]{johnstone2002sketches}. In particular, it has coequalisers. On the other hand, constructively $\Cnt$ is not necessarily closed under coequalisers. For instance, it is not hard to see that $\Cnt$ having coequalisers implies the \emph{Markov principle}.
\end{remark}

At the end of this section, we observe that any notion of size $\kappa$ induces a \emph{dominance} in the sense of~\citet{rosolini1986continuity}:

\begin{definition}[Dominance]\label{def:dominance}
  A dominance $\Sigma$ is a subuniverse of $\pp$ that is closed under dependent sums, i.e.\ $\mb 1\in\Sigma$, and for any $\varphi\in\Sigma$ together with a family $\psi \colon \varphi \to \Sigma$, the dependent sum $\sum_{x\in\varphi}\psi(x)$ also belongs to $\Sigma$.
\end{definition}

\begin{definition}[$\kappa$-overt propositions]\label{def:kappaopen}
  A proposition is \emph{$\kappa$-overt}, if it is of the form $\pss X$ for some $\kappa$-small set $X$. The subuniverse of $\kappa$-overt propositions will be denoted as $\mb{O}_\kappa$. 
\end{definition}

\begin{proposition}\label{prop:sizedominance}
  For any notion of size $\kappa$, $\mb{O}_\kappa$ will be a dominance closed under finite joins.
\end{proposition}
\begin{proof}
  Note that $\mb 1\in\mb{O}_\kappa$. It is closed under finite joins because $\kappa$-small types are closed under finite coproduct (cf.~\zcref{rem:closurepropofkappa}). Now suppose we have $\varphi \in \mb{O}_\kappa$, and consider a family $\psi \colon \varphi \to \mb{O}_\kappa$ of $\kappa$-overt propositions indexed by $\varphi$. By definition, this means we have
  \[ \fa{i}{\varphi}\ex{Y}{\kappa}\psi_i = \pss{Y}. \]
  To show the dependent sum $\sum_{i\in\varphi}\psi_i$ belongs to $\mb{O}_\kappa$, by assumption we may assume $\varphi$ is of the form $\pss X$ for some $X \in \kappa$. In particular, there is a surjection $q \colon X \surj \varphi$. Hence,
  \[ \fa{x}{X}\ex Y{\kappa} \psi_{qx} = \pss{Y}. \]
  Since $X\in\kappa$, $X$ is projective, thus there exists an $X$-indexed family of $\kappa$-small sets $Y \colon X \to \kappa$, where $\psi_{qx} = \pss{Y_x}$ for all $x\in X$. This way, the dependent sum can be computed by
  \[ \sum_{i\in\varphi}\psi_i = \pss{\sum_{x\in X}Y_x}. \]
  Hence, $\mb{O}_\kappa$ is closed under dependent sums, and forms a dominance.
\end{proof}

\begin{remark}[Relation to the dominance of semi-decidable propositions]
  It is shown by~\citet[Prop. 5.2.1]{rosolini1986continuity} (with the proof is attributed to Dana Scott) that the following subuniverse of \emph{semi-decidable propositions} is a dominance,
  \[ \Sigma_{\mr{s.d.}} \coloneq \scomp{p\in\pp}{\ex f{\N^\N} p = \ex{n}{\N}fn = 0}. \]
  In fact, it is not hard to see that the universe $\mb O_{\omega_1}$ of $\omega_1$-overt propositions is identical to $\Sigma_{\mr{s.d.}}$. The proof in \emph{loc.\ cit.}\ also crucially relies on countable choice. In fact, the proof of~\zcref{prop:sizedominance} can be seen as a suitable generalisation of this proof that works for any notion of size. For $\omega$ we have $\mb O_\omega = \mb 2$, recovering the dominance of decidable propositions.
\end{remark}

\section{Filtered diagrams}

In this section we fix a \emph{separated size} $\kappa$. We say a category $\mc A$ is \emph{$\kappa$-small} if its set of objects is $\kappa$-small, and for any $a,b \in \mc A$, the hom-set $\mc A(a,b)$ is $\kappa$-small. 

\begin{example}[Discrete groupoids from $\kappa$-smalls sets are $\kappa$-small]\label{exa:discgroupoidkappasmall}
  The requirement for $\kappa$ to be separated is to make $\kappa$-small sets also $\kappa$-small discrete groupoids. For any $X\in\kappa$, the hom-sets in the associated discrete groupoid is given by $X(x,y) \cong (x = y)$ for any $x,y\in X$. By assumption the identity proposition $x = y$ is $\kappa$-small since $\kappa$ is separated, hence $X$ as a discrete groupoid will be a $\kappa$-small category.
\end{example}

\begin{definition}[$\kappa$-(co)completeness]
  A category $\mc A$ is $\kappa$-(co)complete if it has $\kappa$-small (co)limits, i.e.\ if for any $\kappa$-small category $\mc D$ and diagram $F \colon \mc D \to \mc A$, the (co)limit of $F$ exists.
\end{definition}

By~\zcref{exa:discgroupoidkappasmall}, diagrams indexed by $\kappa$-small sets will in particular be instances of $\kappa$-small (co)limits, thus $\kappa$-(co)complete categories have $\kappa$-small (co)products. Since any $\kappa$ also contains all finite sets, $\kappa$-(co)complete categories also have (co)equalisers. In fact, the existence of these two classes of (co)limits classifies $\kappa$-(co)completeness:

\begin{proposition}\label{prop:smallcompiffeqcoandeq}
  $\mc A$ is $\kappa$-(co)complete if it has (co)equalisers and $\kappa$-small (co)products.
\end{proposition}
\begin{proof}
  The only if direction is trivial. For the if direction, for a $\kappa$-small diagram $F \colon \mc D \to \mc A$, its limit is given by the following equaliser,
  \[ 
  \begin{tikzcd}
    \lt F \ar[r, dashed, tail] & \prod_{d\in\Ob(\mc D)}Fd \ar[r, shift left] \ar[r, shift right] & \prod_{\alpha\in{\mc D}(d_0,d_1)}Fd_1
  \end{tikzcd}
  \]
  It then suffices to observe that $\prod_{d\in\Ob(\mc D)}Fd$ and $\prod_{\alpha\in\mc D(d_0,d_1)}Fd_1$ are $\kappa$-small products: the former by assumption that $\Ob(\mc D)$ is $\kappa$-small, and the latter by the fact that the dependent sum $\sum_{d_0,d_1\in\Ob(\mc D)}\mc D(d_0,d_1)$ is $\kappa$-small, since by~\zcref{rem:closurepropofkappa} $\kappa$-small sets are closed under finite products and dependent sums.
\end{proof}

\begin{definition}[$\kappa$-filtered category]
  We say a category $\mc I$ is \emph{$\kappa$-filtered}, if it is small and every $\kappa$-small diagram $\mc D \to \mc I$ admits a cocone.
\end{definition}

\begin{remark}[$\kappa$-cocomplete categories are $\kappa$-filtered]\label{rem:kappacompeltefiltered}
  If $\mc I$ has all $\kappa$-small colimits, then in particular every $\kappa$-small diagram has a universal cocone, hence it is $\kappa$-filtered.
\end{remark}

As usual, $\kappa$-filtered categories can be characterised in a more concrete fashion:

\begin{lemma}\label{lem:concretefilter}
  A category $\mc I$ is $\kappa$-filtered if and only if the following hold:
  \begin{itemize}
    \item For any $\kappa$-small set $I$ and any diagram $f \colon I \to \mc I$, there exists $x\in\mc I$ and a family of morphisms $\alpha \in \prod_{i\in I}\mc I(fi,x)$;
    \item For any $\kappa$-small set $I$, $x,y\in\mc I$, and any $I$-indexed family of arrows $\alpha \colon I \to \mc J(x,y)$, there exists $z\in\mc J$ and $\gamma \colon y \to z$ that coequalises them, i.e.\ for all $i,j\in I$, $\gamma \circ \alpha_i = \gamma \circ \alpha_j$.
  \end{itemize}
\end{lemma}
\begin{proof}
  Necessity is evident. For sufficiency, we essentially follow the proof strategy for the finitary case; see e.g.~\citet[Lem. 2.13.2]{borceux1994handbookI}. Let $F \colon \mc D \to \mc I$ be a $\kappa$-small diagram in $\mc I$. We construct a cocone as follows: 
  \begin{itemize}
    \item First pick $k_0 \in \mc I$ with a family $\alpha \in \prod_{d\in\mc D}{\mc I}(Fd,k_0)$;
    \item For any $f \in {\mc D}(d,e)$, pick $k_f$ with $\beta_f \in {\mc I}(k_0,k_f)$, coequalising the pair $\alpha_d \circ Ff$ and $\alpha_e$;
    \item We now again have a family $k_f$ indexed by $\sum_{d,e\in\mc D}{\mc D}(d,e)$, thus we may find some $k_1$ with a family of maps $\gamma_f \in {\mc I}(k_f,k_1)$.
    \item Thus, we get a $\kappa$-small family of maps $\gamma_-\circ\beta_- \colon k_0 \to k_1$ in $\mc I$, indexed by $\sum_{d,e\in\mc D}{\mc D}(d,e)$. By assumption, we can find $k$ with $\delta \in {\mc I}(k_1,k)$ coequalising them all.
  \end{itemize}
  We then define the cocone to be consisting of the vertex $k$, with component for each $d\in\mc D$
  \[ \epsilon_d \coloneq \delta \circ \gamma_{1_d} \circ \beta_{1_d} \circ \alpha_d \in {\mc I}(Fd,k). \]
  This is a cocone by construction. 
\end{proof}

In $\Set$, colimits indexed by $\kappa$-filtered categories are easier to compute, since one can explicitly characterise the equivalence relation generating the colimit:

\begin{lemma}\label{lem:filteredquotient}
  Let $F \colon \mc I \to \Set$ be a $\kappa$-filtered diagram. Then the following relation $\sim$ defined on $\sum_{i\in\mc I}Fi$ is an equivalence relation,
  \[ (i,x) \sim (j,y) \coloneq \ex{k}{\mc I}\ex{\alpha}{\mc I(i,k)}\ex{\beta}{\mc I(j,k)} F\alpha(x) = F\beta(y). \]
\end{lemma}
\begin{proof}
  Reflexivity and symmetry hold by construction. Transitivity holds by the fact that any $\kappa$-filtered category admits finite cocones.
\end{proof}

This leads to the following well-known result:

\begin{theorem}\label{thm:filtercommute}
  $\kappa$-filtered colimits commutes with $\kappa$-small limits in $\Set$: for any diagram $F \colon \mc I \times \mc D \to \Set$ where $\mc I$ is $\kappa$-filtered and $\mc D$ is $\kappa$-small, the canonical map below is an isomorphism,
  \[ \ct_{i\in\mc I}\lt_{d\in\mc D} F(i,d) \to \lt_{d\in\mc D}\ct_{i\in\mc I}F(i,d). \]
\end{theorem}
\begin{proof}
  The map sends the equivalence class $[(i,x)]$ to the family $\ld{d}{\mc D}[(i,x_d)]$. To show it is injective, we may assume we have $[(i,x)]$ and $[(j,y)]$ agreeing in the target,
  \[ \fa{d}{\mc D}[(i,x_d)] = [(j,y_d)]. \]
  By Lemma~\ref{lem:filteredquotient}, this explicitly means that
  \[ \fa{d}{\mc D}\ex{k}{\mc I}\ex{\alpha}{\mc I(i,k)}\ex{\beta}{\mc I(j,k)} F(d,\alpha)(x_d) = F(d,\beta)(y_d). \]
  Since $\mc D$ is $\kappa$-small, its set of objects is projective. Hence we can assume that there exists a function $k \colon \Ob(\mc D) \to \Ob(\mc I)$ and similarly $\alpha,\beta$ picking out the morphisms satisfying the above condition. The resulting diagram on $\mc I$ will be $\kappa$-small, thus by assumption we can pick a cocone of it, which we denote as $(u,\gamma)$. In particular, it satisfies that
  \[ \fa{d,d'}{\mc D} \gamma_d \circ \alpha_d = \gamma_{d'} \circ \alpha_{d'} \conjt \gamma_d \circ \beta_d = \gamma_{d'} \circ \beta_{d'}. \]
  Now we have $u$ with $\gamma\circ\alpha \in \mc I(i,u)$ and $\gamma \circ \beta \in \mc I(j,u)$. To show $F(\gamma\circ\alpha)(x) = F(\gamma\circ\beta)(y)$, it suffices to show
  \[  \fa{d}{\mc D}F(d,\gamma_d\circ\alpha_d)(x_d) = F(d,\gamma_d\circ\beta_d)(y_d), \]
  which holds by construction that $(u,\gamma)$ is a cocone. Hence, $[(i,x)] = [(j,y)]$.

  For surjectivity, suppose we have a term in $\prod_{d\in\mc D}\ct_{iin\mc I}F(i,d)$. Again by $\Ob(\mc D)$ being projective, we may assume we have chosen a representative $[(i_d,x_d)]$ for all $d \in \mc D$, satisfying for any $\alpha \in \mc D(d,e)$,
  \[ [(i_d,F(i_d,\alpha)(x_d))] = [(i_e,x_e)]. \]
  In particular, this means there is $j_{d,e,\alpha} \in \mc I$, $\beta_{d,e,\alpha} \in \mc I(i_d,j_{d,e,\alpha})$, and $\gamma_{d,e,\alpha} \in \mc I(i_e,j_{d,e,\alpha})$, that
  \[ F(\beta_{d,e,\alpha},\alpha)(x_d) = F(\gamma_{d,e,\alpha},e)(x_e). \]
  Now the diagram indexed by $\sum_{d,e\in\mc D}\mc D(d,e)$ is again $\kappa$-small, hence we can find a cocone $k$ with $\delta,\xi$, such that for any $d,e \in \mc D_0$ and $\alpha \in \mc D(d,e)$,
  \[ \delta_{d,e,\alpha} \circ \beta_{d,e,\alpha} = \xi_d, \quad \delta_{d,e,\beta} \circ \gamma_{d,e,\alpha} = \xi_e. \]
  This way, we can pick the representative $(k,x)$ with
  \[ x \coloneq \ld{d}{\mc D} F(\xi_d,d)(x_d). \]
  This is indeed well-defined, because for any $\alpha \in \mc D(d,e)$,
  \[ F(\xi_d,\alpha)(x_d) = F(\delta_{d,e,\alpha}\circ\beta_{d,e,\alpha},\alpha)(x_d) = F(\delta_{d,e,\alpha}\circ\gamma_{d,e,\alpha},e)(x_e) = F(\xi_e,e)(x_e). \]
  Thus, we only need to observe that $[(k,d)]$ is mapped to $\ld d{\mc D}[(i_d,x_d)]$, which holds by construction.
\end{proof}



\section{Locally presentable categories}

Again we fix a separated notion of size $\kappa$.

\begin{definition}[Presentable objects]
  Let $\mc A$ be a category having $\kappa$-filtered colimits. An object $a \in \mc A$ is \emph{$\kappa$-presentable}, if the representable functor $\mc A(a,-) \colon \mc A \to \Set$ preserves $\kappa$-filtered colimits. The full subcategory of $\kappa$-presentable objects will be denoted as $\mc A_\kappa$.
\end{definition}

\begin{remark}[Finitely and countably presentable sets]
  We will see later in~\zcref{exa:presentableobjectsinset} that $\Set_{\omega} = \Fin$. On the other hand, constructively $\Set_{\omega_1}$ may not coincide with $\Cnt$, since the former will be closed under countable colimits (as we will soon see in~\zcref{lem:presentableclosedundercolimits}), in particular coequalisers, while as we have seen in~\zcref{rem:coequalisercntandmarkov} the latter may not be.
\end{remark}

\begin{lemma}\label{lem:presentableclosedundercolimits}
  $\mc A_{\kappa}$ is closed under $\kappa$-small colimits that exists in $\mc A$.
\end{lemma}
\begin{proof}
  This is a formal consequence of Theorem~\ref{thm:filtercommute}.
\end{proof}

To define the notion of locally $\kappa$-presentable categories, we first consider the following notion of \emph{generators} (cf.~\citet[Def. 4.5.4]{borceux1994handbookI}):

\begin{definition}[Dense subcategory]
  For any category $\mc A$, a full subcategory $\mc B \subseteq \mc A$ is \emph{dense} if for any $a \in \mc A$, the canonical map below is an isomorphism,
  \[ \ct_{b\in\mc B/a}b \to a. \]
\end{definition}

There is an equivalent criterion for dense subcategories:

\begin{lemma}\label{lem:densefullyfaithful}
  $\mc B \subseteq \mc A$ is dense iff the restricted Yoneda functor below is fully faithful,
  \[ R \colon \mc A \to \psh(\mc B). \]
\end{lemma}
\begin{proof}
  See~\citet[Prop. 4.5.14]{borceux1994handbookI}; the proof there is constructive.
\end{proof}

\begin{definition}[Locally presentable categories]\label{def:locallypresentablecategory}
  A category $\mc A$ is \emph{locally $\kappa$-presentable}, if it is cocomplete and $\mc A_\kappa$ is dense and essentially small.
\end{definition}

\begin{remark}[Locally presentable categories are locally small]
  If $\mc A$ is locally $\kappa$-presentable then it must be locally small, because it is a full subcategory of $\psh(\mc A_\kappa)$ for the essentially small category $\mc A_\kappa$.
\end{remark}

\begin{example}[Presheaf categories are locally $\omega$-presentable]\label{exm:pshupresentable}
  For any small category $\mc B$, the presheaf category $\psh(\mc B)$ is locally $\omega$-presentable.
\end{example}

In general, we have the following fact:

\begin{proposition}\label{prop:locallypresaccembedpsh}
  Let $\mc A$ be locally $\kappa$-presentable. The restricted Yoneda embedding
  \[ R \colon \mc A \hook \psh(\mc A_\kappa), \]
  preserves $\kappa$-filtered colimits, and it has a left adjoint $L$.
\end{proposition}
\begin{proof}
  For any $\kappa$-filtered colimit $a \cong \ct_{i\in\mc I}a_i$ in $\mc A$ and $b \in \mc A_\kappa$,
  \[ Ra(b) \cong \mc A(b,a) \cong \ct_{i\in\mc I}\mc A(b,a_i) \cong \ct_{i\in\mc I}Ra_i(b). \]
  This uses the fact that $b$ is $\kappa$-presentable, hence it follows that $R$ preserves $\kappa$-filtered colimits. It is also fully faithful because $\mc A_\kappa$ is dense. The left adjoint is simply given by point-wise left Kan extension, which takes a presheaf $F$ to the following colimit,
  \[ LF \coloneq \ct_{b\in\mc A_\kappa/F}b. \]
  It is indeed a left adjoint because we have
  \begin{align*}
    &\mc A(LF,a) \\
    &\quad\cong \lt_{b\in\mc A_\kappa/F}\mc A(b,a) \\
    &\quad\cong \lt_{b\in\mc A_\kappa/F}\ct_{c\in\mc A_\kappa/a}\mc A(b,c) \\
    &\quad\cong \lt_{b\in\mc A_\kappa/F}\ct_{c\in\mc A_\kappa/a}\psh(\mc A_\kappa)(\yon_b,\yon_c) \\
    &\quad\cong \psh(\mc A_\kappa)(\ct_{b\in\mc A_\kappa/F}\yon_b,\ct_{c\in\mc A_\kappa/a}\yon_c) \\
    &\quad\cong \psh(\mc A_\kappa)(F,Ra)
  \end{align*}
  The first isomorphism holds by definition of $F$; the second holds since $b$ is $\kappa$-presentable, and that $\mc A_\kappa/a$ is $\kappa$-filtered by~\zcref{lem:presentableclosedundercolimits} and~\zcref{rem:kappacompeltefiltered}; the third holds by Yoneda; the forth holds since representables preserve all colimits; the last holds by definition and the fact that $R$ preserves $\kappa$-filtered colimits.
\end{proof}

\begin{corollary}
  If $\mc A$ is locally $\kappa$-presentable, then it is complete.
\end{corollary}
\begin{proof}
  By Proposition~\ref{prop:locallypresaccembedpsh}, $\mc A$ is a reflective subcategory of a presheaf category.
\end{proof}

Let $\kappa\hp\Lex$ denote the 2-category of categories with $\kappa$-small limits and functors preserving $\kappa$-small limits. We can characterise the essential image of the embedding $\mc A \hook \psh(\mc A_\kappa)$ for a locally $\kappa$-presentable category $\mc A$ as follows:

\begin{proposition}\label{prop:presentableindlexfunctor}
  For any locally $\kappa$-presentable category $\mc A$, 
  \[ \mc A \simeq \kappa\hp\Lex(\mc A_{\kappa}\op,\Set). \]
\end{proposition}
\begin{proof}
  For any $a\in\mc A$, by construction $Ra \cong \mc A(-,a)$ takes any colimit in $\mc A_\kappa$ to a limit in $\Set$, thus $R$ factors through $\kappa\hp\Lex(\mc A_\kappa\op,\Set)$. On the other hand, if $F \colon \mc A_{\kappa}\op \to \Set$ preserves $\kappa$-small limits, then its category of elements $\mc A_\kappa/F$ will be $\kappa$-filtered; a simple proof can be given through the characterisation of $\kappa$-filtered diagrams in~\zcref{lem:concretefilter}. This way,
  \[ F \cong \ct_{a\in\mc A_\kappa/F}\yon_a \cong R\ct_{a\in\mc A_\kappa/F}a, \]
  where the second isomorphism holds by the fact that $R$ preserves $\kappa$-filtered colimits by~\zcref{prop:locallypresaccembedpsh}. This shows the desired equivalence holds.
\end{proof}

% \begin{proposition}\label{prop:upresencloureflction}
%   If $\mc A$ is locally $\kappa$-presentable, and $\mc B$ is a full reflective subcategory of $\mc A$ closed under $\kappa$-filtered colimits, then $\mc B$ is also locally $\kappa$-presentable.
% \end{proposition}
% \begin{proof}
%   Let us use $R \colon \mc B \to \mc A$ to denote the inclusion and $L$ the reflection. Note that $\mc B$ is cocomplete, and for any $F\colon\mc I \to \mc B$, we have
%   \[ \ct_{u\in\mc I}Fu = L\ct_{u\in\mc I}RFu. \]
%   This means that if $i \in \mc A_{\kappa}$ then $Li \in \mc B_{\kappa}$, because for any $\kappa$-filtered diagram $F \colon \mc I \to \mc B$,
%   \begin{align*}
%     &\mc B(Li,\ct_{u\in\mc I}Fu) \\
%     &\quad\cong \mc A(i,R\ct_{u\in\mc I}Fu) \\
%     &\quad\cong \mc A(i,\ct_{u\in\mc I}RFu) \\
%     &\quad\cong \ct_{u\in\mc I}\mc A(i,RFu) \\
%     &\quad\cong \ct_{u\in\mc I}\mc B(Li,Fu).
%   \end{align*}
%   The second step uses the fact that $R$ preserves $\kappa$-filtered colimits; the other steps are routine. Hence, $\mc B_{\kappa}$ is the image of $\mc A_{\kappa}$ under $L$. 

%   Since we are proving a proposition, we can assume there is a small dense full subcategory $\mc U$ of $\kappa$-presentable objects in $\mc A$. Let $\mc V$ be the full subcategory of $\mc B$ under the image of $L$ of $\mc U$. Now for any $b : \mc B_0$, we have
%   \[ \ct_{i:\mc V/b}i = L\ct_{i:\mc V/b}Ri. \]
%   Note that there is a canonical map
%   \[ Rb = \ct_{j:\mc U/Rb}j \to \ct_{i:\mc V/b}Ri \]
%   induced by the functor $L : \mc U/Rb \to \mc V/b$ and the unit of the adjunction $L \dashv R$. This map is indeed localised by $L$, because the unit $j \to RLj$ are localised by $L$. This way, we have
%   \[ b = LRb = L\ct_{j:\mc U/Rb}j = L\ct_{i:\mc V/b}Ri = \ct_{i:\mc V/b}LRi = \ct_{i:\mc V/b}i. \]
%   Hence, $\mc B$ is locally $\kappa$-presentable.
% \end{proof}

% \begin{theorem}
%   The following are equivalent for a category $\mc A$:
%   \begin{enumerate}
%     \item $\mc A$ is locally $\kappa$-presentable;
%     \item $\mc A$ is a full reflective subcategory of $\psh(\mc B)$ for a small category $\mc B$, which are closed under $\kappa$-filtered colimits.
%   \end{enumerate}
% \end{theorem}
% \begin{proof}
%   (1) $\nt$ (2) is Proposition~\ref{prop:locallypresaccembedpsh}. For (2) $\nt$ (1), it follows by Example~\ref{exm:pshupresentable} and Proposition~\ref{prop:upresencloureflction}.
% \end{proof}
\chapter{The internal logic of classifying topoi for Horn theories with higher arity}

\section{Horn theories with higher arity}\label{sec:higherhorntheory}

In this section we recall some basic facts about infinitary Horn theories from~\citet[\S3.C]{adamek1994locally} and from~\citet{ADAMEK1998379}. Through this section, fix some regular cardinal $\kappa$.

\begin{remark}[Constructive notion of regular cardinal]
  For the constructively minded reader, we note that there is nothing special about the classical notion of a regular cardinal that we use. To interpret such a $\kappa$ in a constructive context, it suffices to assume it is a universe of sets closed under \emph{finite sets} and \emph{dependent sums}.
\end{remark}

\begin{definition}[Sorted signatures]
  Let $S$ be a set. An \emph{$S$-sorted signature} is a set $\Sigma$ with an arity function $a$ assigning each $\sigma\in\Sigma$ a collection of sorts $(s_i)_{i<\lambda}$ with some cardinal $\lambda$, and a sort $s$. In notation, we write 
  \[ \sigma \colon \prod_{i<\lambda}s_i \to s. \]
  $\Sigma$ is \emph{$\kappa$-ary} if the associated cardinality $\lambda$ for each $\sigma\in\Sigma$ is strictly smaller than $\kappa$.
\end{definition}

\begin{definition}[Structures for a signature]
  Let $\Sigma$ be an $S$-sorted signature. A \emph{$\Sigma$-structure} $A$ is a collection of sets $(A_s)_{s\in S}$ indexed by sorts in $S$, together with an operation $\ass\sigma_A$ for every $\sigma \colon \prod_{i<\lambda}s_i \to s$ in $\Sigma$, of type
  \[ \ass\sigma_A \colon \prod_{i<\lambda}A_{s_i} \to A_s. \]
  A \emph{homomorphism} of $\Sigma$-structures $f \colon A \to B$ is an $S$-indexed family of functions $f_s \colon A_s \to B_s$, making the following diagram commute for any $\sigma \colon \prod_{i<\lambda}s_i \to s$ in $\Sigma$,
  \[
  \begin{tikzcd}
    \prod_{i<\lambda}A_{s_i} \ar[d, "\prod_{i<\lambda}f_{s_i}"'] \ar[r, "\ass\sigma_A"] & A_s \ar[d, "f_s"] \\
    \prod_{i<\lambda}B_{s_i} \ar[r, "\ass\sigma_B"'] & B_s
  \end{tikzcd}
  \]
  The category of $\Sigma$-structures will be denoted as $\mmod{\Sigma}$.
\end{definition}

\begin{definition}[Terms of a signature]
  Let $\Sigma$ be an $S$-sorted signature. For any $S$-indexed family of sets $X = (X_s)_{s\in S}$, the $S$-indexed set of \emph{$\Sigma$-terms with variables in $X$} is defined as usual: it is the smallest $S$-indexed set such that (1) every $x\in X_s$ is a term of sort $s$; and (2) for any $\sigma \colon \prod_{i<\lambda}s_i \to s$ in $\Sigma$ and any terms $t_i$ of sort $s_i$, $\sigma((t_i)_{i<\lambda})$ is a term of sort $s$. This set will be denoted as $\tm{\Sigma}(X) = (\tm{\Sigma}(X)_{s})_{s\in S}$.
\end{definition}

\begin{remark}[Interpretation of terms in a structure]
  Let $A$ be a $\Sigma$-structure, and let $t$ be a $\Sigma$-term with free variables $\ov x$ of some sorts, given a list of elements $\ov a$ with the corresponding sorts, it is straight forward to define the interpretation of $t(\ov a)$ in $A$, written as
  \[ \ass{t(\ov a)}_A \coloneq \ass{t}_A(\ov a). \]
\end{remark}

\begin{definition}[$\kappa$-Horn and $\kappa$-algebra theory]
  Let $\Sigma$ be an $S$-sorted signature of arity $\kappa$. A \emph{$\kappa$-Horn clause} in signature $\Sigma$ is a sequent of the form 
  \[ \bigwedge_{i\in I}t_i = t_i' \vdash_{\ov x} t = t', \]
  where $t_i,t_i'$ and $t,t'$ are $\Sigma$-terms of the same sorts respectively, whose variables are among $\ov x$. We require $\geo{\ov x},\geo{I} < \kappa$. A \emph{$\kappa$-Horn theory} is simply a family of $\kappa$-Horn clauses. If all the Horn clauses in a theory have trivial antecedents, then we say such a theory is \emph{$\kappa$-algebraic}.
\end{definition}

\begin{definition}[Algebras for a $\kappa$-Horn theory]
  Let $\T$ be a $\kappa$-Horn theory over an $S$-sorted signature $\Sigma$. A $\Sigma$-structure $A$ is a \emph{$\T$-algebra}, if for any $\kappa$-Horn clause $\bigwedge_{i\in I}t_i = t_i' \vdash_{\ov x} t = t'$ in $\T$, and any lists of elements $\ov a$ with the corresponding sorts of $\ov x$,
  \[ \prth{\fa iI \ass{t_i(\ov a)} = \ass{t_i'(\ov a)}} \nt \ass{t(\ov a)} = \ass{t'(\ov a)}. \]
  The full subcategory of $\T$-algebras in $\mmod{\Sigma}$ will be denoted as $\mmod{\T}$.
\end{definition}

\begin{remark}[Variety and quasi-variety]
  The category of algebras for a ($\kappa$-ary) algebraic theory or Horn theory is often called a ($\kappa$-ary) \emph{variety} or \emph{quasi-variety} in the literature on universal algebra.
\end{remark}

\begin{example}[Finitary algebraic and Horn theories]\label{exa:finitealgebratheories}
  All the usual finitary algebraic and Horn theories will be $\aleph_0$-Horn theories. These include the theory of objects $\mbb O$, of meet-semi-lattices $\mbb M$, of distributive lattices $\mbb D$, of de Morgan algebras $d\mbb{M}$, of Heyting algebras $\mbb H$, of Boolean algebras $\mbb B$, etc. They are all 1-sorted algebra theories.
\end{example}

\begin{example}[The $\aleph_1$-Horn theory of posets with countable joins]\label{exa:theoryofcountablejoins}
  The theory $\omega\mbb J$ of posets with countable joins is described in~\citet[Exm. 3.26]{adamek1994locally}. The theory $\omega\mbb J$ is 1-sorted, with a constant 0 and an $\omega$-ary operator $\bigvee_{i<\omega}$, together with the following axioms:
  \begin{itemize}
    \item $\bigvee_{i<\omega}(x,x,\cdots) = x$;
    \item $\bigvee_{i<\omega}x_i = \bigvee_{j<\omega}y_j$ whenever $\set{x_i}_{i<\omega}\backslash\set{0} = \set{y_j}_{j<\omega}\backslash\set{0}$;
  \end{itemize}
  One may also define a binary join operator from this
  \[ x \vee y \coloneq \bigvee_{i<\omega}(x,y,y,\cdots), \]
  and under the usual definition of the order $a \le b \coloneq a \vee b = b$, each $\omega\mbb J$-algebra will be a poset with countable joins defined by the interpretation of $\bigvee_{i<\omega}$.
\end{example}

\begin{example}[The $\aleph_1$-Horn theory of $\sigma$-frames]\label{exa:theoryofsigmaframe}
  The $\aleph_1$-algebraic theory $\omega\mbb F$ of \emph{$\sigma$-frames} is of vital importance to later applications in synthetic domain theory (cf.~\zcref{ch:domain}). It is an $\aleph_1$-algebraic theory extending $\omega\mbb J$, with an additional constant $1$, a binary operator $\wedge$, and the following additional axioms:
  \begin{itemize}
    \item $(1,\wedge)$ forms an idempotent commutative monoid;
    \item The absorption laws hold, making $(1,\wedge,0,\vee)$ into a lattice,
    \begin{align*}
      &1 \vee x = 1, \quad x \vee (x \wedge y) = x, \\
      &0 \wedge x = 0, \quad x \wedge (x \vee y) = x.
    \end{align*}
    \item The countable distributivity law holds,
    \[ x \wedge \bigvee_{i<\omega}y_i = \bigvee_{i<\omega}x \wedge y_i. \]
    
  \end{itemize}
  Algebras for $\omega\mbb F$ are posets with countable joins and finite meets, where finite meets distributes over countable joins. Homomorphisms between them are functions preserving finite meets and countable joins. The category of $\sigma$-frames will be denoted as $\sFrm$.
\end{example}

\begin{example}[Coslices of $\kappa$-Horn theories]\label{exm:coslicehorn}
  $\kappa$-Horn theories are closed under coslices in the following sense. Let $\T$ be a $\kappa$-Horn theory with an $S$-sorted signature $\Sigma$. Let $A$ be a $\Sigma$-structure. One may construct a new $\kappa$-Horn theory $A \ltimes \T$ under a new $S$-sorted signature $A+\Sigma$, where for each sort $s\in S$ and each element $a\in A_s$, a new constant $c_a$ will be added for the sort $s$. The axioms of $A \ltimes \T$ is the union of $\T$ with the closed equations that are valid in $A$. This way, an $A \ltimes \T$-algebra is a pair $(B,f)$, with $B$ a $\T$-algebra together with a homomorphism $f \colon A \to B$. Hence, the natural functor $\mmod{A\ltimes\mbb T} \to A/\mmod{\T}$ will be an equivalence.
\end{example}

Finally, we record the following general results for $\kappa$-Horn theories. Let $\mmod{\T}_\kappa$ be the full subcategory of $\kappa$-presentable objects in $\mmod{\T}$. Recall we also write $\kappa\hp\Lex$ as the 2-category of categories with $\kappa$-small limits.

\begin{theorem}\label{thm:factsaboutkappahorntheory}
  Let $\T$ be an $\kappa$-Horn theory in an $S$-sorted signature $\Sigma$. The following holds:
  \begin{enumerate}
    \item $\mmod{\T}$ is locally $\kappa$-presentable: $\mmod{\T}_\kappa$ has $\kappa$-small colimits and there is a canonical equivalence 
    \[ \mmod{\T} \simeq \kappa\hp\Lex(\mmod{\T}_\kappa\op,\Set). \]
    \item $\mmod{\T}$ has free algebras, in the sense that the forgetful functor to $\Set^S$ has a left adjoint,
    \[\begin{tikzcd}
      {\mmod{\T}} & {\Set^S}
      \arrow[""{name=0, anchor=center, inner sep=0}, "U"', curve={height=18pt}, from=1-1, to=1-2]
      \arrow[""{name=1, anchor=center, inner sep=0}, "T"', curve={height=18pt}, from=1-2, to=1-1]
      \arrow["\dashv"{anchor=center, rotate=-90}, draw=none, from=1, to=0]
    \end{tikzcd}\]
    \item $\mmod{\T}$ is closed under $\kappa$-filtered limits in $\mmod{\Sigma}$.
  \end{enumerate}
\end{theorem}
\begin{proof}
  See~\citet[Rem. 3.32]{adamek1994locally}. 
\end{proof}

The most important fact about Horn theories is that it admits the construction of \emph{free} algebras:

\begin{lemma}\label{lem:freealgebraexist}
  For any type $X$ and any family $\varphi : R \to \eqn{\Sigma}(X)$ of equations, we can construct a $\T$-model $T\gr XR$, which has the following universal property,
  \[ \mmod\T(T\gr XR,A) = \scomp{f : X \to A}{\fa rR A,f \models \varphi_r}. \]
\end{lemma}
\begin{proof}
  We define the algebra $T\gr XR$ as a set-quotient of $\tm{\Sigma}(X)$. We consider the smallest relation $\sim$ on $\tm{\Sigma}(X)$ satisfying
  \begin{itemize}
    \item For any $r:R$, $\sim$ contains $\varphi_r$;
    \item For all operation $\sigma : \Sigma$ and any $f,g : a\sigma \to \tm{\Sigma}(X)$, 
    \[ \fa{x}{a\sigma} fx \sim gx \nt \sigma f \sim \sigma g, \]
    or in other words $\sim$ is a congruence;
    \item For all $\Phi := (U,V,\set{\varphi_u:=\pair{s_u,t_u}}_{u:U},\varphi:=\pair{s,t})$ in $\T$ and all $f : V \to \tm{\Sigma}(X)$,
    \[ \fa uU \ass{s_u}_f \sim \ass{t_u}_f \nt \ass{s}_f\sim\ass{t}_f. \]
  \end{itemize}
  $T\gr XR$ is then the set-quotient $\tm{\Sigma}(X)/\sim$. The fact that $T\gr XR$ is a $\T$-model and it satisfies the universal property directly follows from the construction.
\end{proof}

In fact, the above free construction also indicates how to construct general colimits by taking the congruence. The argument is standard, and we record that

\begin{proposition}\label{prop:hornmodelaccessunderset}
  For any $\ms U$-Horn theory $\T$, $\mmod\T$ is complete and cocomplete. The forgetful functor creates limits and $\ms U$-filtered colimits.
\end{proposition}
\begin{proof}
  We verify that $U : \mmod\T \to \Set$ preserves $\ms U$-filtered colimits. Suppose we have a diagram $F : \mc I \to \mmod\T$ with $\mc I$ $\ms U$-filtered. We can define a $\Sigma$-structure on the colimit $A := \ct_{i:\mc I}UF$. For $\sigma : \Sigma$, since $a\sigma : \ms U$ is 1-projective, there merely exists $f : a\sigma \to \sum_{i:\mc I_1}Fi$. This gives us a $\ms U$-small diagram in $\mc I$ indexed by $a\sigma : \ms U$, thus there merely exist a cocone $(k,\alpha)$. This way, we can define $\ass{\sigma}([f])$ to be $[(k,\ass{\sigma}_k(\ld x{a\sigma}F\alpha_x(fx)))]$. The point here is that the element $\ass\sigma([f])$ is \emph{uniquely} defined, thus does not depend on the choice of $(k,\alpha)$, which justifies turning mere existence into an actual existence. 

  To see the above $\Sigma$-structure is a $\T$-model, consider any Horn clause $\Phi := (U,V,\set{\varphi_u}_{u:U},\varphi)$ in $\T$. Suppose we have a family $f : V \to A$ with $A,f \models \varphi_u$ for all $u : U$. Completely similarly, since both $U$ and $V$ are $\ms U$-small, there merely exists a cocone $(k,\alpha)$ where the terms $\ass{s_u}_f$ and $\ass{t_u}_f$ are identified. Since $F_k$ is a $\T$-model, it follows that $A,f \models \varphi$ as well. The universal property can be verified directly.
\end{proof}

We want to furthermore show that $\mmod\T$ is locally $\ms U$-presentable. For this we study $\ms U$-presentable $\T$-models.


\section{Presentable Algebras}\label{sec:presentable_algebras}

Let $\T$ be a $\ms U$-Horn theory for some notion of size $\ms U$. In this section we will characterise $\ms U$-presentable objects in $\mmod\T$. 

\begin{definition}
  We say a $\T$-model $A$ is \emph{$\ms U$-presented}, if there merely exists $\ms U$-small $X,R$ such that $A = T\gr XR$. The full subcategory of $\ms U$-presented algebras will be denoted as $\mmod\T^{\ms U}$.
\end{definition}

\begin{lemma}
  $\ms U$-presented $\T$-models are $\ms U$-presentable.
\end{lemma}
\begin{proof}
  Let $A$ be a $\ms U$-presented model. Since we are showing a proposition, we assume $A = T\gr XR$ for some $\ms U$-small $X,R$. Let $F : \mc I \to \mmod\T$ be a $\ms U$-filtered diagram. Consider the canonical map
  \[ \ct_{i:\mc I}\mmod\T(A,Fi) \to \mmod\T(A,\ct_{i:\mc I}Fi). \]
  To show it is injective, suppose we have $[(i,f)]$ and $[(j,g)]$ in the domain which are identified in the codomain. By construction, this means that
  \[ \fa xX\ex{k}{\mc I_1}\ex{\alpha}{\mc I(i,k)}\ex{\beta}{\mc I(j,k)} F\alpha(fx) = F\beta(gx). \]
  We may assume $X$ to be 1-projective, thus there merely is $k : X \to \mc I_1$ with $\alpha,\beta$ satisfying the above condition. In particular, $[(i,f)] = [(j,g)]$.

  To show surjectivity, suppose we are given a map in the codomain, which we may view as a map $f : X \to \ct_{i:\mc I}Fi$ that respects $R$. Again by 1-projectivity, we may first find a cocone $k$ that $f$ factors through and also respects $R$.
\end{proof}

\begin{lemma}
  $\mmod\T^{\ms U}$ is closed under $\ms U$-small colimit.
\end{lemma}
\begin{proof}
  By Proposition~\ref{prop:smallcompiffeqcoandeq} it suffices to show $\mmod\T^{\ms U}$ has $\ms U$-small equivariant colimit and coequalisers. For a diagram $F : I \to \mmod\T^{\ms U}$ from a $\ms U$-small set $I$, to show $\ct_{i:I}Fi$ is $\ms U$-presented, by 1-projectivity and the classifier 1-type $U$, we may assume there merely exists for any $i : I$ a presentation $Fi = T\gr{X_i}{R_i}$ with $X_i,R_i$ being $\ms U$-small. Then the colimit is given by $T\gr{\sum_{i:I}X_i}{\sum_{i:\mc I}R_i}$, which is again $\ms U$-presented since $\ms U$ is closed under dependent sum.

  For coequaliser, suppose we have $f,g : T\gr XR \to T\gr YS$. The coequaliser is the same when composed with the quotient $T[X] \surj T\gr XR$, thus we may assume $R$ is empty, and $f,g : X \to T\gr YS$. In this case, by 1-projectivity of $X$ we merely get $\qsi f,\qsi g : X \to \tm{\Sigma}(Y)$, and we may view $X$ as the set of equations
  \[ x \mapsto \pair{\qsi f,\qsi g} : X \to \eqn\Sigma(Y). \]
  This way, the coequaliser is $T\gr{Y}{S+X}$.
\end{proof}

\begin{lemma}
  For any $X : \mmod\T$, it is the $\ms U$-filtered colimit of the canonical diagram $\mmod\T^{\ms U}/X$.
\end{lemma}
\begin{proof}
  Since $\mmod\T^{\ms U}$ has $\ms U$-small colimits, the diagram $\mmod{\T}^{\ms U}/X$ is indeed $\ms U$-filtered. The map $\ct_{A:\mmod\T^{\ms U}/X}A \to X$ is evidently surjective. Injectivity follows from the fact that $\mmod\T^{\ms U}/X$ has pushouts.
\end{proof}

\begin{theorem}\label{thm:hornmodelpresentable}
  For any $\ms U$-Horn theory $\T$, the category $\mmod\T$ is locally $\ms U$-presentable, and $\mmod\T_{\ms U} = \mmod\T^{\ms U}$.
\end{theorem}
\begin{proof}
  We have already seen that $\mmod\T^{\ms U} \subseteq \mmod\T_{\ms U}$. On the other hand, if $A$ is $\ms U$-presentable in $\mmod\T$, then the identity on $A$ merely factors through the colimit diagram $\mmod\T^{\ms U}/A$, thus it must be $\ms U$-presented.
\end{proof}

\begin{example}[Presentable objects in $\Set$]\label{exa:presentableobjectsinset}
  Note that by viewing $\Set$ as $\mmod{\mbb O}$, a set is $\ms U$-presented iff it is the set-coequaliser of a pair $R \rightrightarrows X$ for some $\ms U$-small types $X,R$. This implies that $\Set_{\fst} = \Fin$, and $\Set_{\cnt}$ coincide with coequalisers of countable sets.
\end{example}


\section{Classifying topoi for Horn theories with higher arity}


Let $\kappa$ be a regular cardinal and $\T$ a $\kappa$-Horn theory. $\T$ is not a geometric theory in the usual sense, since it contains operations with possibly higher arity and axioms possibly using infinitary conjunctions. However, it is a \emph{$\kappa$-geometric theory} in the sense of~\citet[\S8.3]{makkai1977first}. In fact, via functorial semantics, we can define its models in an arbitrary topos $\mc E$ as follows; cf.\ \emph{loc.\ cit.}:

\begin{definition}[Models of a $\kappa$-Horn theory in topoi]
  The category of $\T$-models in a topos $\mc E$ is defined as follows,
  \[ \mmod{\T}(\mc E) \coloneq \kappa\hp\Lex(\mmod{\T}_\kappa\op,\mc E). \]
\end{definition}

However, the 2-functor $\mc E \mapsto \mmod{\T}(\mc E)$ is \emph{not} necessarily representable in the 2-category $\Top$ of topoi. In other words, $\T$ may fail to have a classifying topos in the usual sense, since it may contain operations of higher arity, and the usual inverse images of geometric morphisms only preserve finite, instead of \emph{$\kappa$-small}, limits. However, it has a \emph{classifying $\kappa$-topos} following~\citet{espindola_every_nodate}:

\begin{definition}[$\kappa$-topos]
  A topos $\mc E$ is a \emph{$\kappa$-topos} if it is an $(\infty,\kappa)$-coherent category in the sense of~\citet[Def.\ 2.7]{espindola_every_nodate}. A \emph{$\kappa$-geometric morphism} between $\kappa$-topoi is a geometric morphism whose inverse image preserves $\kappa$-small limits. The 2-category of $\kappa$-topoi will be denoted as $\kappa\hp\Top$.
\end{definition}

For us, the precise definition of $\kappa$-topos does not matter that much -- we need them simply to state the universal property of the classifying $\kappa$-topos of $\T$ we are going to introduce shortly. What matters more is the following family of examples:

\begin{example}[Presheaf topoi are $\kappa$-topoi]\label{exm:presheafiskappatopos}
  Every presheaf topos is a $\kappa$-topos for any $\kappa$; see~\citet[Exm.\ 3.3]{espindola_every_nodate}.
\end{example}

\begin{definition}[The classifying $\kappa$-topos for a $\kappa$-Horn theory]\label{def:kappaclassifying}
  We define the classifying $\kappa$-topos for $\T$ to be the following presheaf topos,
  \[ \Set[\T_\kappa] \coloneq [\mmod{\T}_\kappa,\Set]. \]
  By~\zcref{exm:presheafiskappatopos}, $\Set[\T_\kappa]$ is indeed a $\kappa$-topos.
\end{definition}

\begin{remark}[The notion for the classifying $\kappa$-topos]
  Note here we have chosen to denote the classifying $\kappa$-topos as $\Set[\T_\kappa]$. This choice reflects the fact that it is the classifying gadget when we view $\T$ as a \emph{$\kappa$-geometric theory}, thus the subscript $\kappa$ is on $\T$. Of course, by definition for any regular cardinal $\lambda > \kappa$, $\T$ is also $\lambda$-ary, and its classifying topos as a \emph{$\lambda$-Horn theory} will be different, and enjoys a universal property in a different 2-category, viz.\ $\lambda\hp\Top$.
\end{remark}

The classifying $\kappa$-topos of $\T$ then enjoys the following universal property:

\begin{theorem}[Representability of the 2-functor of models for a $\kappa$-Horn theory]\label{thm:classtoposofkappahorn}
  There is a natural equivalence for any $\kappa$-topos $\mc E$, 
  \[ \mmod{\T}(\mc E) \simeq \kappa\hp\Top(\mc E,\Set[\T_\kappa]). \]
\end{theorem}
\begin{proof}
  This is~\citet[Thm.\ 3.15]{espindola_every_nodate}.
\end{proof}

\begin{remark}[The need for $\kappa$-topos]
  One may wonder whether the universal property stated in~\zcref{thm:classtoposofkappahorn} may hold for a larger 2-category $\kappa\hp\Top$, viz.\ the 2-category of \emph{all} topoi with geometric morphisms whose inverse images preserve $\kappa$-small limits. However, in essence~\zcref{thm:classtoposofkappahorn} requires the following equivalence: a functor $F \colon \mmod{\T}_\kappa\op \to \mc E$ preserves $\kappa$-small limits iff its left Kan extension $\lan_\yon F \colon \Set[\T_\kappa] \to \mc E$ does so. This does \emph{not} hold for all topoi, hence restricting to $\kappa$-topoi is necessary for~\zcref{thm:classtoposofkappahorn} to hold.
\end{remark}

At the end of this section we note that the classifying $\kappa$-topos $\Set[\T_\kappa]$ of $\T$ has many nice properties. Recall from~\citet{johnstone_remarks_2011} the notion of a \emph{totally connected geometric morphism}, which is a geometric morphism whose inverse image has a further left exact left adjoint, hence in effect having a \emph{right adjoint} in the 2-category of topoi. $\Set[\T_\kappa]$ in fact satisfies the following stronger condition:

\begin{proposition}
  $\Set[\T_\kappa]$ is \emph{$\kappa$-totally connected}, in the sense that the constant sheaf functor $\Delta \colon \Set \to \Set[\T_\kappa]$ has a left adjoint $\Pi \colon \Set[\T_\kappa] \to \Set$ that preserves $\kappa$-small limits. In other words, the $\kappa$-geometric morphism $\Gamma \colon \Set[\T_\kappa] \to \Set$ has a right adjoint in $\kappa\hp\Top$.
\end{proposition}
\begin{proof}
  Since $\Set[\T_\kappa]$ is a presheaf category, the constant sheaf functor $\Delta$ is simply the constant functor, thus having a left adjoint $\Pi$ by computing the colimit as follows,
  \[ \Pi X \cong \ct_{A\in\mmod{\T}_\kappa}X(A). \]
  The indexing category $\mmod{\T}_\kappa$ has $\kappa$-small colimits, thus in particular is $\kappa$-filtered. It follows that $\Pi$ preserves $\kappa$-small limits.
\end{proof}

\section{Lifting of universes in presheaf categories}\label{sec:HSlifting}

Working towards the synthetic quasi-coherence principle in the classifying $\kappa$-topos $\Set[\T_\kappa]$, we need to make sense of Horn theories with higher arity \emph{internally} to the presheaf topos $\Set[\T_\kappa]$. For this, we need a working notion of \emph{universe of size} in $\Set[\T_\kappa]$, to replace the role played by $\kappa$ in $\Set$.

A convenient technical tool is the \emph{universe lifting construction} specified by~\citet{10.1093/oso/9780198501275.003.0008}; for a more categorical understanding of this construction, see~\citet{Awodey_2024} and~\citet{lmcs:17229}. This construction describes how to lift a universe within the base category $\Set$, to any presheaf topos.

For simplicity, by a \emph{universe} in $\Set$, we simply mean a small full subcategory $\mc U$ of $\Set$.\footnote{For us we do not need our universe $\mc U$ to be closed under dependent function types, thus we can work fully with \emph{small} categories, instead of assuming a Grothendieck universe or an inaccessible cardinal.} It comes equipped with the structure of the classifier of $\mc U$-small presheaves, as specified by the fibration 
\[ \pi_{\mc U} \colon \mc U_* \to \mc U, \]
where $\mc U_*$ is the category of \emph{pointed} $\mc U$-small sets.

Let $\mc C$ be a small category. The lifted universe $\pf{\mc U}$ in the presheaf category $\psh(\mc C)$ is defined as the following presheaf,
\[ \pf{\mc U}(c) \coloneq \Cat((\mc C/c)\op,\mc U), \]
together with the discrete fibration $\pf{\mc U_*}$ above it as 
\[ \pf{\mc U_*}(c,X) \coloneq A(\id_c). \]
Here we have viewed the map $\pf{\mc U_*} \to \pf{\mc U}$ in $\psh(\mc C)$ as a presheaf over the category of elements $\elem_{\mc C}\mc U$ of the presheaf $\mc U$.

The important fact about the universe lifting construction is that it preserves the relevant logical properties of the universe; for us, the most important structure is dependent sums:

\begin{proposition}\label{prop:HSpreservesigma}
  The Hofmann-Streicher universe lifting construction preserves universes closed under dependent sum types.
\end{proposition}
\begin{proof}
  See~\citet[\S4.9]{10.1093/oso/9780198501275.003.0008}; though the proof there is written for the groupoid model, the strategy clearly works for any presheaf category.
\end{proof}

\section{Synthetic quasi-coherence for Horn theories with higher arity}

In this section we prove the main theorem of this chapter, viz.\ the synthetic quasi-coherence principle in the classifying $\kappa$-topos $\Set[\T_\kappa]$ for a $\kappa$-Horn theory $\T$. Our proof strategy here will essentially follow that of finitary Horn theories stated in~\citet{blechschmidt_general_2020}, though with appropriate generalisations towards higher arity.

Let $\pf{\kappa}$ be the lifted universe of $\kappa$-small sets in the presheaf category $\Set[\T_\kappa]$. This way, internally in $\Set[\T_\kappa]$ we can speak of a \emph{$\pf{\kappa}$-Horn theory}, by interpreting the specifications in~\zcref{sec:higherhorntheory} internally in $\Set[\T_\kappa]$, with $\kappa$ replaced by $\pf{\kappa}$ as the chosen universe. In particular, the external theory $\T$ can be pulled back, via the constant sheaf functor $\Delta$, to an internal $\pf\kappa$-Horn theory in $\Set[\T_\kappa]$. 

Also recall that there is a generic $\T$-model in the classifying $\kappa$-topos $\Set[\T_\kappa]$, which we denote as $V_\T$. As a model, its interpretation of the sort $s\in S$ in the signature is simply the evaluation functor at $s$: for any $A\in\mmod\T_\kappa$, 
\[ (V_\T)_s(A) \coloneq A_s. \]
In other words, $V_\T$ is the corepresentable functor on the free $\T$-algebra generated by a single element in sort $s\in S$; in accordance with the notation of the left adjoint in~\zcref{thm:factsaboutkappahorntheory}, we write this free algebra as $T[\ms x_s]$, for any $A\in\mmod{\T}$,
\[ \mmod{\T}(T[\ms x_s],A) \cong A_s \cong (V_\T)_s(A). \]
$V_\T$ being an internal $\T$-algebra follows from the fact that each $V_\T(A) = A$ is a $\T$-algebra.

The next step is to realise that there is an internal $\pf\kappa$-Horn theory in $\Set[\T_\kappa]$ of \emph{$V_\T$-algebras}, obtained by applying the construction in~\zcref{exm:coslicehorn} to the $\T$-model $V_\T$. We will denote this internal theory as $\VT$. We will use $\mmod\VT$ to denote the locally small internal category of $\VT$-models in $\Set[\T_\kappa]$. In particular, for any $\VT$-models $M,N$ in $\Set[\T_\kappa]$, there is an object of homomorphisms between them
\[ {\mmod\VT}(M,N) \coloneq \scomp{f \colon M \to N}{f \text{ is a $\T$-homomorphism over $V_\T$}}, \]
which is an internal object in $\Set[\T_\kappa]$. We also use $\geo{\mmod\VT}$ to denote the \emph{externalised} category over $\Set$, which is the image along the global section functor $\Gamma \colon \Set[\T_\kappa] \to \Set$. 

It is easy to observe that there exists a canonical functor internalising an external $\T$-model to an internal model of $\VT$,
\[ \uv - \colon \mmod\T \to \geo{\mmod\VT}. \]
Given any $\T$-algebra $X$, the interpretation of $\uv{X}$ at sort $s\in S$ is given by the following presheaf: for any $A\in\mmod{\T}_\kappa$,
\[ \uv X_s(A) \coloneq (X \otimes A)_s, \]
where $\otimes$ is the \emph{coproduct} in $\mmod\T$. This admits an evident homomorphism $V_\T \to \uv X$, where levelwise is given by the canonical coproduct inclusion $A \to X \otimes A$.

\begin{proposition}\label{prop:internalrep}
  $\uv- \colon \mmod\T \to \geo{\mmod\VT}$ makes $\mmod\T$ a coreflexive subcategory of $\geo{\mmod\VT}$, with the right adjoint retraction given by the global section functor
  \[ \Gamma \colon \geo{\mmod\VT} \to \mmod\T, \]
  which takes a $\VT$-model $M$ in $\Set[\T_\kappa]$ to its global point.
\end{proposition}
\begin{proof}
  We first observe that $\Gamma$ is equivalent to evaluating at $T[\emp]$, the \emph{initial} $T$-model, since the corepresentable functor takes $T[\emp]$ to the terminal object in $\Set[\T_\kappa]$. Hence, given any $f \colon \geo{\mmod\VT}(\uv X,M)$, by applying $\Gamma$ we get
  \[
  \begin{tikzcd}
    {\uv X} && M && X && {\Gamma M} \\
    & {V_\T} &&& {} & T
    \arrow["f", from=1-1, to=1-3]
    \arrow["\mapsto"{description}, draw=none, from=1-3, to=2-5]
    \arrow["{\Gamma f}", from=1-5, to=1-7]
    \arrow[from=2-2, to=1-1]
    \arrow[from=2-2, to=1-3]
    \arrow[from=2-6, to=1-5]
    \arrow[from=2-6, to=1-7]
  \end{tikzcd}
  \]
  Since $\Gamma$ preserves limits, $\Gamma f$ is again a homomorphism between $\T$-models. On the other hand, given $g \colon X \to \Gamma M$, we first obtain a $\T$-homomorphism for any $A \colon \mmod\T\fp$,
  \[ M? \circ g \colon \mmod\T(X,M(A)), \]
  where $? \colon \mmod\T(T,A)$ is the unique map from the initial object. This way, we obtain a map
  \[ [M?\circ g,M_A] \colon \uv X(A) = X \otimes A \to M(A), \]
  where $M_A \colon A \to M(A)$ is given by the map $V_\T \to M$. Verifying these are inverses to each other is routine. 
\end{proof}

\begin{corollary}\label{cor:cpalgebraabsolute}
  $\uv-$ restricts to an equivalence
  \[ \uv - \colon \mmod\T\cp = \geo{\mmod\VT\cp}. \]
\end{corollary}
\begin{proof}
  Since $\uv -$ is a coreflexive embedding, it is automatically fully faithful. Thus, it suffices to show it is essentially surjective, which holds by absoluteness of countable sets in $\Set[\T]$.
\end{proof}

For any $\VT$-model $M$, we define its internal spectrum as follows,
\[ \spec M \coloneq \mmod\VT(M,V_\T). \]

\begin{lemma}\label{lem:represpecrepre}
  For any $A \colon \mmod\T\fp$, we have
  \[ \spec \uv A = \yon^A. \]
\end{lemma}
\begin{proof}
  By construction, for any $B:\mmod\T\fp$,
  \begin{align*}
    \spec \uv A(B) 
    &= \Set[\T](\yon^B,\mmod\VT(\uv A,V_\T)) \\
    &= \geo{\mmod\VT}(\uv A,V_\T^{\yon^B}) \\
    &= \mmod\T(A,B) = \yon^A(B)
  \end{align*}
  The last steps is due to \zcref{prop:internalrep}, and that $\Gamma V_\T^{\yon^B} = V_\T(B) = B$. It follows that $\spec A = \yon^A$.
\end{proof}

\begin{lemma}\label{lem:externalqc}
  For any $A \colon \mmod\T\fp$, the transpose of the evaluation map
  \[ \uv A \to V_T^{\spec \uv A} \]
  is an equivalence between $\VT$-models.
\end{lemma}
\begin{proof}
  For any $B:\mmod\T\fp$, we have
  \[ V_\T^{\spec\uv A}(B) = \Set[\T](\yon^A \times \yon^B,V_\T) = A \otimes B = \uv A(B). \]
  This holds since $\yon^A \times \yon^B = \yon^{A \otimes B}$.
\end{proof}

\begin{theorem}[Quasi-Coherence]\label{thm:qc}
  For any $X \colon \mmod\VT\fp$, the transpose of the evaluation map
  \[ X \to V_\T^{\spec X} \]
  is an equivalence between $\VT$-models.
\end{theorem}
\begin{proof}
  We apply \zcref{lem:externalqc} to the internal theory $\VT$ in $\Set[\T]$. Let $\Set[\T][\VT]$ be the classifying topos of $\VT$ over $\Set[\T]$, and let $V_{V_\T}$ be the generic $\VT$-model in $\Set[\T][\VT]$. By \zcref{lem:externalqc}, we have an equivalence in $\Set[\T][\VT]$
  \[ \uv X = V_{\VT}^{\spec \uv X}, \]
  and by \zcref{lem:represpecrepre}, $\spec\uv X = \yon^{X}$. Now if we apply the global section $\Gamma \colon \Set[\T][\VT] \to \Set[\T]$, on one hand we get
  \[ \Gamma\uv X = \uv X(V_\T) = X \otimes V_\T = X, \]
  and on the other hand, by the computation in \zcref{lem:externalqc},
  \[ \Gamma V_{\VT}^{\spec\uv X} = V_{\VT}^{\spec\uv X}(V_\T) = X \otimes V_\T = X. \]
  This completes the proof.
\end{proof}

\begin{example}\label{exm:pshdlqc}
  Let $\mbb D$ be the theory of distributive lattices, and let $\DL$ be the category of distributive lattices. By \zcref{thm:finitedualityfordl}, we have a duality $\DL\fp\op \simeq \Pos\fp$, which implies
  \[ \Set[\mbb D] \simeq \psh(\Pos\fp). \]
  Thus internally in $\Set[\mbb D]$, we have quasi-coherence for finitely presentable distributive lattices.
\end{example}

\begin{example}\label{exm:pshdmqc}
  Let $d\mbb{M}$ be the theory of de Morgan lattice, and let $d\mb{M}$ be the category of de Morgan lattices. This theory is still locally finite, thus all finitely presented de Morgan lattices are again finite.
\end{example}

\begin{example}\label{exm:pshhaqc}
  Let $\mbb H$ be the theory of Heyting algebras, and let $\HA$ be the category of Heyting algebras. The difference in this case is that finitely presented Heyting algebras are no longer finite.
\end{example}

\begin{example}\label{exm:pshbaqc}
  Let $\mbb B$ be the theory of Boolean algebras, and let $\BA$ be the category of Boolean algebras. Similarly, we have
  \[ \Set[\mbb B] \simeq \func(\BA\fp,\Set) \simeq \psh(\Fin). \]
\end{example}

\begin{example}\label{exm:pshmlqc}
  Let $\mbb J$ be the theory of join-semi-lattices, and let $\JL$ be the category of join-semi-lattices. The category $\JL\fp$ is in fact self-dual,
  \[ \Set[\mbb J] \simeq \func(\JL\fp,\Set) \simeq \psh(\JL\fp). \]
  Completely similarly we can define the theory $\mbb M$ of meet-semi-lattices. The theory $\mbb M$ and $\mbb J$ are isomorphic, hence they generates equivalent models and classifying topoi.
\end{example}

\begin{example}\label{exm:pshcnqc}
  Let $\mbb C$ be the theory of \emph{connection algebra}. A model of $\mbb C$ is a set equipped with two monoid structure $(0,\sqcup)$ and $(1,\sqcap)$, such that $0 \sqcap x = 0 = 0 \sqcap x$ and $1 \sqcup x = 1 = x \sqcup 1$. In particular, this is not necessarily a lattice.
\end{example}


\section{Transferring to local subuniverses}

\section{A resulting internal duality}\label{sec:internalduality}
\chapter{The internal logic of classifying topoi for coherent theories}

\section{The path topos of a coherent topos}

\section{Synthetic quasi-coherence for coherent theories}


\part{The applications of synthetic quasi-coherence}\label{prt:appl}
\NewDocumentCommand\AxiomSQCI{}{SQCI}
\NewDocumentCommand\PrintAxiomSQCI{}{
  \begin{axiom}[\AxiomSQCI]
    $\mbb{I}$ is stably spatial, i.e.\ the slice $\mbb{I}/i$ is spatial for every $i \in \mbb I$.
  \end{axiom}
}

\NewDocumentCommand\AxiomSQCID{}{SQCID}
\NewDocumentCommand\PrintAxiomSQCID{}{
  \begin{axiom}[\AxiomSQCID]
    Both $\mbb{I}$ and $\mbb{I}\op$ are stably spatial.
  \end{axiom}
}

\NewDocumentCommand\AxiomSQCP{}{SQCP}
\NewDocumentCommand\PrintAxiomSQCP{}{
  \begin{axiom}[\AxiomSQCP]
    The polynomial $\mbb{I}$-algebra $\mbb{I}[\ms{i}]$ is spatial.
  \end{axiom}
}

\NewDocumentCommand\AxiomNT{}{NT}
\NewDocumentCommand\PrintAxiomNT{}{
  \begin{axiom}[\AxiomNT]
    We have $0 \neq 1$ in $\mbb{I}$.
  \end{axiom}
}

\NewDocumentCommand\AxiomL{}{L}
\NewDocumentCommand\PrintAxiomL{}{
  \begin{axiom}[\AxiomL]
    $\mbb{I}$ is local, i.e.\ $0 \neq 1$, and $\ist{i\vee j} \eq \ist{i} \vee \ist{j}$ for $i,j \in \mbb{I}$. 
  \end{axiom}
}

\NewDocumentCommand\AxiomCL{}{cL}
\NewDocumentCommand\PrintAxiomCL{}{
  \begin{axiom}[\AxiomCL]
    $\mbb{I}$ is colocal, i.e.\ $0 \neq 1$ and $\isf{i\wedge j} \eq \isf{i} \vee \isf{j}$ for all $i,j \in \mbb{I}$.
  \end{axiom}
}

\NewDocumentCommand\AxiomSL{}{SL}
\NewDocumentCommand\PrintAxiomSL{}{
  \begin{axiom}[\AxiomSL]
    $\mbb{I}$ is a \emph{strict linear order}, i.e.\ $0 \neq 1$ and $i \le j \vee j \le i$ for all $i,j\in\mbb{I}$. 
  \end{axiom}
}

\NewDocumentCommand\AxiomOneCS{}{1cS}
\NewDocumentCommand\PrintAxiomOneCS{}{
  \begin{axiom}[\AxiomOneCS]
    $\mbb{I}$ is \emph{1-coskeletal}, i.e.\ (\AxiomSL) holds and $i \le j \eq \isf{i} \vee (i = j)\vee \ist{j}$ for all $i,j \in \mbb{I}$.
  \end{axiom}
}

\NewDocumentCommand\AxiomSQCF{}{SQCF}
\NewDocumentCommand\PrintAxiomSQCF{}{
  \begin{axiom}[\AxiomSQCF]
    All finitely presented $\mbb{I}$-algebras are spatial.
  \end{axiom}
}

\NewDocumentCommand\AxiomSQCC{}{SQCC}
\NewDocumentCommand\PrintAxiomSQCC{}{
  \begin{axiom}[\AxiomSQCC]
    All countably presented $\mbb{I}$-algebras are spatial.
  \end{axiom}
}

\chapter{Distributive lattice, sigma-frames, and synthetic domain theory}\label{ch:domain}

This chapter is based on~\citet{sterling2025domains}, which has systematically investigated the new connection between the $\omega$-Horn theory of distributive lattices (cf.~\zcref{exa:finitealgebratheories}) and the $\omega_1$-Horn theory of $\sigma$-frames (cf.~\zcref{exa:theoryofsigmaframe}), and \emph{synthetic domain theory} (\gls{SDT}), through the lens of \gls{QC}.

The proposal to develop a synthetic theory of domains was raised by Dana Scott in the early 1980s. The thesis is that (pre)domains should be viewed simply as some special sets in a suitable universe, and any function between them should automatically respect the computational data associated to domains. Later \citet{hyland_first_2006} gave an extensive list of the properties such a universe might satisfy.

A universe for synthetic domains should contain an interval object $\mbb{I}$. From the perspective of \gls{QC}, as we shall see, a good candidate for such an interval object is the \emph{generic model} of distributive lattices or $\sigma$-frames, in their corresponding classifying topoi. The reason is that, the \emph{logical/algebraic} structure of distributive lattices and $\sigma$-frames, via \gls{QC}, translates exactly to the synthetic behaviours of the interval $\mbb I$ required for a good notion of synthetic domains. 

As a first instance, the interval $\mbb I$ induces an observational preorder (or \emph{specialisation} preorder from a topological perspective) on domains for which functions should be automatically monotone. An internal axiom that characterises the synthetic nature of this fact is the so-called \emph{Phoa principle}: the function space $\mbb{I}^\mbb{I}$ should classify the \emph{order} on $\mbb{I}$. In other words, functions from $\mbb{I}$ to itself are completely determined by the order structure on $\mbb{I}$. 

From the perspective of \gls{QC}, we already know that $\mbb I$ is a spectrum (cf.~\zcref{exm:cubesober}), thus via \gls{QC} for Horn theories the exponential $\mbb I^{\mbb I}$ then can be computed as the \emph{free algebra} $\mbb I[\ms x]$ (but even without \gls{QC}, this observation is already contained in~\citet{RN879}). \citet{gratzer2024directed} has shown that when $\mbb{I}$ is the generic distributive lattice, the isomorphism $\mbb{I}^\mbb{I} \cong \mbb{I}[\ms{x}]$ implies Phoa's principle---because for distributive lattices, the polynomial $\mbb{I}$-algebra $\mbb{I}[\ms{x}]$ always classifies the partial order on $\mbb{I}$. This is a striking example of how the properties of a theory impact the internal logic of its classifying topos, which is the running theme of~\zcref{prt:appl} of this thesis.

Additionally for \gls{SDT}, $\mbb{I}$ should also form a subuniverse of propositions closed under dependent sums, so as to form a dominance in the sense of \citet{rosolini1986continuity}. The dominance structure is used to parametrise partial functions through a \emph{partial map classifier} functor $L$ constructed from $\mbb{I}$. 

Besides these elementary axioms, the main additional axiom of synthetic domain theory is the \emph{chain completeness} of the interval: 
\begin{itemize}
  \item[]
  Let $\omega$ and $\ov\omega$ be the carriers of the internal initial algebra and final coalgebra to the partial map classifier functor $L$ respectively. There is a canonical inclusion $\omega \hook \ov\omega$, and the chain completeness axiom states that $\mbb{I}$ is \emph{right orthogonal} to this inclusion in the sense that the canonically induced restriction map $\mbb{I}^{\ov\omega} \to \mbb{I}^{\omega}$ should be an isomorphism. 
\end{itemize} 
This can be viewed as the synthetic version of the $\omega$-completeness of the Sierpi\'nski space in traditional domain theory.

As we shall see, \gls{QC} for distributive lattices and $\sigma$-frames suffices to explain the behaviour of the above mentioned axioms in synthetic domain theory all at once, including not only the dominance property and Phoa's principle, but also the chain completeness of the interval and the inductive construction of the initial $L$-algebra. In fact, both the dominance property and the Phoa's principle are \emph{equivalent} to certain \gls{QC} axiom (cf.~\zcref{thm:chardominance,obs:sqcp-phoa-gen}).

As a first indication of why \gls{QC} for these theories might be useful for domain theory, we have observed that by taking the interval $\mbb{I}$ to be a distributive lattice or $\sigma$-frame internal to some topos, many important objects for domain theory can be written as \emph{spectra}. We have mentioned the example of $\mbb I$ itself; this includes furthermore the cubes $\mbb{I}^n$ (\zcref{exm:cubesober}), the simplices $\Delta^n$ (\zcref{exm:simplicessober}), and the final coalgebra $\ov\omega$ (\zcref{exm:ovomegasober}). Here we outline the main results of this chapter, and discuss some of their potential applications.

\begin{enumerate}[leftmargin=*]
  \item \emph{\gls{QC} produces \gls{SDT}.} We show that several instances of the principle of \gls{QC} suffice to account for \emph{all} essential axioms for \gls{SDT}. Specifically:
  \begin{enumerate}
    \item\label{contribution:dominance} In the case of both distributive lattices and $\sigma$-frames, the \emph{dominance property} is identified as a \gls{QC} principle with respect to quotient algebras $\mathbb{I}/(i=1)$ in \zcref{prop:Idominance}; the co-dominance principle is dually identified with \gls{QC} with respect to quotient algebras $\mathbb{I}/(i=0)$ in \zcref{cor:dualisdominance}. This leads to an explicit computation of the partial map classifiers of spatial $\mbb{I}$-algebras (\zcref{prop:liftingofalgebra}), and of spectra (\zcref{prop:liftofsober}).
    \item We give an explicit computation of the \emph{observational preorder} (\zcref{defn:specialisation}) for sober spaces (\zcref{lem:specorderofsober}) and spatial $\mbb{I}$-algebras (\zcref{cor:specordonalgiscan}).
    \item A generalisation of Phoa's principle that applies to an arbitrary $\mbb{I}$-algebra $A$ is stated (\zcref{def:gen-phoa}), and is shown to coincide with \gls{QC} with respect to polynomial algebras (\zcref{obs:sqcp-phoa-gen}). As a corollary, we see that spatiality of finitely generated free $\mbb{I}$-algebras is equivalent to Phoa's principle (\zcref{cor:phoa-vs-quasicoherence}).
    \item In the case of $\sigma$-frames, \gls{QC} with respect to countably presented algebras implies chain completeness of $\mbb{I}$ (\zcref{thm:complete}) assuming $\mbb{I}$ is non-trivial. This is the only place where it is important to work with $\sigma$-frames and not distributive lattices (cf.~\zcref{rem:whynotdis}). This gives a full account of the axioms for \gls{SDT} as specified by \citet{hyland_first_2006}. 
  \end{enumerate}

  \item \emph{Local properties of spectra.} Though in general we do not assume the generic interval $\mbb{I}$ to satisfy any further geometric property such as linearity, we do discuss many instances of them and investigate their consequences.
  \begin{enumerate}

    \item \zcref{sec:locality} introduces various local conditions for $\mbb{I}$, all of which will be compatible with \gls{QC}. This showcases the flexibility of this framework to encompass different flavours of domain theory.

    \item  More interestingly, without assuming the local properties globally, one can still show that $\mbb{I}$ will be \emph{right orthogonal} to the maps that classifies these local properties (\zcref{specisnontrivial,specistriangulated,specis1t}). In general, any limiting diagram of spatial algebras will induce a localisation class containing $\mbb{I}$; cf.\ \zcref{rem:limofalgloc}. This exhibits a new type of techniques in reasoning about domains in this framework.
    
    \item As a special case of locality, we also connect to the recent approach of synthetic (higher) category theory~\citep{riehl2017type,buchholtz2021synthetic,gratzer2024directed}. In particular, we show spectra will be synthetic categories (\zcref{thm:soberposet}). As another example, we also show $\omega$ is a synthetic category, in fact it satisfies \emph{all} the orthogonality conditions discussed in this paper except chain completeness (\zcref{thm:omegaortho}).
  \end{enumerate}

  \item \emph{New models for synthetic domain theory.} Consequently, we uncover a large family of new models for synthetic domain theory based on sites induced by countably presented $\sigma$-frames. We will discuss these models in \zcref{sec:newmodelforsdt}, and compare them to the existing sheaf models for synthetic domain theory~\citep{FIORE1997151} in \zcref{subsec:compare}.
\end{enumerate}

\begin{remark}[A note on realisability models of \gls{SDT}]
  As \gls{QC} are based on classifying topoi, our models of \gls{SDT} will be \emph{sheaf models}. The most studied models of \gls{SDT} are in fact those arising from \emph{realisability}---in which the interval $\mbb{I}$ is taken to be Rosolini's dominance of semidecidable propositions in the Effective Topos $\mathbf{Eff}$~\citep{hyland_first_2006,PhoaWesleyKym-Son1991DtiR}, which happens to coincide there with the \emph{initial $\sigma$-frame} by virtue of countable choice. The model of \gls{SDT} in $\mathbf{Eff}$ ultimately interprets domain theoretic concepts in terms of Turing machines, which makes them applicable to the classical theory of computation~\citep{bauer_first_2006}. 

  On the other hand, realisability models do not provide a direct connection to classical theories of domains. For the latter, researchers returned to the original suggestion of Scott and devised sheaf models of \gls{SDT} in which the site is obtained from some existing category of domains or a dense generator thereof~\citep{fiore-plotkin:1996,FIORE1997151,fiore2001domains}.

  Much effort has been put toward identifying a common axiomatics that accounts for both realisability and sheaf models of \gls{SDT}; cf.\ the work of \citet{reus-streicher:1999,simpson:2004}. Beyond these basic axioms, the realisability and sheaf models of \gls{SDT} behave very differently, as exemplified in the failure of the initial $L$-algebra to be inductive in the former, as \citet{VANOOSTEN2000233} pointed out.
\end{remark}

Our motivation for domain theory encourages us to work in a context as general as possible, so as to allow future developments that connect more specific flavours of domain theory existing in the literature. This means that besides \zcref{sec:newmodelforsdt} where we discuss models, throughout the chapter we will work constructively in a language enriched with a generic model $\mbb{I}$. We do not assume any additional assumption globally (even \gls{QC}). Whenever we introduce an assumption under the \textbf{Axiom} environment, it should be viewed as introducing the \emph{content} of that assumption, rather than assuming it directly afterwards. In particular, the development is completely modular, and any additional assumption will be explicitly mentioned in each result.

In fact, we even take a step further. For the first half of this chapter we will not work with distributive lattices specifically. Instead, we assume to work with an arbitrary \emph{propositionally stable} theory in the sense of \zcref{defn:propositional}. In particular, the construction of the dominance and the computation of partial map classifier mentioned in (\ref{contribution:dominance}) earlier in this chapter works in this generality. Only from \zcref{sec:locality} onward will we then work more specifically with theories based on distributive lattices.

\section{Open propositions and the dominance property}\label{sec:dominance}

We start by assuming we have a Horn theory $\T$ (which might be infinite) which is \emph{propositionally stable} in the following sense: 

\begin{definition}[Propositionally stable theory]\label{defn:propositional}
  We say $\T$ is \emph{propositionally stable}, if it extends the theory of meet-semi-lattices (1, $\wedge$), and quotient by truth of an element is computed by \emph{slicing}: for any $\T$-model $A$ and element $a\in A$, the quotient $A/a=1$ is given by the restriction map
  \[ {a \wedge -} \colon A \surj A/a\text{,} \]
  where the slice $A/a$ by definition is ${\cv} a \coloneq \scomp{b\in A}{b\le a}$, with $\le$ the partial order on $A$ induced by the meet-semi-lattice structure.
\end{definition}

\begin{remark}[Examples of propositionally stable theories]
  $\mbb M$ of meet-semi-lattices, $\mbb D$ of distributive lattices, $\mbb H$ of Heyting algebras, and $\mbb B$ of Boolean algebras are all examples of propositionally stable theories. In fact, all finitary quotients of $\mbb H$ or $\mbb B$ will again be propositionally stable. More generally, for any propositionally stable theory $\T$ and any $\T$-model $D$, the theory of $D$-algebras will again be so, because quotients of $D$-algebras are computed the same as quotients of models of $\T$.
\end{remark}

\begin{remark}[The theory of $\sigma$-frames]\label{rem:sigmaframe}
  Recall the $\omega_1$-Horn theory $\sigma\mbb F$ of $\sigma$-frames in~\zcref{exa:theoryofsigmaframe}. Equivalently, a model of $\sigma\mbb F$ is a lattice with finite meets and countable joins, where finite meets distributes over countable joins. A notable fact about $\sigma$-frames is that their finitary behaviours are exactly the same as distributive lattices, in the sense that for any $\sigma$-frame $A$, the finitely presented $\sigma$-frame $A[n]/R$ coincides with the finitely presented \emph{distributive lattice} $A[n]/R$. Thus, there will be no difference between working with $\sigma$-frames or distributive lattice when we work with finitely presented $\mbb{I}$-algebras. The salience of countable joins will emerge only in \zcref{sec:infdomain} when we prove chain completeness of $\mbb{I}$.
\end{remark}

\begin{remark}[Propositionally op-stable theories]\label{rem:opprop}
  There is an evident ``opposite'' notion of being propositionally stable, which requires $\T$ to extend the theory of join-semi-lattices and the quotient $A/a = 0$ for any $a\in A$ is computed by ${a \vee -} \colon A \surj a/A$. Every propositionally stable theory corresponds to a propositionally op-stable theory by dualising: the theories of join-semi-lattices and coHeyting algebras are propositionally op-stable; Boolean algebras and distributive lattices dualise to themselves; the theory of $\sigma$-frames is also propositionally op-stable.
\end{remark}

For a propositionally stable theory $\T$, we think of the generic model $\mbb{I}$ as certain interval object, since it is equipped with a partial order. In this case, more sets can be realised as spectra. The important examples are \emph{simplices}:

\begin{example}[Simplices as spectra]\label{exm:simplicessober}
  For any $n \in \N$, let $\Delta^n \hook \mbb{I}^n$ be the $n$-simplex
  \[ \Delta^n \coloneq \scomp{i \colon n \to \mbb{I}}{i_1 \ge \cdots \ge i_n}\text{.} \]
  This set is indeed a spectrum, since by definition we have
  \[ \Delta^n \cong \spec\mbb{I}[\ms{i}_1,\cdots,\ms{i}_n]/\ms{i}_1\ge\cdots\ge \ms{i}_n\text{.} \]
  To simplify later discussion, we might also introduce the following sets isomorphic to simplices above,
  \[ \Delta_n \coloneq \scomp{i \colon n \to \mbb{I}}{i_1 \le \cdots \le i_n}\text{.} \]
\end{example}

The constant $1 \in \mbb{I}$, which is the top element in $\mbb{I}$, induces a predicate
\[ \istsym\colon \mbb{I} \to \pp\text{,} \]
which takes $i \in \mbb{I}$ to the proposition $i = 1$. Using $\istsym$, we can think of the spectrum $\spec A$ of an $\mbb{I}$-algebra $A$ as its set of \emph{models}, with a satisfaction relation $\models$ as follows:

\begin{definition}[Satisfaction]
  For any $\mbb{I}$-algebra $A$ and elements $a\in A$ and $x\in \spec A$, we say $x$ \emph{satisfies} $a$, written $x \models a$, if $\ist{xa}$.
\end{definition}

For instance, the satisfaction relation can be used to carve out sober subspaces consisting of models satisfying some element in $A$:

\begin{example}
  For any $\mbb{I}$-algebra $A$ and $a\in A$, the subset $D_a\subseteq \spec A$ defined below is again a spectrum:
  \[ D_a \coloneq \scomp{x \in X}{x \models a} \cong \spec A/a\text{.} \]
\end{example}

The first observation is that $\istsym$ takes $i \in \mbb I$ to a spectrum:

\begin{lemma}\label{lem:openpropsober}
  The spectrum $\spec\mbb{I}/i$ for any $i \in \mbb I$ is equivalent to the proposition $\ist{i}$.
\end{lemma}
\begin{proof}
  We have $\spec \mbb{I}/i \cong \mbb{I}\alg(\mbb{I}/i,\mbb{I})$ by definition. Since $\mbb{I}$ is the initial $\mbb{I}$-algebra and $\mbb{I}/i$ is a quotient of $\mbb{I}$, any homomorphism $\mbb{I}/i \to \mbb{I}$ will be unique, hence $\spec\mbb{I}/i$ is a proposition. It thus suffices to show $\spec\mbb{I}/i$ implies $\ist{i}$. Because $\T$ is propositionally stable, an $\mbb{I}$-algebra morphism $f \colon \mbb{I}/i \to \mbb{I}$ is a commutative diagram as follows,
  \[
  \begin{tikzcd}
    \mbb{I}/i \ar[rr, dashed, "f"] & & \mbb{I} \\ 
    & \mbb{I} \ar[ul, two heads, "i \wedge -"] \ar[ur, equal]
  \end{tikzcd}
  \]
  In this case, we have
  \[ i = f(i \wedge i) = f(i) = 1\text{.} \]
  The final identity holds since $f$ preserves the top element. 
\end{proof}

The goal of this section is to show that, under a suitable \gls{QC} assumption, $\istsym\colon \mbb{I} \to \pp$ makes $\mbb{I}$ a \emph{dominance} in the sense of \citet{rosolini1986continuity} (cf.~\zcref{def:dominance}). In fact, at the end of this section we will provide an exact characterisation of the dominance property of $\mbb{I}$ as a principle of \gls{QC}. It turns out that it suffices to require all slices of $\mbb{I}$ to be spatial:

\begin{definition}[Stable spatiality and sobriety]
  An $\mbb{I}$-algebra $A$ is \emph{stably spatial}, if for any $i \in \mbb I$, the slice $A/i$ is spatial. A sober space $X \cong \spec A$ is \emph{stably sober} if $A$ is stably spatial.
\end{definition}

\PrintAxiomSQCI

As a first consequence, we show the following important lemma:

\begin{lemma}[\AxiomSQCI]\label{lem:intconserve}
  The interval $\mbb{I}$ is \emph{conservative}: For any $i,j \in \mbb{I}$
  \[ (\ist{i} \eq \ist{j}) \eq i = j\text{.} \]
\end{lemma}
\begin{proof}
  If $i \le j$ then trivially $\ist{i} \to \ist{j}$. Thus it suffices to show
  \[ (\ist{i} \to \ist{j}) \to i \le j\text{.} \]
  Take $i,j$ with $\ist{i} \to \ist{j}$. By \zcref{lem:openpropsober} and (\AxiomSQCI), each $\ist{i}$ and $\ist{j}$ will be sober. Then we get a restriction map between $\mbb{I}$-algebras,
  \[ \mbb{I}/j \cong \mbb{I}^{\ist{j}} \to \mbb{I}^{\ist{i}} \cong \mbb{I}/i\text{.} \]
  Similar to the proof of \zcref{lem:openpropsober}, the existence of such an $\mbb{I}$-algebra homomorphism implies $i \le j$. 
\end{proof}

Thus, under (\AxiomSQCI), we may view the generic algebra $\mbb{I}$ as a subuniverse of \emph{open} propositions via the embedding $\istsym$:

\begin{definition}[Open proposition]
  We say that a proposition $p$ is \emph{open} when $p$ is equivalent to $\ist{i}$ for some $i$ in $\mbb{I}$.
\end{definition}

Under (\AxiomSQCI), by \zcref{lem:intconserve} the $i$ that $p \eq \ist{i}$ will be \emph{unique}, thus being open is itself a proposition. Also, open propositions are evidently closed under finite conjunctions, since $\istsym$ preserves them. 

More generally, we may define the notion of open subsets:

\begin{definition}[Open subset]
  A subset $U$ of $X$ is \emph{open} if for any $x\in X$ the proposition $x\in U$ is open.
\end{definition}

If $X$ is sober, open subsets are easy to classify:

\begin{lemma}[\AxiomSQCI]\label{lem:openofsobergivesalgebra}
  Let $X \cong \spec A$ be sober. Then any open subset of $X$ is of the form $D_a$ for some $a\in A$.
\end{lemma}
\begin{proof}
  Let $U$ be an open subset of $X$, so that for any $x\in X$ there exists some (by conservativity, necessarily unique) element $i \in \mbb I$ that $x \in U \eq \ist{i}$.
  Hence we have a map $\qsi a \colon X \to \mbb{I}$ with $x \in U \eq \ist{\qsi ax}$. Because $X$ is sober, we may identify $\qsi a$ with an element $a\in A$, where $\qsi ax = xa$. This way,
  \[ \fa xX x \in U \eq x \models a\text{,} \]
  and thus $U$ is the subset $D_a$.
\end{proof}

\begin{proposition}[\AxiomSQCI]\label{prop:Idominance}
  The set of opens $\mbb{I}$ forms a dominance.
\end{proposition}
\begin{proof}
  Since $\istsym\colon \mbb{I} \hook \pp$ by definition is closed under finite meets, it suffices to show that an open subset of an open proposition is again an open proposition. Suppose $p$ is an open proposition and $q$ is an open subset of $p$. By definition $p \eq \ist{i}$ for some $i \in \mbb I$. Since $\ist{i}$ is sober, \zcref{lem:openofsobergivesalgebra} implies that $q$ is $D_j$ for some $j \in \mbb I/i$. Since $\T$ is propositionally stable, equivalently $j$ can be viewed as an element $j \in \mbb I$ with $j \le i$. Hence,
  \[ q \eq D_j \eq \spec(\mbb{I}/i)/j \eq \spec\mbb{I}/j\text{,} \]
  which implies $q$ is also an open proposition.
\end{proof}

\begin{remark}[Dominance without choice]\label{rem:dominancewithoutchoice}
  The related work of \citet{Cherubini_Coquand_Hutzler_2024} contains a similar construction of a dominance from a generic \emph{ring}. However, the proof there relies on certain classical axiom, which they call \emph{Zariski local choice}. The reason that the construction here does not require a similar assumption is precisely due to conservativity of $\mbb{I}$: the dominance is $\mbb{I}$ itself, rather than an image of $\mbb{I}$. In this way, we can transform an open subset of $X$ to a map $X \to \mbb{I}$, which by \gls{QC} can be further identified as an element of the algebra if $X$ is sober. 
\end{remark}

We end this section by showing that $\mbb{I}$ being stably spatial, i.e.\ all slices of $\mbb{I}$ are spatial, is a \emph{characterisation} the dominance property of $\mbb{I}$. Under this assumption, we can prove the (surprising) fact that all spatial algebras will indeed be stably spatial:

\begin{theorem}\label{thm:chardominance}
  The following are equivalent:
  \begin{enumerate}
    \item (\AxiomSQCI) holds, i.e.\ $\mbb{I}$ is stably spatial;
    \item $\mbb{I}$ form a dominance;
    \item All spatial $\mbb{I}$-algebras are stably spatial.
  \end{enumerate}
\end{theorem}
\begin{proof}
  We have shown (1) $\nt$ (2) in~\zcref{prop:Idominance}. Evidently, (3) $\nt$ (1) since by~\zcref{exm:intervalqc} $\mbb{I}$ is spatial. We first show (2) $\nt$ (1). Suppose $\mbb{I}$ is a dominance and take any $i \in \mbb I$. We construct an explicit inverse $s \colon \mbb{I}^{\ist{i}} \to \mbb{I}/i$ by defining $s$ as follows,
  \[ s(j) \coloneq \sum_{p\in\ist i} j(p). \]
  By construction $s(j) \le i$, since if $(p,x) \in s(j)$, then $p \in \ist i$ which implies $i = 1$, hence $x \le i$. To show this is an inverse of $\eta \colon \mbb{I}/i \to \mbb{I}^{\ist i}$, for any $j \le i$,
  \[ s\eta(j) \eq \sum_{p\in\ist i}j \eq j. \]
  On the other hand, for any $j\in\mbb{I}^{\ist i}$, to show $j = \eta s(j)$ by function extensionality it suffices to assume $i = 1$, viz.\ there is $p\in\ist i$. In this case, 
  \[ \eta s(j) = \eta j(p) = \ld{q}{\ist i} j(p). \]
  But in this case, for any $q\in\ist i$, $q = p$, thus $j(q) = j(p)$, which by function extensionality again it follows that 
  \[ \eta s(j) = j. \]
  This shows $s$ is indeed an inverse of $\eta$, hence $\mbb{I}/i$ is spatial.

  Now we show (1) $\nt$ (3). Suppose $A$ is a spatial $\mbb{I}$-algebra. Since the slice $A/i$ is the quotient algebra $A/i = 1$, it follows that $A/i \cong A \otimes \mbb{I}/i$. By~\zcref{specrightadj}, we have
  \[ \spec(A/i) \cong \spec A \times \spec{\mbb{I}/i} \cong \spec A \times \ist i. \]
  Since $\mbb{I}/i$ is spatial, the following canonical isomorphisms hold,
  \[ \opens{\spec(A/i)} \cong (\mbb{I}/i)^{\spec A} \cong \scomp{U\in\opens{\spec A}}{\fa x{\spec A} U(x) \le i}. \]
  Since $A$ is spatial, $A \cong \opens{\spec A}$. In fact, spatiality also implies that there are ``enough points'' in the spectrum to detect its order (cf.~\zcref{lem:canon-beh-coincide}),
  \[ \opens{\spec(A/i)} \cong \scomp{a\in A}{\fa x{\spec A} x(a) \le i} \cong \scomp{a\in A}{a \le i} \cong A/i. \]
  This shows $A/i$ is also spatial.
\end{proof}

\section{Partial map classifier}\label{sec:lifting}
Given the interval $\mbb{I}$, we can construct internally the partial map classifier. For any set $X$, it is given by
\[ L X \coloneq \sum_{i \in \mbb I}X^{\ist{i}}\text{.} \]
The functoriality is easy to express: For any $f \colon X \to Y$, we have
\[ L f(i,x) \coloneq (i,\ld{w}{\ist{i}}fxw)\text{.} \]
There is an evident unit $\eta_L \colon X \to L(X)$, where
\[ \eta_L(x) \coloneq (1,\ld\_ 1 x)\text{.} \]
The unit $\eta_L$ then classifies the partial maps out of $X$ with an \emph{open} domain. If $\mbb{I}$ forms a dominance, then there is also a multiplication $\mu \colon L(L X) \to L X$, where $\mu$ takes any $(i,u)$ with $i \in \mbb I$ and $u \colon \ist{i} \to L X$ first to $(j,x)$, where $j$ is the dependent sum
\[ j \coloneq \sum_{w\in\ist{i}} (uw)_0\text{,} \]
and $x \colon \ist{j} \to X$ is the partial element such that for $w \in \ist{i}$ and $v \in (uw)_0$
\[ x(w,v) \coloneq (uw)_1(v)\text{.} \]

\begin{example}
  By definition, it is easy to see that
  \[ L 1 \cong \sum_{i \in \mbb I}\ist{i} \cong \mbb{I}\text{.} \]
\end{example}

For synthetic domain theory, the object of particular importance is the partial map classifier of the interval $\mbb{I}$ itself. The following computation works more generally for all stably spatial algebras:

\begin{proposition}\label{prop:liftingofalgebra}
  If $A$ is stably spatial, then we have
  \[ L A \cong \sum_{i \in \mbb I}A/i\text{.} \]
\end{proposition}
\begin{proof}
  By assumption for $i \in \mbb I$, the quotient $A/i$ is again spatial. Now notice that we do have
  \[ A/i \cong A \otimes \mbb{I}/i\text{.} \]
  By \zcref{specrightadj} and spatiality of $A/i$ and $A$,
  \[ A/i \cong \mbb{I}^{\spec(A \otimes \mbb{I}/i)} \cong \mbb{I}^{\spec A \times \spec\mbb{I}/i} \cong A^{\spec \mbb{I}/i} \cong A^{\ist{i}}\text{.} \]
  This way, it follows that 
  \[ L A \cong \sum_{i \in \mbb I}A^{\ist{i}} \cong \sum_{i \in \mbb I}A/i. \qedhere \]
\end{proof}

In other words, an element $(i,a) \in L A$ is simply a pair $i \in \mbb I$ and $a \in A$ that $a \le i$. This way, we can easily compute $L \mbb{I}$:

\begin{corollary}[\AxiomSQCI]
  $L \mbb{I} \cong \Delta^2$.
\end{corollary}
\begin{proof}
  By \zcref{prop:liftingofalgebra} and the assumption (\AxiomSQCI),
  \[ L\mbb{I} \cong \sum_{i \in \mbb I}\mbb{I}/i \cong \scomp{i,j \in \mbb{I}}{i \ge j} \cong \Delta^2\text{.} \]
  The second equivalence is again due to $\T$ being propositionally stable.
\end{proof}

\begin{remark}[Algebra vs.\ geometry]\label{rem:alggeoI}
  One interesting fact to note here is that, though $L\mbb{I}$ by computation is a dependent sum of algebras, it is naturally equivalent to a \emph{spectrum}. In some sense the source is the assumption that $\T$ is propositionally stable, which allows us to identify the algebraic object $\mbb{I}/i$ as a subset $\scomp{j \in \mbb I}{j \le i}$. There will be more examples of such nature once we work more specifically with distributive lattices; cf.\ \zcref{prop:simplicesasalgebra}.
\end{remark}

More generally, for domain theoretic applications we would want to compute the partial map classifier of $\Delta^2$, or $\Delta^n$ for any $n\in\N$. For this purpose, we observe that we can more generally compute them for any spectra (not only sober spaces):

\begin{proposition}[\AxiomSQCI]\label{prop:liftofsober}
  For an $\mbb{I}$-algebra $A$, we have
  \[ L\spec A \cong \sum_{i \in \mbb I}\mbb{I}\alg(A,\mbb{I}/i)\text{.} \]
\end{proposition}
\begin{proof}
  By (\AxiomSQCI) and \zcref{lem:openpropsober}, $\ist{i}$ is sober for any $i \in \mbb I$. By the duality result in \zcref{prop:duality}, we have
  \[ L\spec A \cong \sum_{i \in \mbb I}(\spec A)^{\ist{i}}\cong \sum_{i \in \mbb I}\mbb{I}\alg(A,\mbb{I}/i). \qedhere \]
\end{proof}

\begin{corollary}[\AxiomSQCI]
  For any $n\in\N$, we have
  \[ L\Delta^n \cong \Delta^{n+1}\text{,} \]
  and the unit $\eta_L \colon \Delta^n \hook L\Delta^n \cong \Delta^{n+1}$ takes $i_1 \ge \cdots \ge i_n$ in $\Delta^n$ to $1 \ge i_1 \ge \cdots \ge i_n$ in $\Delta^{n+1}$. 
\end{corollary}
\begin{proof}
  By \zcref{prop:liftofsober}, for the spectrum $\Delta^n \cong \spec\mbb{I}[n]/\ms{i}_1 \ge \cdots \ge \ms{i}_n$,
  \begin{align*}
    & L\Delta^n \\
    &\quad\cong \sum_{i \in \mbb I}\mbb{I}\alg(\mbb{I}[n]/\ms{i}_1\ge\cdots\ge \ms{i}_n,\mbb{I}/i) \\
    &\quad\cong \scomp{i_1,\cdots,i_n\in\mbb{I}/i}{i_1 \ge \cdots \ge i_n} \\
    &\quad\cong \scomp{i,i_1,\cdots,i_n\in\mbb{I}}{i \ge i_1 \ge \cdots \ge i_n} \\
    &\quad\cong \Delta^{n+1}
  \end{align*}
  Again the third steps holds since $\T$ is propositionally stable.
\end{proof}

\section{Distributive lattices and locality}\label{sec:locality}

One important finitary axiom for synthetic domain theory is \emph{Phoa's principle}, which we have mentioned at the beginning of this chapter is a consequence of \gls{QC} for distributive lattices. Thus, we now work with the theory $\mbb D$ of distributive lattices, or more generally the theory of $D$-algebra for some distributive lattice $D$. As indicated in \zcref{rem:sigmaframe}, computation for finitely presented algebras will not change if we replace distributive lattices by $\sigma$-frames. Hence, the results here and in \zcref{sec:order-theoretic-structure} will be applicable to $\sigma$-frames as well, where only the properties of finitely presented $\mbb{I}$-algebras matter.

As mentioned in \zcref{rem:opprop}, the theory of distributive lattices is also \emph{propositionally op-stable}, and similarly for $\sigma$-frames. Hence, dual versions of the previous results will also hold by symmetry, when we exchange $1$ for $0$ and $\wedge$ for $\vee$. Thus, we first record some simple corollaries for the dual structure.

\subsection{The dual dominance and co-partial map classifier}

In this case, $\mbb{I}$ has a minimal element $0$. Notice that the constant $0$ also induces a predicate on $\mbb{I}$,
\[ \isfsym\colon \mbb{I} \to \pp\text{,} \]
which takes any $i \in \mbb I$ to $i = 0$. Equivalently, we can view it as the $\ms t$ map for $\mbb{I}\op$.

\begin{corollary}
  If $\mbb{I}\op$ is stably spatial, then $\isfsym\colon \mbb{I} \to \pp$ is an embedding.
\end{corollary}

We will now call propositions in the image of $\isfsym$ \emph{closed}, and a subset classified by $\isfsym$ a \emph{closed subset}. Analogously to the previous \zcref{prop:Idominance}, closed propositions also form a dominance if $\mbb{I}\op$ is stably spatial:

\begin{corollary}\label{cor:dualisdominance}
  If $\mbb{I}\op$ is stably spatial, then $\isfsym$ forms a dominance.
\end{corollary}

There is an accompanying ``co-partial map classifier'' $T$ for partial maps with \emph{closed} domains, as \citet{hyland_first_2006} pointed out. Concretely, for any set $X$,
\[ T X \coloneq \sum_{i \in \mbb I} X^{\isf{i}}\text{.} \]
And analogously to \zcref{prop:liftingofalgebra}, under the assumption that $\mbb{I}\op$ is stably spatial, we can explicitly compute 
\[ T\Delta^n \cong \Delta^{n+1}\text{,} \]
where now the unit $\eta_T \colon \Delta^n \hook T\Delta^n \cong \Delta^{n+1}$ takes a sequence $i_1 \ge \cdots \ge i_n$ to $i_1 \ge \cdots \ge i_n \ge 0$ in $\Delta^{n+1}$. 

To maintain this symmetry, we also introduce the following axiom:

\PrintAxiomSQCID

The remaining part of this section will introduce various locality axioms for the interval $\mbb{I}$, and discuss some of their consequences with the \gls{QC} axiom.

\subsection{Non-triviality}

As a start, a minimal requirement for synthetic domain theory to model divergent computation is for the dominance $\mbb{I}$ to be closed under falsum. For this, we introduce the following minimal amount of non-triviality:

\PrintAxiomNT

Geometrically, the proposition $0 = 1$ can be identified with the spectrum $\spec\mbb{I}/0=1$. Hence, (\AxiomNT) states that the unique map $\emp \hook \spec\mbb{I}/0=1$ is an equivalence.
 %
This way, $\emp \eq \ist{0}$ now will be sober, and it is both an open and closed proposition. As mentioned in introduction, (\AxiomNT) together with a strong enough \gls{QC} principle will imply all of Hyland's axioms for synthetic domain theory. Before we discuss that in \zcref{sec:infdomain}, a minimal amount of \gls{QC} provided by (\AxiomSQCI) already implies a lot of elementary properties. As a first consequence of (\AxiomNT) with (\AxiomSQCI), we can show that $\mbb{I}$ is almost the Boolean algebra $2$ in the following sense:

\begin{proposition}[\AxiomNT, \AxiomSQCID]\label{prop:filed}
  For any $i \in \mbb I$ we have:
  \[ \neg \ist{i} \eq \isf{i}, \quad \neg\isf{i} \eq \ist{i}\text{.} \]
  In particular, the embedding $2 \hook \mbb{I}$ induced by $0,1$ is $\neg\neg$-dense:
  \[ \fa i{\mbb{I}} \dneg(\ist{i} \vee \isf{i})\text{.} \]
\end{proposition}
\begin{proof}
  If $\isf{i}$ then $\neg\ist{i}$ by (\AxiomNT). On the other hand, by conservativity in \zcref{lem:intconserve} and (\AxiomNT), if $i \neq 1$ then $i = 0$ since $0 \neq 1$. The dual case also follows by symmetry. Now given $i\in\mbb{I}$, if $\neg(\ist{i} \vee \isf{i})$, which implies $\neg\ist{i} \vee \neg\isf{i}$, equivalently $\isf{i} \wedge \ist{i}$, contradictory to (\AxiomNT). Hence, $\neg\neg(\ist{i} \wedge \isf{i})$.
\end{proof}

This allows us to observe that the open and closed propositions have a bijective correspondence between them:

\begin{corollary}[\AxiomNT, \AxiomSQCID]\label{cor:opendnegclose}
  For any proposition $p$, $p$ is open iff $\neg p$ is closed and vice versa. Furthermore, open and closed propositions are classical (cf.~\zcref{def:propertiesofpropositions}).
\end{corollary}

On the other hand, \zcref{sec:order-theoretic-structure} will show that under the axiom (\AxiomNT), \gls{QC} implies the interval $\mbb{I}$ will \emph{never} be isomorphic to $2$. 

\subsection{Locality}

A slightly stronger axiom which is common in the practice of domain theory is that the dominance $\mbb{I}$ should furthermore be closed under all finite disjunctions. For this, we consider the following axiom:

\PrintAxiomL

Geometrically, the second condition states that the following two inclusions are jointly surjective. 
\[ 
\begin{tikzcd}
  \mbb{I} \ar[r, hook, "{i \mapsto (i,1)}"] & \scomp{i,j \in \mbb{I}}{\ist{i\vee j}} & \mbb{I} \ar[l, hook', "{j \mapsto (1,j)}"']
\end{tikzcd}
\]

Under (\AxiomL), the dominance $\mbb{I}$ will be closed under finite disjunctions. Furthermore, it also implies more elementary properties about $\mbb{I}$:

\begin{lemma}[\AxiomL, \AxiomSQCID]\label{lem:intisnotBoolean}
  $\mbb{I}$ has no non-trivial complemented elements, i.e.\ if $i \in \mbb I$ is complemented, then $\ist{i} \vee \isf{i}$. 
\end{lemma}
\begin{proof}
  Suppose $i$ has a complement $j$. Then by (\AxiomL) we have
  \[ 1 \eq \ist{i\vee j} \eq \ist{i} \vee \ist{j}\text{,} \]
  and similarly,
  \[ \emp \eq \ist{i\wedge j} \eq \ist{i} \wedge \ist{j}\text{.} \]
  It follows that $\ist{j} \eq \neg \ist{i}$, and by \zcref{prop:filed}, $\ist{j} \eq \isf{i}$. Hence, $\ist{i} \vee \isf{i}$.
\end{proof}

\zcref{lem:intisnotBoolean} above allows us to realise 2 as a \emph{spectrum}:

\begin{example}\label{exm:2issober}
  Consider the algebra $B \coloneq \mbb{I}[\ms{i},\ms{j}]/{\ms{i}\wedge \ms{j} =0,\ms{i}\vee \ms{j} = 1}$. By construction, $B$ classifies complemented elements in $\mbb{I}$-algebras. Thus, assuming (\AxiomL) and (\AxiomSQCI), by \zcref{lem:intisnotBoolean}
  \[ \spec B \cong \scomp{i,j \in \mbb{I}}{i \wedge j = 0 \wedge i \vee j = 1} \cong 2\text{.} \]
  In particular, this means that for any sober space $X \cong \spec A$, $2^X$ will be equivalent to the set of complemented elements in $A$ by \zcref{prop:duality}.
\end{example}

As another example of a new spectrum (\AxiomL) allows us to define:

\begin{example}\label{exm:hornsober}
  Consider the right outer horn $\Lambda^2_2$, which we may describe by the following pushout:
  \[
    \begin{tikzcd}
      1 \ar[d, "1"'] \ar[r, "1"] & \mbb{I} \ar[d] \\
      \mbb{I} \ar[r] & \Lambda^2_2
      \arrow["\lrcorner"{anchor=center, pos=0.125, rotate=180}, draw=none, from=2-2, to=1-1]
    \end{tikzcd}
  \]
  By viewing $\Lambda^2_2$ as a subspace of $\mbb{I}^2$, we may identify it as follows:
  \[ \Lambda^2_2 \cong \scomp{(i,j) : \mbb{I}^2}{\ist{i} \vee \ist{j}}\text{.} \]
  Assuming (\AxiomL), $\ist{i} \vee \ist{j} \eq \ist{i\vee j}$, it follows that 
  \[ \Lambda^2_2 \cong \spec \mbb{I}[\ms{i},\ms{j}]/\ms{i} \vee \ms{j}\text{.} \]
\end{example}

Axiom (\AxiomL) clearly has a dual counterpart:


\PrintAxiomCL

Similarly to \zcref{exm:hornsober}, this allows one to identify the left outer horn $\Lambda^2_0$ as a spectrum. Of course we can combine the axioms (\AxiomL) and (\AxiomCL). 

\subsection{Linearity and coskeletality}\label{sec:linearity-and-coskeletality}

An axiom even stronger than both (\AxiomL) and (\AxiomCL) is the simplicial axiom (\AxiomSL):

\PrintAxiomSL

Geometrically, the simplicial axiom states that the two simplices $\Delta^2,\Delta_2$ covers the square $\Delta^2 \cup \Delta_2 \cong \mbb{I}^2$.
The strongest locality axiom states that the simplicial structure is truncated to level 1:


\PrintAxiomOneCS

Geometrically, (\AxiomOneCS) additionally requires the canonical inclusion $\partial\Delta_2 \hook \Delta_2$ from the boundary $\partial\Delta_2$ to the 2-simplex $\Delta_2$ is an equivalence, where $\partial\Delta_2$ is simply
\[ \partial\Delta_2 \coloneq \scomp{i,j \in \mbb{I}}{\isf{i} \vee (i = j) \vee \ist{i}}\text{.} \]

In general we will at most work with the weakest locality axiom (\AxiomNT). However, we will show in \zcref{sec:local} that $\mbb{I}$, hence all spatial algebras and spectra, will be \emph{right orthogonal} to the maps that define these local principles as shown above. We will see in \zcref{sec:newmodelforsdt} that these axioms corresponds to various coverages we can put on the presheaf classifying topos, and the orthogonality result particularly implies that $\mbb{I}$ will belong to these subtopoi.


\section{Order-theoretic structure on algebras and spaces}\label{sec:order-theoretic-structure}

In this section, we describe several important local orderings on algebras and spaces.


\subsection{Local ordering of algebras}

Every $\mbb{I}$-algebras has a built-in ``canonical'' partial ordering.

\begin{definition}[Canonical partial order]
  We shall write $\le_A$ for the \emph{canonical} partial ordering of an $\mbb{I}$-algebra $A$ defined equivalently by meets or joins:
  \[ 
    a\le_A b \eq (a\wedge b = a) \eq (a\vee b = b)
  \]
\end{definition}

\begin{example}
  The canonical partial ordering of an algebra of observations $\opens{X}$ is given pointwise because $\opens{X}$ is the product of algebras $\prod_{x\in X}\mbb{I}$:
  \begin{align*}
    U\le_{\opens{X}} V 
    &\eq \fa{x}{X} U(x)\le_\mbb{I} V(x)
  \end{align*}
\end{example}

In addition to the canonical partial order, there is a local preordering of algebras that mirrors the observational order for spaces.

\begin{definition}[Behavioural preorder]
  For any $\mbb{I}$-algebra $A$, we define the \emph{behavioural preorder} on $A$ as follows:
  \begin{align*}
    a \behle[A] b &\eq \fa{x}{\spec{A}} x(a) \le_{\mbb{I}} x(b)
  \end{align*}
\end{definition}


To relate the behavioural preorder to the canonical partial ordering, we first define a slight weakening of the \gls{QC} principle.

\begin{definition}
  We say that an $\mbb{I}$-algebra $A$ is \emph{pre-spatial} when the counit component $\iota_A\colon A\to \opens\spec{A}$ is an embedding.
\end{definition}

\begin{proposition}\label{lem:canon-beh-coincide}
  The canonical partial order of any pre-spatial $\mbb{I}$-algebra $A$ is its behavioural preorder.
\end{proposition}

\begin{proof}
  For $a,b\in A$ we have the following sequence of equivalences:
  \begin{align*}
    & a\le_A b \\
    &\quad\eq \iota_A(a) \le_{\opens\spec{A}} \iota_A(a)
    \\ 
    &\quad\eq \fa{x}{\spec{A}} x(a) \le_{\mbb{I}} x(b)
    \\ 
    &\quad\eq a\behle[A]b
  \end{align*}
  The first equivalence requires that $A$ is pre-spatial.
\end{proof}


\subsection{Local ordering of spaces}

The \emph{observational preorder} is a classical notion in synthetic domain theory (cf.\ \citet{PhoaWesleyKym-Son1991DtiR,hyland_first_2006}):

\begin{definition}[Observational preorder]\label{defn:specialisation}
  For any set $X$, we define the \emph{observational preorder} on $X$ as follows:
  \begin{align*}
    x\obsle[X] y 
    &\eq \fa{U}{\opens{X}} U(x)\le_\mbb{I} U(x)
  \end{align*}
\end{definition}

By definition, the observational preorder is reflexive and transitive. As already observed in \emph{loc.\ cit.}, one important property of the observational preorder is that \emph{every} map is monotone w.r.t.\ this order:

\begin{lemma}\label{lem:anymapmonotoneintriscorder}
  For $f \colon X \to Y$, $x \obsle y$ in $X$ implies $fx \obsle fy$ in $Y$.
\end{lemma}
\begin{proof}
  This simply follows from compositionality of functions.
\end{proof}

We can very easily classify the observational preorder on a decidable set.

\begin{lemma}[\AxiomNT]\label{lem:obs-preord-on-set-with-decidable-equality}
  If $M$ is decidable, then the observational preorder on $M$ is discrete, in the sense that for $n,m \in M$,
  \[ n \obsle[M] m \to n = m\text{.} \]
\end{lemma}
\begin{proof}
  Since $M$ has decidable equality, we can define functions by case distinction. If $n \neq m$, we can construct a function $U \colon M \to \mbb{I}$ with $U(n) = 1,U(m) = 0$, contradicting with $n \obsle[M] m$. Hence, $\neg\neg (n=m)$, and with decidable equality this implies $n = m$.
\end{proof}


\begin{corollary}[\AxiomNT]\label{cor:connectedpreservediscrete}
  If the observational preorder on $X$ has a least or greatest element, then $X$ is \emph{internally connected} in the sense that $X$ is right orthogonal to $M \to 1$ for any set $M$ with decidable equality.
\end{corollary}
\begin{proof}
  Since $X$ is inhabited, the restriction map $M \to M^X$ is an embedding. It suffices to show that this map is also surjective. Fixing $f\colon X\to M$, we must show that there exists $m:M$ such that $f(x) = m$ for all $x\in X$. Let $x_0$ be the least element of $X$ in its observational preorder, so that we have $x_0\obsle[X] x$; by \zcref{lem:anymapmonotoneintriscorder} we have $f(x_0) \obsle f(x)$ which implies $f(x_0) = f(x)$ by \zcref{lem:obs-preord-on-set-with-decidable-equality}. Therefore, we may choose $m\coloneq f(x_0)$. The case for a maximal element is analogous.
\end{proof}

\begin{remark}
  With greater efforts, the above proof would also work with the weaker assumption that the observational preorder on $X$ is \emph{connected}. This generalises a similar result given by \citet[Prop.~4.4.1]{hyland_first_2006}.
\end{remark}


\subsection{The satisfaction order of a spectrum}

\zcref{cor:connectedpreservediscrete} demonstrated the usefulness of the observational preorder. We would like to use it to show e.g.\ that $\mbb{I}$ is internally connected. For this, we need to characterise the observational preorder on sober spaces. It turns out that for any sober space, this is simply its \emph{satisfaction order} in the sense defined below:


\begin{definition}[Satisfaction order]
  We shall write $\satle[\spec{A}]$ for the \emph{satisfaction order} on the spectrum $\spec{A}$ of an $\mbb{I}$-algebra as defined below:
  \begin{align*}
    x\satle[\spec{A}] y 
    &\eq \fa{a}{A}
    x(a) \le_\mbb{I} y(a) 
  \end{align*}
\end{definition}


\begin{observation}\label{lem:specorderofsober}
  When $A$ is a spatial $\mbb{I}$-algebra (and thus $\spec{A}$ is sober), the observational and satisfaction orders on $\spec{A}$ coincide.
\end{observation}
\begin{proof}
  For any $x,y:\spec{A}$ we have the following chain of equivalences:
  \begin{align*}
    &x\obsle[\spec{A}] y \\
    &\quad\eq 
    \fa{U}{\opens\spec{A}} U(x) \le_{\mbb{I}} U(y)
    \\ 
    &\quad\eq 
    \fa{a}{A} 
    \iota_A(a)(x)
    \le_\mbb{I}
    \iota_A(a)(y)
    \\ 
    &\quad\eq 
    \fa{a}{A} 
    x(a)
    \le_\mbb{I}
    y(a)
    \\ 
    &\quad\eq 
    x\satle[\spec{A}] y
    \qedhere
  \end{align*}
\end{proof}

\begin{observation}
  The satisfaction order on the spectrum $\spec{\mbb{I}[X]}$ of a free $\mbb{I}$-algebra is induced by the canonical order on the \emph{observation algebra} $\opens{X}$ via the canonical bijection $\spec{\mbb{I}[X]}\cong \opens{X}$.
\end{observation}

In the above, $\opens{X}$ is playing two roles simultaneously: it is both a spectrum of $\mbb{I}[X]$ and the observational algebra of $X$.

\begin{proof}
  For any $u,v\in\spec{\mbb{I}[X]}$ we have the following equivalences:
  \begin{align*}
    &u \satle[\spec{\mbb{I}[X]}] v\\
    &\quad\eq \fa{p}{\mbb{I}[X]} u(p) \le_{\mbb{I}} v(p)
    \\ 
    &\quad\eq \fa{x}{X} u(x) \le_{\mbb{I}} v(x)
    \\ 
    &\quad\eq (\ld{x}{X}u(x)) \le_{\opens{X}} (\ld{x}{X}v(x))
    \qedhere
  \end{align*}
\end{proof}


\begin{lemma}\label{lem:cancoincide}
  For any surjective homomorphism $q \colon \mbb{I}[Y] \surj A$ of a free algebra $\mbb{I}[Y]$, the satisfaction order on $\spec A$ is induced by the satisfaction order on $\spec{\mbb{I}[Y]}$ via the inclusion ${\spec{q}}\colon\spec{A}\hook \spec{\mbb{I}[Y]}$.
\end{lemma} 

\begin{proof}
  Let $x,x' \in \spec A$. We have the following equivalences:
  \begin{align*}
    &(\spec{q})(x)\satle[\spec{\mbb{I}[Y]}] (\spec{q})(x') \\
    &\quad\eq 
    \fa{p}{\mbb{I}[Y]}
    x(qp) \le_\mbb{I} x'(qp)
    \\ 
    &\quad\eq 
    \fa{a}{A} x(a) \le_\mbb{I} x'(a) 
    \\ 
    &\quad\eq 
    x\satle[\spec{A}] x'
  \end{align*}
% 
  The second equivalence holds because $q\colon \mbb{I}[Y]\surj A$ is surjective.
\end{proof}


\section{Phoa's principle, finitary QC, and homotopy}\label{sec:phoaandhomotopy}

The interval $\mbb{I}$ is playing multiple roles so far:
\begin{enumerate}
  \item Mapping a set $X$ into $\mbb{I}$ computes its algebra of observations $\opens{X} = \mbb{I}^X$.
  \item Mapping an $\mbb I$-algebra $A$ into $\mbb{I}$ computes its spectrum $\spec{A} = \mbb{I}\alg(A,\mbb{I})$.
\end{enumerate}

In both cases above, $\mbb{I}$ is viewed as an algebra of observations and not as a geometrical figure. For the latter, we may consider mappings from $\mbb{I}$ into \emph{either} an algebra or a space as defining a notion of (directed) homotopy. In particular, the function space $X^\mbb{I}$ classifies \emph{paths} drawn in $X$; this geometrical role of $\mbb{I}$ as a figure shape for paths is the primary one in applications to synthetic category theory~\citep{riehl2017type}, and it is also upon the geometry of the interval that the notion of \emph{chain completeness} in synthetic domain theory is founded (cf.\ \zcref{sec:infdomain}).

\subsection{The generalised Phoa principle and spatiality of polynomial algebras}

As soon as we have path spaces $X^\mbb{I}$, we immediately wish to begin characterising them for specific $X$. For example, the path space of a function space $A\to B$ is given pointwise in terms of the path space of $B$, which follows immediately from the laws of exponentials. To characterise paths in spaces like $\mbb{I}^n$ and $\Delta^n$, however, we need additional axioms. The simplest such axiom asserts that the polynomial algebra $\mbb{I}[\ms{i}]$ is spatial---which, in the case of distributive lattices, is equivalent to Phoa's principle.

\begin{definition}\label{def:gen-phoa}
  We say that an $\mbb{I}$-algebra $A$ satisfies the \emph{generalised Phoa principle} when either of the following equivalent conditions hold:
  \begin{enumerate}
    \item The path space $A^\mbb{I}$ classifies the canonical order on $A$.
    \item For any function $\alpha\colon \mbb{I}\to A$, we have $\alpha(i) = \alpha(0)\vee i \wedge \alpha(1)$.
  \end{enumerate}
\end{definition}

\begin{proof}
  Suppose that $A^\mbb{I}$ classifies the canonical order on $A$; we must show that for any $\alpha\colon \mbb{I}\to A$ we have $\alpha(i) = \alpha(0)\vee i \wedge \alpha(1)$. By our assumption that $A^\mbb{I}$ classifies the canonical order on $A$, we know that that any such function must be uniquely determined by its values on $0$ and $1$ and moreover that $\alpha(0) \le_A \alpha(1)$; therefore, it suffices to observe for arbitrary such $\alpha$ that
  \begin{align*}
    \alpha(0) &= \alpha(0) \vee 0 = \alpha(0)\vee 0 \wedge \alpha(1)\text{,}\\
    \alpha(1) &= \alpha(0)\vee\alpha(1) = \alpha(0)\vee 1 \wedge \alpha(1)\text{.}
  \end{align*}

  Conversely, suppose that every path $\alpha\colon \mbb{I}\to A$ is of the form $\alpha(i) = \alpha(0)\vee i \wedge \alpha(1)$. We must show that $a\le_A b$ holds if and only if there exists a unique path $\alpha\colon \mbb{I}\to A$ sending $0$ and $1$ to $a$ and $b$ respectively.
  \begin{enumerate}
    \item If $a\le_A b$ holds, we define $\alpha(i) \coloneq a\vee i \wedge b$; this path is unique with the required property by our assumption. 
    \item If there exists a path $\alpha\colon\mbb{I}\to A$ from $a$ to $b$, we know that it has the form $\alpha(i) = a \vee i \wedge b$. Therefore we have
    \[ 
      b = \alpha(1) = a\vee 1\wedge b = a \vee b.
      \qedhere
    \] 
  \end{enumerate}
\end{proof}

The standard Phoa principle is the generalised Phoa principle for $\mbb{I}$ itself.

\begin{theorem}\label{thm:gen-phoa-poly}
  If $A$ satisfies the generalised Phoa principle, then so does the polynomial $\mbb{I}$-algebra $A[\ms{j}]$.
\end{theorem}

\begin{proof}
  Fixing $\alpha\colon \mbb{I}\to A[\ms{j}]$, we must check that $\alpha(i) = \alpha(0) \vee i \wedge \alpha(1)$. 
  %
  The normal form theorem for polynomials of distributive lattices~\citep[Ch.\ 1, Thm.\ 10.11]{lausch2000algebra} implies that any function $\alpha\colon \mbb{I}\to A[\ms{j}]$ must be determined by functions $\alpha_0,\alpha_1\colon \mbb{I}\to A$ such that in $A[\ms{j}]$ we have $\alpha(i) = \alpha_0(i)\vee \ms{j} \wedge \alpha_1(i)$. Meanwhile, Phoa's principle for $A$ characterises $\alpha_0$ and $\alpha_1$:
  \begin{align*}
    \alpha_0(i) &= \alpha_0(0) \vee i \wedge \alpha_0(1)
    \\ 
    \alpha_1(i) &= \alpha_1(0) \vee i \wedge \alpha_1(1)
  \end{align*}

  Hence we have $\alpha(i) = (\alpha_0(0) \vee i \wedge \alpha_0(1))\lor x \wedge (\alpha_0(0) \vee i \wedge \alpha_0(1))$. On the other hand, we can use the normal form to compute $\alpha(0)$ and $\alpha(1)$ as polynomials in $\ms{j}$: 
  \begin{align*}
    \alpha(0) &= \alpha_0(0) \vee \ms{j}\wedge \alpha_1(0)\\
    \alpha(1) &= \alpha_0(1) \vee \ms{j}\wedge \alpha_1(1)
  \end{align*}

  Hence $\alpha(0)\lor i \wedge \alpha(1) = (\alpha_0(0) \vee \ms{j}\wedge \alpha_1(0))\vee i \wedge (\alpha_0(1) \vee \ms{j}\wedge \alpha_1(1))$. Using elementary lattice algebra we deduce that $\alpha(i) = \alpha(0) \vee i \wedge \alpha(1)$.
\end{proof}

\begin{theorem}\label{obs:sqcp-phoa-gen}
  Let $A$ be a spatial $\mbb{I}$-algebra. Then the following are equivalent:
  \begin{enumerate}
    \item The polynomial $\mbb{I}$-algebra $A[\ms{j}]$ is spatial.
    \item The generalised Phoa principle holds for $A$.
  \end{enumerate}
\end{theorem}

\begin{proof}
  We note that $\spec{A[\ms{j}]} \cong \spec{A}\times \mbb{I}$ unconditionally, and so the algebra of observations $\opens\spec{A[\ms{j}]}$ is actually the path space of $\opens\spec{A}$; taking spatiality of $A$ into account, the counit component $A[\ms{j}]\to \opens\spec A[\ms{j}]$ is the evaluation map $A[\ms{j}] \to A^\mbb{I}$, and we wish to prove that the latter is an equivalence if and only if for each $\alpha\colon \mbb{I}\to A$, $\alpha(i)=\alpha(0)\vee i\wedge\alpha(1)$. 
  %
  This follows immediately from the normal form for polynomials of distributive lattices~\citep[Ch.\ 1, Thm.\ 10.11]{lausch2000algebra}: any polynomial $p\in A[\ms{j}]$ is uniquely of the form $p = \ev_0(p) \vee \ms{j} \wedge \ev_1(p)$ with $\ev_0(p)\le_A \ev_1(p)$.
\end{proof}


\begin{corollary}\label{cor:A-and-Ax-qc-to-Axy-qc}
  If both $A$ and $A[\ms{j}]$ are spatial, then so is $A[\ms{j},\ms{k}]$.
\end{corollary}
\begin{proof}
  By \zcref{obs:sqcp-phoa-gen}, the generalised Phoa principle holds for $A$ because $A$ and $A[\ms{j}]$ are spatial; by \zcref{thm:gen-phoa-poly}, the generalised Phoa principle holds for $A[\ms{j}]$ because it holds for $A$.  By \zcref{thm:gen-phoa-poly} again, the generalised Phoa principle holds for $A[\ms{j},\ms{k}] = A[\ms{j}][\ms{k}]$ because it holds for $A[\ms{j}]$.
\end{proof}

\begin{corollary}\label{cor:phoa-vs-quasicoherence}
  The following are equivalent:
  \begin{enumerate}
    \item The polynomial $\mbb{I}$-algebra $\mbb{I}[\ms{i}]$ is spatial.
    \item Phoa's principle holds.
  \end{enumerate}
\end{corollary}

\begin{proof}
  By \zcref{obs:sqcp-phoa-gen}, because $\mbb{I}$ is unconditionally spatial.
\end{proof}


We therefore identify the axiom (\AxiomSQCP) below accordingly, which we have seen to be equivalent to Phoa's principle.

\PrintAxiomSQCP


The \gls{QC} axiom for polynomials in one variable is surprisingly strong:

\begin{lemma}[\AxiomSQCP]\label{lem:fg-qc}
  Every finitely generated free $\mbb{I}$-algebra is spatial.
\end{lemma}

\begin{proof}
  The free $\mbb{I}$-algebra on $0$ generators is just $\mbb{I}$, which is automatically spatial. The free $\mbb{I}$-algebra $\mbb{I}[\ms{i}]$ on one generator is spatial by assumption. In the case of $\mbb{I}[\ms{i},\ms{j}_0,\ldots,\ms{j}_n]$ we recall \zcref{cor:A-and-Ax-qc-to-Axy-qc} and see that $\mbb{I}[\ms{i}]$ spatial implies $\mbb{I}[\ms{i},\ms{j}_0]$ spatial, which implies $\mbb{I}[\ms{i},\ms{j}_0,\ms{j}_1]$ spatial and so on.
\end{proof}


\begin{theorem}[\AxiomSQCP]\label{thm:phoa-spectrum}
  For any $\mbb{I}$-algebra $A$, the path space ${(\spec A)}^\mbb{I}$ classifies the satisfaction order on $\spec A$.
\end{theorem}
\begin{proof}
  By assumption $\mbb{I}[\ms{i}]$ is spatial, so from \zcref{prop:duality} may characterise the path space of $\spec{A}$ algebraically as follows:
  \[ (\spec A)^\mbb{I} \cong (\spec{A})^{\spec{\mbb{I}[\ms{i}]}}\cong \mbb{I}\alg(A,\mbb{I}[\ms{i}])\text{.} \]
  
  By the normalisation theorem for $\mbb{I}[\ms{i}]$, the pair of evaluation maps
  \[ \pair{\ev_0,\ev_1} \colon \mbb{I}[\ms{i}] \to \mbb{I} \times \mbb{I} \]
  classifies the canonical order on $\mbb{I}$.
  Thus every homomorphism $h\colon A\to \mbb{I}[\ms{i}]$ is canonically induced by elements $h_0,h_1\in\spec{A}$ with $h_0\satle[\spec{A}] h_1$ such that $h(a) = h_0(a) \vee \ms{i} \wedge h_1(a)$.
\end{proof}


\begin{theorem}[\AxiomSQCP]\label{thm:algebraphoa}
  If both $A$ and $A[\ms{i}]$ are spatial $\mbb{I}$-algebras, then the path space $A^\mbb{I}$ classifies the canonical order on $A$.
\end{theorem}
\begin{proof}
  By assumption and \zcref{specrightadj}, since $A \otimes \mbb{I}[\ms{i}] \cong A[\ms{i}]$, we have
  \[ A^\mbb{I} \cong (\opens\spec A)^\mbb{I} = (\mbb{I}^{\spec A})^\mbb{I} \cong \mbb{I}^{\spec A \times \spec\mbb{I}[\ms{i}]} \cong \mbb{I}^{\spec A[\ms{i}]} \cong A[\ms{i}]\text{.} \]
  We have already seen that $A[\ms{i}]$ classifies the canonical order of $A$ by the normalisation result for polynomial $\mbb{I}$-algebras.
\end{proof}

\begin{remark}[Algebraic properties in classifying topoi]\label{rem:normalalgebra}
  We emphasise that the above proofs are purely the consequences of the algebraic fact that $A[\ms{i}]$ classifies the canonical order of any $\mbb{I}$-algebra $A$, plus \gls{QC}. This is a perfect example of how an algebraic property of a theory has a non-trivial effect on the internal logic of its classifying topos.
\end{remark}

\begin{corollary}[\AxiomSQCP]\label{lem:cubes-are-sober}
  Each cube $\mbb{I}^n$ is sober.
\end{corollary}
\begin{proof}
  $\mbb{I}^n$ is the spectrum of the free $\mbb{I}$-algebra $\mbb{I}[n]$, which is spatial by \zcref{lem:fg-qc}.
\end{proof}

\begin{corollary}[\AxiomNT, \AxiomSQCP]\label{cor:cubes-internally-connected}
  The $n$-cube $\mbb{I}^n$ is internally connected.
\end{corollary}

In particular, $\mbb{I}$ is not the Boolean set $2$.

\begin{proof}
   By \zcref{lem:specorderofsober,lem:cancoincide}, the observational preorder space $\mbb{I}^n$ is induced by the satisfaction order on $\spec\mbb{I}[n]$ because $\mbb{I}^n$ is sober (\zcref{lem:cubes-are-sober}); the satisfaction order is simply the pointwise order, so it has a top element. Thus by \zcref{cor:connectedpreservediscrete}, $\mbb{I}^n$ is internally connected. So are $\Delta^n$ because they are retracts of $\mbb{I}^n$.
\end{proof}

\subsection{Comparing the different orders}

An immediate application of Phoa principle is to compare path structure with the observational preorder.

\begin{observation}[\AxiomSQCP]
  For any set $X$, the boundary evaluation map 
  \[ \pair{\ev_0,\ev_1}\colon X^\mbb{I}\to X\times X \] 
  factors through the observational preorder relation ${\obsle[X]}\hook X\times X$ as displayed below:
  \[ 
    \begin{tikzcd}
      X^\mbb{I} 
        \arrow[rr,densely dashed, "\exists"]
        \arrow[dr]
      && 
      {\obsle[X]}
        \arrow[dl,hookrightarrow]
      \\ 
      &
      X\times X
    \end{tikzcd}
  \] 
\end{observation}
\begin{proof}
  Fix a path $\alpha\colon \mbb{I}\to X$. We have $0\obsle[\mbb{I}] 1$ by Phoa's principle; because every function is monotone in the observation preorder, we conclude $\alpha(0)\obsle[X] \alpha(1)$.
\end{proof}

\begin{definition}
  A set $X$ for which $X^\mbb{I}\to {\obsle[X]}$ is an surjection is called \emph{linked} (following \citet{PhoaWesleyKym-Son1991DtiR}); when the map is an equivalence, $X$ is called \emph{strongly linked}.
\end{definition}

\begin{lemma}[\AxiomSQCP]\label{lem:sober-strongly-linked}
  Any sober space is strongly linked.
\end{lemma}

\begin{proof}
  Let $X=\spec{A}$ be the spectrum of a spatial algebra $A$. 
  By \zcref{thm:phoa-spectrum}, the path space $X^\mbb{I}$ classifies the satisfaction order on $\spec{A}$, which is to say that there is a (necessarily unique) path from $x$ to $y$ if and only if \[\fa{a}{A} x(a)\le_\mbb{I} y(a)\text{.}\] Because $A$ is spatial, we have $A\cong \opens X$ and there is a unique path from $x$ to $y$ if and only if \[\fa{U}{\opens{X}} U(x)\le_\mbb{I} U(y)\text{,}\] but this is the observational order of $X$.
\end{proof}

\begin{lemma}[\AxiomSQCP]\label{lem:int-strong-link}
  The interval is strongly linked.
\end{lemma}

\begin{proof}
  The interval is sober because it is the spectrum of $\mbb{I}[\ms{i}]$, which is spatial by (\AxiomSQCP).
\end{proof}

\begin{lemma}\label{lem:strong-link-exp-ideal}
  Being strongly linked is an exponential ideal.
\end{lemma}

\begin{proof}
  \citet[Prop.\ 5.4.4]{PhoaWesleyKym-Son1991DtiR} implies that being linked is an exponential ideal, because an exponential is an internal limit. A set $X$ is strongly linked if and only if it is linked \emph{and} it is $\mbb{I}$-separated in the sense that $X^\mbb{I}\to X\times X$ is an embedding; but $\mbb{I}$-separated sets also form an exponential ideal (in fact, a reflective exponential ideal).
\end{proof}

\begin{lemma}[\AxiomSQCP]\label{lem:O-strongly-linked}
  Any algebra of observations is strongly linked.
\end{lemma}

\begin{proof}
  We assume that $A=\opens{X}$ for some not necessarily sober $X$; then $\opens{X}=\mbb{I}^X$ is strongly linked by \zcref{lem:strong-link-exp-ideal,lem:int-strong-link}.
\end{proof}


\begin{corollary}\label{cor:strong-link-obs-pw}
  When $Y$ is strongly linked, the observational preorder on $Y^X$ is induced pointwise by the observational preorder on $Y$.
\end{corollary}

\begin{proof}
  Let $Y$ be strongly linked; by \zcref{lem:strong-link-exp-ideal}, so is $Y^X$. Thus $\obsle[Y^X]$ is precisely given by $(Y^X)^\mbb{I}\cong (Y^\mbb{I})^X\cong (\obsle[Y])^X$.
\end{proof}

\begin{proposition}[\AxiomSQCP]\label{prop:specordopens}
  On any observation algebra $\opens{X}$, the observational preorder coincides with the canonical order.
\end{proposition}

\begin{proof}
  By Phoa's principle, $\mbb{I}$ is strongly linked; thus by \zcref{cor:strong-link-obs-pw}, the observational preorder on $\mbb{I}^X$ is pointwise induced by that of $\mbb{I}$. In particular, we have the following equivalences involving the function space $\opens{X}=\mbb{I}^X$:
  \begin{align*}
    &U\obsle[\mbb{I}^X]V \\
    &\quad\eq \fa{x}{X} U(x)\obsle[\mbb{I}] V(x)
    \\
    &\quad\eq \fa{x}{X} U(x)\le_{\mbb{I}} V(x)
    \\
    &\quad\eq U\le_{\opens{X}} V \qedhere
  \end{align*}
  The second equivalence uses Phoa's principle and \zcref{lem:O-strongly-linked} to identify the canonical order on $\mbb{I}$ with its observational preorder.
% 
\end{proof}

\begin{proposition}[\AxiomSQCP]\label{cor:specordonalgiscan}
  For a spatial $\mbb{I}$-algebra $A$, the observational preorder on the underlying set of $A$ is induced by the canonical order on $\opens\spec{A}$ by the counit isomorphism $\iota_A\colon A\to \opens\spec{A}$. Hence the observational order and canonical order of $A$ coincide.
\end{proposition}

\begin{proof}
  Because $A$ is spatial, the counit component $\iota_A$ gives us an order-isomorphism 
  \[ {(A,\le_A)} \cong (\opens\spec{A},\le_{\opens\spec{A}}). \] 
  By \zcref{prop:specordopens}, this is an order-isomorphism ${(A,\le_A)}\to {(\opens\spec{A},\obsle[\opens\spec{A}])}$. Meanwhile, any bijective function tracks an isomorphism of observational preorders, so the inverse $\iota_A^{-1}\colon \opens\spec{A}\to A$ tracks an order-isomorphism  ${(\opens\spec{A},\obsle[\opens\spec{A}])}\to {(A,\obsle[A])}$. 
\end{proof}

\begin{corollary}[\AxiomSQCP]
  The canonical partial order, observational preorder, and behavioural preorder of a spatial $\mbb{I}$-algebra all coincide.
\end{corollary}

\begin{proof}
  By \zcref{cor:specordonalgiscan,lem:canon-beh-coincide}.
\end{proof}


\subsection{Spatiality of finitely presented algebras}

The (\AxiomSQCP) axiom does not imply that the simplices $\Delta^n$ are sober; for this, we need a stronger \gls{QC} principle that applies to all finitely \emph{presented} algebras.


\PrintAxiomSQCF

\begin{remark}
  In particular, the axiom (\AxiomSQCF) implies all our previously mentioned \gls{QC} axioms, including (\AxiomSQCP) and (\AxiomSQCID) (hence also (\AxiomSQCI)).
\end{remark}

It is immediate that $\Delta^n\cong \spec\mbb{I}[\ms{i}_1,\cdots,\ms{i}_n]/\ms{i}_1\geq \cdots \geq \ms{i}_n$ is sober under (\AxiomSQCF).

\begin{corollary}[\AxiomSQCF]
  Each $n$-simplex $\Delta^n$ is internally connected.
\end{corollary}

\begin{proof}
  This follows in the same way as \zcref{cor:cubes-internally-connected}.
\end{proof}

\subsection{Classification of simplices by algebras}

At the end of this section, we describe another interesting perspective arising from the proof of \zcref{thm:phoa-spectrum}. We have seen that the dualising object $\mbb{I}$ has a double role: It is both an algebra and a spectrum. The mixture of algebraic and geometric object has also been observed in \zcref{rem:alggeoI}. The proof of \zcref{thm:phoa-spectrum} gives us many more such examples. For instance, $\mbb{I}[i]$ classifies the order on $\mbb{I}$, which by definition is the spectrum $\Delta_2$. In fact, \emph{all} the simplices are classified by some algebra.

Just as we described $\Delta^n$ in terms of descending sequences in $\mbb{I}$, we can do the same replacing $\mbb{I}$ with another algebra $A$.
For any $n\in\N$ and $\mbb{I}$-algebra $A$, we define $\Delta[A]^{n}$ to be the set of lists of decreasing elements of length $n$ in $A$, namely
\[ \Delta[A]^{n} \coloneq \scomp{a_1,\cdots,a_n \in A}{a_1 \ge \cdots \ge a_n}\text{.} \]
Of course $\Delta[\mbb{I}]^n$ is simply $\Delta^n$. We have a general algebraic description of $\Delta[A]^n$.

\begin{proposition}\label{prop:simplicesasalgebra}
  For any $n\in\N$ and $\mbb{I}$-algebra $A$, there is an equivalence 
  \[ \Delta[A]^{n+1} \cong A[n]/\ms{i}_1 \le \cdots \le \ms{i}_n\text{,} \]
  sending any $a_0 \ge \cdots \ge a_n$ in $A$ to the polynomial 
  \[ a_0 \wedge \ms{i}_1 \vee a_1 \wedge \cdots \wedge \ms{i}_n \vee a_n\text{.} \]
\end{proposition}

We omit the proof here, which again follows from a general normal form for polynomials with finitely many variables for distributive lattices; see \citet[Ch.\ 1, Thm.\ 10.21]{lausch2000algebra}. We only note here that the above polynomial is well-defined, in the sense that the expression is associative and its value does not depend on how one add parenthesises. The reason for this is that for any expression $a \vee b \wedge c$ to have a definite meaning in a distributive lattice, i.e.\ the two value $(a \vee b) \wedge c$ and $a \vee (b \wedge c)$ coincide, it suffices to have $a \le c$. Then the condition $a_0 \ge \cdots \ge a_n$ for the parameters, and the fact that $\ms{i}_1 \le \cdots \le \ms{i}_n$, makes sure the above polynomial has a definite meaning in the quotient.


\section{Chain completeness and infinitary domain theory}\label{sec:infdomain}

Until this point, we have seen that elementary axioms for synthetic domain theory follow from (SQCF) for distributive lattices and thus $\sigma$-frames. In this section we show that the final axiom of synthetic domain theory, viz.\ chain completeness of the interval $\mbb{I}$, is also a consequence of \gls{QC} for $\sigma$-frames. In fact, it implies the main infinitary axiom of synthetic domain theory in~\citet{fiore-plotkin:1996}---the initial algebra $\omega$ for $L$ is inductive.

\subsection{Chain completeness and inductivity of the initial algebra}

Chain completeness is again specified as an orthogonality condition. Let $\omega,\ov\omega$ denote the initial algebra and final coalgebra for the partial map classifier $L$, respectively. Intuitively, $\omega$ behaves as an infinite chain, and $\ov\omega$ as an infinite chain with a top element. The chain completeness condition is thus intuitively stating a set has joins for infinite sequences internally:

\begin{definition}[Chain completeness]
  A set $X$ is \emph{chain complete} if it is right orthogonal to the inclusion $\omega \hook \ov\omega$.
\end{definition}

The importance for chain completeness of $\mbb{I}$ for domain theory is that, as observed in~\citet{hyland_first_2006}, it produces fixed points of endo-morphisms on a suitably defined class of objects, which could be viewed as \emph{domains}. 

As $L$ by construction is a polynomial functor, it preserves connected limits. This implies the final coalgebra $\ov\omega$ can always be constructed as a sequential limit as follows, 
\[\begin{tikzcd}
	{\ov\omega\cong\lt_{n\in\N}L^n1} \ar[r] & \cdots & {L^21} & L1 & 1
	\arrow[from=1-2, to=1-3]
	\arrow["{{L!}}"', from=1-3, to=1-4]
	\arrow["{{!}}"', from=1-4, to=1-5]
\end{tikzcd}\]
However, the dual statement fails for the initial algebra $\omega$, i.e.\ it is \emph{not} always the following sequential colimit,
\[\begin{tikzcd}
	\emp & {L(\emp)} & {L^2(\emp)} & \cdots
	\arrow["{{?}}"', from=1-1, to=1-2]
	\arrow["{{L(?)}}"', from=1-2, to=1-3]
	\arrow[from=1-3, to=1-4]
\end{tikzcd}\]

\begin{definition}
  $\omega$ is \emph{inductive} if it is the sequential colimit above.
\end{definition}
The failure of inductivity of $\omega$ has been observed in various realisability models for synthetic domain theory~\citep{VANOOSTEN2000233}. As shown by \citet{fiore-plotkin:1996}, the inductivity of $\omega$ is indeed a desirable property: If it holds, then there will be a much closer correspondence between models for \emph{synthetic domain theory} and models for \emph{axiomatic domain theory}. Specifically, for a model of synthetic domain theory where $\omega$ is inductive, one can construct a corresponding model of axiomatic domain theory, where the domains are the \emph{well-complete sets}\footnote{A set $X$ is well-complete if $LX$ is chain complete.}. The inductivity of $\omega$ makes $\ov\omega$ an \emph{inductive fixed point object} in the category of domains in the sense of \emph{loc.\ cit.}, which is part of their requirement for an axiomatic model of domain theory. 

The remaining part of this section is dedicated to showing that under a suitable \gls{QC} assumption, we can indeed show that $\mbb{I}$ is chain complete (\zcref{thm:complete}) and $\omega$ is inductive (\zcref{prop:omegacolimit}). In fact, the proof of \zcref{thm:complete} is the only place we where have used the essential properties of $\sigma$-frames, and we will also indicate why it fails for distributive lattices in \zcref{rem:whynotdis}.

\subsection{Spatiality of countably presented algebras and infinitary domain theory}
The connection between the infinitary aspect of synthetic domain theory and \gls{QC} starts from the observation that the final coalgebra $\ov\omega$ for the functor $L$ can be described as a \emph{spectrum}:

\begin{example}[$\ov\omega$ is a spectrum]\label{exm:ovomegasober}
  As mentioned, the final coalgebra $\ov\omega$ can be characterised as the sequential limit,
  \[\begin{tikzcd}
    \cdots & {\Delta^2} & {\Delta^1} & {\Delta^0} \\
    \cdots & {L^21} & L1 & 1
    \arrow[from=1-1, to=1-2]
    \arrow[from=1-2, to=1-3]
    \arrow["\cong"{description}, from=1-2, to=2-2]
    \arrow[from=1-3, to=1-4]
    \arrow["\cong"{description}, from=1-3, to=2-3]
    \arrow["\cong"{description}, from=1-4, to=2-4]
    \arrow[from=2-1, to=2-2]
    \arrow["{L!}"', from=2-2, to=2-3]
    \arrow["{!}"', from=2-3, to=2-4]
  \end{tikzcd}\]
  where under the isomorphisms, the transition map $\Delta^{n+1} \to \Delta^n$ takes the sequence $i_0 \ge \cdots \ge i_n$ to the final segment $i_1 \ge \cdots \ge i_n$. As observed by \citet[Sec.\ 5.2]{hyland_first_2006}, it also has an equivalent set-theoretic description as the object of infinite descending sequences in $\mbb{I}$,
  \[ \ov\omega \cong \scomp{i : \N \to \mbb{I}}{\fa n\N i_n \ge i_{n+1}}\text{.} \]
  This way, we have a natural equivalence
  \[ \ov\omega \cong \spec\mbb{I}[\N]/\pair{\ms{i}_n \ge \ms{i}_{n+1}}_{n\in\N}\text{.} \]
  Here now $\mbb{I}[\N]/\pair{\ms{i}_n \ge \ms{i}_{n+1}}_{n\in\N}$ is a \emph{countably presented} $\mbb{I}$-algebra.
\end{example}

In general, a countably presented $\mbb{I}$-algebra is of the form $\mbb{I}[\N]/\pair{s_n = t_n}_{n\in\N}$ for some $\N$-indexed lists of terms $s,t \colon \N \to \mbb{I}[\N]$. In particular, all finitely presented $\mbb{I}$-algebras will also be countably presented. Motivated by the above characterisation of $\ov\omega$, we naturally consider the following stronger \gls{QC} principle:


\PrintAxiomSQCC

\begin{remark}
  Our modular development in the previous sections can be now immediately applied to $\ov\omega$ if we assume (\AxiomSQCC). For instance, the description of the observational preorder for sober spaces in \zcref{lem:specorderofsober} now applies to $\ov\omega$, which shows this again coincides with its pointwise order as a subspace of $\mbb{I}^\N$. In particular, it also has both a top and bottom element, thus \zcref{cor:connectedpreservediscrete} now implies $\ov\omega$ is also internally connected. 
\end{remark}

(\AxiomSQCC) combined with the non-triviality axiom (\AxiomNT) has many logical consequences. The crucial observation is that (\AxiomNT) implies a weak form of \emph{Nullstellensatz} result, as already noted by several authors~\citep{blechschmidt:2017,blechschmidt_general_2020,Cherubini_Coquand_Hutzler_2024}:

\begin{lemma}[\AxiomNT]\label{lem:nulls}
  For an sober space $X \cong \spec A$, $X \cong \emp$ iff $A$ is trivial, viz.\ $0=1$ in $A$.
\end{lemma}
\begin{proof}
  The backward direction holds since because $\mbb{I}$ is non-trivial (\AxiomNT), thus there is no homomorphism from a trivial algebra to $\mbb{I}$. For the forward direction: by assumption $A \cong \mbb{I}^{\spec A} \cong \mbb{I}^\emp$, which implies $A$ is trivial. 
\end{proof}

Together with (\AxiomSQCC), this implies the following form of Markov's principle; a similar result is shown by \citet{cherubini2024foundation}:

\begin{lemma}[\AxiomNT, \AxiomSQCC]\label{lem:markov}
  For any $i \in \ov\omega$, we have
  \[ \neg\fa{n}{\N}\ist{i_n} \to \ex n\N\isf{i_n}\text{.} \]
\end{lemma}
\begin{proof}
  Let $i \in \ov\omega$. Note that similar to \zcref{lem:openpropsober}, the proposition $\fa n\N \ist{i_n}$ by construction is the following sober space, 
  \[ \spec\mbb{I}/i \cong \mbb{I}\alg(\mbb{I}/i,\mbb{I}) \cong \fa n\N \ist{i_n}\text{,} \]
  where we have abbreviated the countably presented $\mbb{I}$-algebra $\mbb{I}/\pair{i_n=1}_{n\in\N}$ as $\mbb{I}/i$. Now if we have $\neg\fa n\N \ist{i_n}$, then $\spec\mbb{I}/i \cong \emp$ which by \zcref{lem:nulls} implies $\mbb{I}/i$ is trivial. But this algebra is trivial iff $\ex n\N \isf{i_n}$.
\end{proof}

Equipped with the above result, we can now proceed to study the initial algebra $\omega$ for $L$. \citet{JIBLADZE1997185} has given a beautiful formula for a set-theoretic description of $\omega$ as the following subset of $\ov\omega$, 
\[ \omega \coloneq \scomp{i : \ov\omega}{\fa\phi{\pp} (\fa n{\N} (\ist{i_n} \to \phi) \to \phi) \to \phi}\text{.} \]
For another proof, see e.g.~\citet{VANOOSTEN2000233}. In the presence of \zcref{lem:markov}, this description can be greatly simplified:

\begin{proposition}[\AxiomNT, \AxiomSQCC]\label{prop:omegacolimit}
  $\omega$ is equivalent to the following subset of $\ov\omega$,
  \[ \omega \cong \scomp{i \in \ov\omega}{\ex n{\N} \isf{i_n}}\text{.} \]
  In particular, $\omega$ is the colimit of the following sequence, thus is inductive:
  \[\begin{tikzcd}
    \emp & {\Delta^0} & {\Delta^1} & \cdots \\
    \emp & {L(\emp)} & {L^2(\emp)} & \cdots
    \arrow[from=1-1, to=1-2]
    \arrow["\cong"{description}, from=1-1, to=2-1]
    \arrow[from=1-2, to=1-3]
    \arrow["\cong"{description}, from=1-2, to=2-2]
    \arrow[from=1-3, to=1-4]
    \arrow["\cong"{description}, from=1-3, to=2-3]
    \arrow["{?}"', from=2-1, to=2-2]
    \arrow["{L(?)}"', from=2-2, to=2-3]
    \arrow[from=2-3, to=2-4]
  \end{tikzcd}\]
\end{proposition}

The transition maps $\Delta^n\to \Delta^{n+1}$ above are given by appending $0$ to the end of a descending sequence. Under the identification $\Delta^{n+1}\cong T(\Delta^n)$, the corresponding map $\Delta^n\to T(\Delta^n)$ is the unit component $\eta_T\colon \Delta^n\to T(\Delta^{n})$.

\begin{proof}
  Let $i \in \ov\omega$. It suffices to show that
  \[ \prth{\fa\phi{\pp} (\fa n{\N} (\ist{i_n} \to \phi) \to \phi) \to \phi} \to \ex{n}\N \isf{i_n}\text{.} \]
  Assume the premise. We can instantiate $\phi$ to $\emp$. By assumption, $\neg\fa n\N \neg\neg\ist{i_n}$, which by \zcref{cor:opendnegclose} is equivalent to $\neg\fa n\N \ist{i_n}$. Then \zcref{lem:markov} implies this is $\ex n\N \isf{i_n}$. The fact that this makes $\omega$ into the above sequential colimit follows from \citet[Cor.~1.10]{VANOOSTEN2000233}.
\end{proof}

Using this colimit description of $\omega$ and the fact that $\ov\omega$ is sober, we can show $\mbb{I}$ is chain complete. This again crucially depends on a normal form result for a countably presented $\sigma$-frame, which generalises the finitary version for distributive lattices given in \zcref{prop:simplicesasalgebra}:

\begin{lemma}\label{app:normalsigma}
  For the countably presented $\sigma$-frame $\mbb{I}[\N]/\pair{\ms{i}_n \ge \ms{i}_{n+1}}_{n\in\N}$, an element $p$ can be uniquely written as 
  \[ p = p_0 \vee \bigvee_{n\in\N} p_{n+1} \wedge \ms{i}_n\text{,} \]
  with $p_n \le p_{n+1}$ for all $n$. In other words, $\mbb{I}[\N]/\pair{\ms{i}_n \ge \ms{i}_{n+1}}_{n\in\N}$ is isomorphic to the following $\sigma$-frame with the pointwise order induced by $\mbb{I}$,
  \[ \Delta_\infty \coloneq \scomp{i : \N \to \mbb{I}}{\fa n\N i_n \le i_{n+1}}\text{.} \]
\end{lemma}
\begin{proof}
  We directly prove that $\Delta_\infty$ satisfies the universal property of the countably presented $\sigma$-frame $\mbb{I}[\N]/\pair{\ms{i}_n \ge \ms{i}_{n+1}}_{n\in\N}$. We pick the generators in $\Delta_\infty$ as follows,
  \[ i_n \coloneq (\underbrace{0,\cdots,0}_{n+1 \text{ times}},1,1,\cdots)\text{.} \]
  For any $\sigma$-frame $A$ with $a_0 \ge a_1 \ge \cdots$, we define a map $f_a \colon \Delta_\infty \to A$, which sends $j = (j^0,j^1,\cdots) \in \Delta_\infty$ to
  \[ f_a(i) \coloneq j_0 \vee \bigvee_{n\in\N} j^{n+1} \wedge a_{n}\text{.} \]
  By construction it is easy to see $f_a$ is a $\sigma$-frame morphism. Evidently for any $n\in\N$, 
  \[ f_a(i_n) = \bigvee_{m:\N} i_n^{m+1} \wedge a_m = a_n\text{.} \]
  Furthermore, for any $\sigma$-frame map $f \colon \Delta_\infty \to A$ with $f(i_n) = a_n$, we must have $f = f_a$ because any $j$ in $\Delta_\infty$ can be written as
  \[ j = j^0 \vee \bigvee_{n\in\N} j^{n+1} \wedge i_n\text{,} \]
  which implies that
  \[ f(j) = j^0 \vee \bigvee_{n\in\N}j^{n+1} \wedge a_n = f_a(j)\text{.} \]
  This completes the proof.
\end{proof}

\begin{theorem}[\AxiomNT, \AxiomSQCC]\label{thm:complete}
  $\mbb{I}$ is right orthogonal to the inclusion $\omega\hook\ov\omega$.
\end{theorem}
\begin{proof}
  Since $\ov\omega$ is now sober, we have
  \[ \mbb{I}^{\ov\omega} \cong \mbb{I}[\N]/\pair{\ms{i}_n \ge \ms{i}_{n+1}}_{n\in\N}\text{.} \]
  On the other hand, since $\omega$ is the colimit of $\Delta^n$ and they are sober, we have
  \[ \mbb{I}^\omega \cong \lt_{n\in\N}\mbb{I}^{\Delta^n} \cong \lt_{n\in\N}\mbb{I}[n]/\ms{i}_0\ge\cdots\ge \ms{i}_{n-1}\text{.} \]
  Note that the transition maps induced by $\eta_T \colon \Delta^n \to T(\Delta^n)\cong \Delta^{n+1}$ under \gls{QC} gives us the following maps on algebras:
  \[
  \begin{tikzcd}
    \mbb{I}^{\Delta^{n+1}} \ar[r, "\mbb{I}^{\eta_T}"] & \mbb{I}^{\Delta^n} \\ 
    \mbb{I}[n\!+\!1]/\ms{i}_0 \ge \cdots \ge \ms{i}_{n} \ar[u, "\cong"] \ar[r, "\ms{i}_{n} \mapsto 0"'] & \mbb{I}[n]/\ms{i}_1 \ge \cdots \ge \ms{i}_{n-1} \ar[u, "\cong"']
  \end{tikzcd}
  \]
  Hence, it suffices to show that $\mbb{I}[\N]/\pair{\ms{i}_n \ge \ms{i}_{n+1}}_{n\in\N}$ is indeed the sequential limit of the above $\mbb{I}$-algebras,
  \[\begin{tikzcd}
    \cdots & {\mbb{I}[n\!+\!1]/\ms{i}_0 \ge \cdots \ge \ms{i}_{n}} & {\mbb{I}[n]/\ms{i}_0 \ge \cdots \ge \ms{i}_{n-1}} & \cdots \\
    & {\mbb{I}[\N]/\pair{\ms{i}_n \ge \ms{i}_{n+1}}_{n\in\N}}
    \arrow[from=1-1, to=1-2]
    \arrow["{i_{n} \mapsto 0}", from=1-2, to=1-3]
    \arrow[from=1-3, to=1-4]
    \arrow[dashed, from=2-2, to=1-1]
    \arrow["{f_{n+1}}"{description}, from=2-2, to=1-2]
    \arrow["{f_n}"{description}, from=2-2, to=1-3]
    \arrow[curve={height=15pt}, dashed, from=2-2, to=1-4]
  \end{tikzcd}\]
  where the map $f_n$ takes $\ms{i}_k$ to itself for $k\le n$, and takes $\ms{i}_k$ to $0$ for $k \ge n$. We can more directly compute the above sequential limit via \zcref{prop:simplicesasalgebra},
  \[\begin{tikzcd}
    \cdots & {\mbb{I}[n\!+\!1]/\ms{i}_0 \ge \cdots \ge \ms{i}_{n}} & {\mbb{I}[n]/\ms{i}_0 \ge \cdots \ge \ms{i}_{n-1}} & \cdots \\
    \cdots & {\Delta_{n+2}} & {\Delta_{n+1}} & \cdots
    \arrow[from=1-1, to=1-2]
    \arrow["{{\ms{i}_{n} \mapsto 0}}", from=1-2, to=1-3]
    \arrow["\cong", from=1-2, to=2-2]
    \arrow[from=1-3, to=1-4]
    \arrow["\cong", from=1-3, to=2-3]
    \arrow[from=2-1, to=2-2]
    \arrow["\pi_{n}", from=2-2, to=2-3]
    \arrow[from=2-3, to=2-4]
  \end{tikzcd}\]
  where $\pi_{n} : \Delta_{n+1} \to \Delta_{n+1}$ forgets the last entry, i.e.\ it takes $\ms{i}_0 \le \cdots \le \ms{i}_{n+1}$ to $\ms{i}_0 \le \cdots \le \ms{i}_{n}$. Hence, the sequential limit is given by
  \[ \lt_{n\in\N} \Delta_{n} \cong \Delta_\infty\text{.} \]
  Now the desired result follows from \zcref{app:normalsigma}.
\end{proof}

\begin{remark}
  As indicated in \zcref{rem:qcreplete,rem:specarereplete}, both spectra and spatial $\mbb{I}$-algebras are replete. The above result then implies that they are also chain complete, i.e.\ orthogonal to $\omega\hook\ov\omega$.
\end{remark}

\begin{corollary}[\AxiomNT, \AxiomSQCC]
  $\mbb{I}^{\ov\omega} \cong \mbb{I}^\omega \cong \Delta_\infty$. In particular, $\omega$ is not sober.
\end{corollary}
\begin{proof}
  By the above proof, we have $\mbb{I}^{\ov\omega} \cong \mbb{I}[\N]/\pair{\ms{i}_n \ge \ms{i}_{n+1}}_{n\in\N} \cong \Delta_\infty$.
\end{proof}

\begin{remark}[$\sigma$-frames vs.\ distributive lattices]\label{rem:whynotdis}
  Note that the inductivity of $\omega$ as shown in \zcref{prop:omegacolimit} is \emph{independent} of working with distributive lattices or $\sigma$-frames. The completeness result shown in \zcref{thm:complete} is the only one in this paper that works for $\sigma$-frames and not for distributive lattices more generally. The reason is exactly because the normal form in \zcref{app:normalsigma} for the countably presented \emph{distributive lattice} $\mbb{I}[\N]/\pair{\ms{i}_n \ge \ms{i}_{n+1}}_{n\in\N}$ will \emph{not} be isomorphic to $\Delta_\infty$. More specifically, for the countably presented distributive lattice $\mbb{I}[\N]/\pair{\ms{i}_n \ge \ms{i}_{n+1}}_{n\in\N}$, since it is a sequential colimit of the following finitely presented distributive lattices
  \[ 
  \begin{tikzcd}
    \cdots \ar[r] & \mbb{I}[n]/\ms{i}_{0} \ge \cdots \ge \ms{i}_{n-1} \ar[r, hook] & \mbb{I}[n\!+\!1]/\ms{i}_0 \ge \cdots \ge \ms{i}_n \ar[r, hook] & \cdots
  \end{tikzcd}
  \]
  and, due to the fact that the theory of distributive lattices is \emph{finitary}, this sequential colimit of \emph{algebras} is computed the same as their \emph{underlying sets}. Via \zcref{prop:simplicesasalgebra}, the result can be identified as the subset of $\Delta_\infty$,
  \[ \Delta_\omega \coloneq \scomp{i \in \Delta_\infty}{\ex n\N \ist{i_n}}\text{,} \]
  where $\Delta_\omega \hook \Delta_\infty$ in some sense is the order-dual to the inclusion $\omega\hook\ov\omega$. The fact that we have \emph{all} infinite increasing sequences in the case for $\sigma$-frames is clearly due to the fact that we have countable disjunctions in $\sigma$-frames, as when identifying $p$ in $\mbb{I}[\N]/\pair{\ms{i}_n \ge \ms{i}_{n+1}}_{n\in\N}$ with $p = p_0 \vee \bigvee_{n\in\N}p_{n+1}\wedge \ms{i}_n$.
\end{remark}

\begin{remark}[Limits of algebras induce locality for $\mbb{I}$]\label{rem:limofalgloc}
  By inspecting the proof of \zcref{thm:complete}, it is clear that $\mbb{I}$ is chain complete \emph{precisely} because there is a specific \emph{limiting diagram} of spatial $\mbb{I}$-algebras. In general, any such limit diagram will induce an orthogonality property satisfied by the all spectra, and we will see many more examples in \zcref{sec:local}. From \zcref{rem:whynotdis} it is clear whether a diagram of $\mbb{I}$-algebras is a limit heavily relies on the underlying algebraic theory. This is again a perfect example of how the algebraic properties of a theory has significant consequences for the internal logic of its classifying topos.
\end{remark}

\section{Local properties for the interval}\label{sec:local}

In this section we review some of the locality axioms we have introduced in \zcref{sec:locality}. As mentioned in the introduction, the observation is that even if we do not assume them to be true globally for the interval $\mbb{I}$, we can still show the maps representing the locality axiom to be \emph{left orthogonal} to $\mbb{I}$.  
Unlike chain completeness in \zcref{thm:complete}, the local conditions in \zcref{sec:locality} are all induced by limits of \emph{f.p.} $\mbb{I}$-algebras. Thus in this section, it does not matter whether we work with distributive lattices or $\sigma$-frames.

\begin{remark}
  Comment that though here we have formulated the results for ``sets'' to keep the foundation aligned with the remaining of this thesis, the same argument works in \gls{HoTT} for arbitrary types, thus this really has applications in synthetic higher category theory.
\end{remark}

Starting from the simplest example, let us recall the local property (\AxiomNT). In general it is \emph{not} necessarily $\emp$, if we do not assume (\AxiomNT). However, we can look at the localisation class that thinks this map is invertible. The following fact shows that, from the perspective of $\mbb{I}$, (\AxiomNT) indeed holds:

\begin{proposition}[\AxiomSQCI]\label{specisnontrivial}
  $\mbb{I}$ is right orthogonal to $\emp \hook (0 = 1)$.
\end{proposition}
\begin{proof}
  Observe the proposition $0 = 1$ is propositionally equivalent to the (sober) spectrum $\spec(\mbb{I}/0=1)$. By \zcref{prop:duality} we have:
  \[ 
    \mbb{I}^{0=1} \cong 
    (\spec \mbb{I}[\ms{i}])^{\spec{\mbb{I}/0=1}}
    \cong 
    \mbb{I}\alg(\mbb{I}[\ms{i}], \mbb{I}/0=1) 
    \cong 
    \mathbf{1}
  \]
  The last isomorphism is due to $\mbb{I}/0=1$ being the terminal $\mbb{I}$-algebra, since its underlying set is the singleton.
\end{proof}

As another example, consider the simplicial axiom (\AxiomSL). We skip the discussion of (\AxiomL) and (\AxiomCL), because (\AxiomSL) is strictly stronger, and it has a better-known geometric meaning. Recall from \zcref{sec:linearity-and-coskeletality} that (\AxiomSL) can be represented geometrically as an embedding $\Delta^2\cup\Delta_2 \hook \mbb{I}^2$.

\begin{definition}[Triangulated set]
  A set $X$ is said to be \emph{triangulated} when it is right orthogonal to the inclusion $\Delta^2 \cup \Delta_2 \hook \mbb{I}^2$.
\end{definition}

Again, (\AxiomSL) holds globally iff the above embedding is an equivalence, thus every set in this case will be triangulated. When it does not hold globally, we can still show:

\begin{proposition}[\AxiomSQCF]\label{specistriangulated}
  $\mbb{I}$ is triangulated.
\end{proposition}
\begin{proof}
  To show that $\mbb{I}$ is right orthogonal to $\Delta^2 \cup \Delta_2 \hook \mbb{I}^2$, it is equivalent to verify that for any $f,g$ indicated below making the solid diagram below commute, there exists a unique lift $h$ making the whole diagram commute:
  \[\begin{tikzcd}
    \mbb{I} \ar[r, "\delta"] \ar[d, "\delta"'] & \Delta^2 \ar[r, "f"] \ar[d] & \mbb{I} \\
    \Delta_2 \ar[urr, "g"{description, near start}] \ar[r] & \mbb{I}^2 \ar[ur, dashed, "h"']
    \arrow["\lrcorner"{anchor=center, pos=0.125}, draw=none, from=1-1, to=2-2]
  \end{tikzcd}\]
  Now by (\AxiomSQCF), since the vertices of the left square are all sober, equivalently it suffices to show we have a pullback of $\mbb{I}$-algebras,
  \[\begin{tikzcd}
    {\mbb{I}[\ms{i},\ms{j}]} \ar[r] \ar[d] & {\mbb{I}[\ms{i},\ms{j}]/\ms{i} \ge \ms{j}} \ar[d, "{\ms{i},\ms{j} \mapsto\ms{k}}"] \\
    {\mbb{I}[\ms{i},\ms{j}]/\ms{i} \le \ms{j}} \ar[r, "{\ms{i},\ms{j}\mapsto\ms{k}}"'] & {\mbb{I}[\ms{k}]}
    \arrow["\lrcorner"{anchor=center, pos=0.125}, draw=none, from=1-1, to=2-2]
  \end{tikzcd}\]
  Recall from \zcref{prop:simplicesasalgebra} that the normal form of elements in $\mbb{I}[\ms{i},\ms{j}]/\ms{i}\ge \ms{j}$. An element $p$ there can be viewed as the polynomial 
  \[ p = p_{0,0} \vee \ms{i}\wedge p_{1,0} \vee \ms{j}\wedge p_{1,1}\text{,} \]
  with $p_{0,0} \le p_{1,0} \le p_{1,1}$. And similarly for $q$ in $\mbb{I}[\ms{i},\ms{j}]/\ms{i} \le \ms{j}$, 
  \[ q = q_{0,0} \vee \ms{j} \wedge q_{0,1} \vee \ms{i} \wedge q_{1,1}\text{.} \]
  They agree in $\mbb{I}[\ms{k}]$ iff 
  \[ p_{0,0} = q_{0,0}, \quad p_{1,1} = q_{1,1}\text{.} \]
  This then gives us a polynomial in $\mbb{I}[\ms{i},\ms{j}]$, which we write as follows,
  \[ (p_{0,0} \vee \ms{i} \wedge p_{1,0}) \vee \ms{j} \wedge (q_{0,1} \vee \ms{i} \wedge q_{1,1})\text{,} \]
  where $p_{0,0} = q_{0,0} \le q_{0,1}$ and $p_{1,0} \le p_{1,1} = q_{1,1}$. 
  By the general normal form for $A[\ms{i}]$ for any $\mbb{I}$-algebra $A$, if we view $\mbb{I}[\ms{i},\ms{j}]$ as $\mbb{I}[\ms{i}][\ms{j}]$, the above exactly corresponds to the normal form of polynomials in $\mbb{I}[\ms{i},\ms{j}]$; see also \citet[Ch.\ 1, Thm.\ 10.21]{lausch2000algebra}. This means the above is a pullback of $\mbb{I}$-algebras.
\end{proof}



Finally, we discuss the even stronger locality condition (\AxiomOneCS). As mentioned in \zcref{sec:linearity-and-coskeletality}, the additional property of (\AxiomOneCS) is characterised by the embedding $\partial\Delta_2 \hook \Delta_2$.

\begin{definition}[1-coskeletal set]
  A set $X$ is said to be \emph{1-coskeletal} when it is both triangulated and right orthogonal to the boundary inclusion $\partial\Delta_2 \hook \Delta_2$.
\end{definition}

\begin{proposition}[\AxiomSQCF]\label{specis1t}
  $\mbb{I}$ is 1-coskeletal.
\end{proposition}
\begin{proof}
  Completely similar to the proof of \zcref{specistriangulated}, it suffices to show that we have a limit diagram of $\mbb{I}$-algebras as follows,
  \[\begin{tikzcd}
    & {\mbb{I}[\ms{i},\ms{j}]/\ms{i}\le \ms{j}} \\
    \mbb{I}[\ms{i}] & \mbb{I}[\ms{k}] & \mbb{I}[\ms{j}] \\
    \mbb{I} & \mbb{I} & \mbb{I}
    \arrow["{\ms{j} \mapsto 1}"', from=1-2, to=2-1]
    \arrow["{\ms{i},\ms{j} \mapsto k}"{description}, from=1-2, to=2-2]
    \arrow["{\ms{i}\mapsto 0}", from=1-2, to=2-3]
    \arrow["\ms{i}\mapsto 1"', from=2-1, to=3-1]
    \arrow["\ms{i}\mapsto 0"{description, pos=0.75}, from=2-1, to=3-2]
    \arrow["\ms{k}\mapsto 1"{description, pos=0.75}, from=2-2, to=3-1]
    \arrow["\ms{k}\mapsto 0"{description, pos=0.75}, from=2-2, to=3-3]
    \arrow["\ms{j}\mapsto 1"{description, pos=0.75}, from=2-3, to=3-2]
    \arrow["\ms{j}\mapsto 0", from=2-3, to=3-3]
  \end{tikzcd}\]
  Now by the normalisation theorem, an element in the limit consists of $p$ in $\mbb{I}[\ms{i}]$, $q$ in $\mbb{I}[\ms{j}]$ and $r$ in $\mbb{I}[\ms{k}]$, such that
  \[ p_1 = r_1, \quad p_0 = q_1, \quad r_0 = q_0\text{.} \]
  This exactly corresponds to a normal form in the algebra $\mbb{I}[\ms{i},\ms{j}]/\ms{i}\le\ms{j}$ with $q_0 \le q_1 = p_0 \le p_1$ as follows (cf.\ \zcref{prop:simplicesasalgebra})
  \[ q_0 \vee \ms{j} \wedge q_1 \vee \ms{i} \wedge p_1\text{.} \]
  Hence, the above is a limit diagram of $\mbb{I}$-algebras.
\end{proof}

Besides the locality principles discussed in \zcref{sec:locality}, we also consider the orthogonality classes emerging from \emph{synthetic category theory}, as introduced by \citet{riehl2017type}. 
A synthetic category will be a set that satisfies the internal orthogonality condition of being \emph{Segal complete}, \emph{Rezk complete}, and (according to some sources) \emph{simplicial}. Being simplicial is a sheaf condition that strengthens the triangulation property by stabilising the left class under pullback.

We start with the Segal completeness condition. Besides the outer horn discussed in \zcref{exm:hornsober}, we can also define the inner horn $\Lambda^2_1$ as a pushout,
\[
  \begin{tikzcd}
    1 \ar[d, "0"'] \ar[r, "1"] & \mbb{I} \ar[d] \\
    \mbb{I} \ar[r] & \Lambda^2_1
    \arrow["\lrcorner"{anchor=center, pos=0.125, rotate=180}, draw=none, from=2-2, to=1-1]
  \end{tikzcd}
\]
As a subset of $\Delta_2$, it can be identified as $\scomp{(i,j) \in \Delta_2}{\isf{i} \vee \ist{j}}$.

\begin{definition}[Segal complete sets]
  $X$ is called \emph{Segal complete} when it is right orthogonal to $\Lambda^2_1 \hook \Delta_2$.
\end{definition}

\begin{remark}[Path transitivity]
  The Segal completeness condition has also been studied independently by \citet{fiore2001domains} under the name of \emph{path transitivity}.
\end{remark}

\begin{proposition}[\AxiomSQCP]
  $\mbb{I}$ is Segal complete.
\end{proposition}
\begin{proof}
  By \zcref{prop:simplicesasalgebra} we have the following canonical isomorphism:
  \[ 
    \Delta^{n+1} \cong 
    \mbb{I}[n]/i_1\leq\cdots\leq i_n
  \]
  
  In particular, we have $\Delta^3\cong \mbb{I}[\ms{i},\ms{j}]/\ms{i}\leq\ms{j}$ and $\Delta^2 \cong \mbb{I}[\ms{i}]$ and under Phoa's principle (SQCP), we have $\mbb{I}^\mbb{I} \cong \Delta^2\cong\mbb{I}[\ms{i}]$. Therefore, to show that $\mbb{I}$ is right orthogonal to the inner horn comparison map, it suffices to show that the following is a fibre product of $\mbb{I}$-algebras:
  \[
  \begin{tikzcd}
    \mbb{I}[\ms{i},\ms{j}]/\ms{i}\le\ms{j} \ar[r, "\ms{j} \mapsto 1"] \ar[d, "\ms{i} \mapsto 0"'] & \mbb{I}[\ms{i}] \ar[d, "\ms{i}\mapsto 0"] \\
    \mbb{I}[\ms{j}] \ar[r, "\ms{j} \mapsto 1"'] & \mbb{I}
  \end{tikzcd}
  \]
  By the normal form theorem, an element in the pullback is given by $p$ in $\mbb{I}[\ms{i}]$ and $q$ in $\mbb{I}[\ms{j}]$ with $p_0 = q_1$. This way, it again corresponds to the following normal form in $\mbb{I}[\ms{i},\ms{j}]/\ms{i} \le \ms{j}$ by \zcref{prop:simplicesasalgebra},
  \[ q_0 \vee \ms{j} \wedge q_1 \vee \ms{i} \wedge p_1\text{.} \]
  Hence, the above is again a pullback.
\end{proof}

Next we consider the Rezk completeness condition. Following \citet{buchholtz2021synthetic}, we can define the set $\mbb E$ classifying categorical equivalences as the colimit of the following diagram,
\[
\begin{tikzcd}
	& \mbb{I} && \mbb{I} && \mbb{I} \\
	1 && {\Delta^2} && {\Delta^2} && 1
	\arrow[from=1-2, to=2-1]
	\arrow["{i \mapsto (i,i)}"{description}, from=1-2, to=2-3]
	\arrow["{i \mapsto (i,0)}"{description}, from=1-4, to=2-3]
	\arrow["{i \mapsto (1,i)}"{description}, from=1-4, to=2-5]
	\arrow["{i \mapsto (i,i)}"{description}, from=1-6, to=2-5]
	\arrow[from=1-6, to=2-7]
\end{tikzcd}
\]


\begin{definition}[Rezk sets]
  A set $X$ is \emph{Rezk complete} when it is right orthogonal to $\mbb E \to 1$.
\end{definition}

\begin{proposition}[\AxiomSQCP]
  $\mbb{I}$ is Rezk complete.
\end{proposition}
\begin{proof}
  Notice by \zcref{thm:phoa-spectrum} $\mbb{I}^\mbb{I}$ classifies the canonical order on $\mbb{I}$, which is antisymmetric by definition. Thus $\mbb{I}$ will be Rezk complete.
\end{proof}

In fact, $\mbb{I}$ is not only a synthetic category but in fact a synthetic \emph{poset}; this property too can be expressed in terms of orthogonality. To that end, we define the ``walking parallel pair'' $\I_{\rightrightarrows}$ to be the following pushout:
\[
\begin{tikzcd}
  2 \ar[r] \ar[d] & \mbb{I} \ar[d] \\ 
  \mbb{I} \ar[r] & \I_{\rightrightarrows}
  \arrow["\lrcorner"{anchor=center, pos=0.125, rotate=180}, draw=none, from=2-2, to=1-1]
\end{tikzcd}
\]

\begin{definition}[$\mbb{I}$-separated sets]
  $X$ is called \emph{$\mbb{I}$-separated} when it is right orthogonal to $\I_{\rightrightarrows} \to \mbb{I}$.
\end{definition}

Equivalently, $X$ is $\mbb{I}$-separated iff $X^\mbb{I} \to X \times X$ is an embedding. It is an immediate consequence of (\AxiomSQCP) that $\mbb{I}$ is $\mbb{I}$-separated.

The notion of synthetic categories and synthetic posets are formulated as these orthogonality classes:

\begin{definition}[Synthetic categories and synthetic posets]
  A set $X$ is a \emph{synthetic category}, if it is Segal and Rezk complete. We say it is a \emph{synthetic poset}, if it is also $\mbb{I}$-separated.
\end{definition}

\begin{theorem}[\AxiomSQCF]\label{thm:soberposet}
  Any spectrum or spatial $\mbb{I}$-algebra will be a synthetic poset.
\end{theorem}
\begin{proof}
  This follows from $\mbb{I}$ being a synthetic poset, and spectra and spatial algebras are \emph{replete} as indicated in \zcref{rem:qcreplete,rem:specarereplete}.
\end{proof}

At the end of this section, we also discuss the example of $\omega$, which as we have seen is \emph{not} sober. We first observe its observational preorder again coincides with its expected pointwise order viewed as a subspace of $\mbb{I}^\N$:

\begin{lemma}[\AxiomNT, \AxiomSQCC]\label{speconomegaiscan}
  The inclusion $\omega \hook \ov\omega$ is an order embedding for the observational preorder. In particular, the observational preorder on $\omega$ is induced pointwise as a subspace of $\ov\omega\hookrightarrow \mbb{I}^\N$.
\end{lemma}
\begin{proof}
  The fact that $\omega\hook\ov\omega$ is an order embedding for the observational preorder follows from chain completeness of $\mbb{I}$ in \zcref{thm:complete}. The fact that the observational preorder on $\ov\omega$ coincides with its satisfaction order as a spectrum follows from \zcref{lem:specorderofsober} by the fact that it is sober under (\AxiomSQCC).
\end{proof}

Furthermore, we can show that $\omega$ also satisfies a version of the Phoa principle, i.e.\ the path space $\omega^\mbb{I}$ again classifies its pointwise order. For this we need to compute $\omega^\mbb{I}$.

Recall from \zcref{prop:omegacolimit} that, under (\AxiomNT) and (\AxiomSQCC), $\omega$ is a colimit $\omega \cong \ct_{n\in\N}\Delta^n$. As indicated in \zcref{rem:whynotdis}, this does not depend on working with distributive lattices or $\sigma$-frames. Type-theoretically, we have also shown that $\omega$ can be realised as the following subspace of $\ov\omega$,
\[ \omega \cong \scomp{i \in \ov\omega}{\ex n\N \isf{i_n}}\text{.} \]
This way, the inclusion $\Delta^n \hook \omega$ can be viewed as follows, 
\[ \Delta^n \cong \scomp{i \in \omega}{\isf{i_n}}\text{,} \]
which implies $\Delta^n \hook \omega$ is \emph{downward closed}. This allows us to directly compute $\omega^\mbb{I}$:

\begin{proposition}[\AxiomNT, \AxiomSQCC]\label{lem:intervalcommuteomega}
  We have a family of isomorphisms
  \[ \omega^\mbb{I} \cong \big({\ct_{n\in\N}\Delta^n}\big)^\mbb{I} \cong \ct_{n\in\N}(\Delta^n)^\mbb{I}\text{,} \]
  and similarly by replacing $\mbb{I}$ with $\mbb{I}^n$ or $\Delta^n$. In particular, $\omega^\mbb{I}$ classifies the pointwise order on $\omega$.
\end{proposition}
\begin{proof}
  We show the following canonical map is an equivalence,
  \[ \ct_{n\in\N}(\Delta^n)^\mbb{I} \to \big({\ct_{n\in\N}\Delta^n}\big)^\mbb{I}\text{.} \]
  It is evident this map is an embedding, hence it suffices to show it is surjective. Given $f \colon \mbb{I} \to \ct_{n\in\N}\Delta^n$. By assumption, there merely exists $n\in\N$ that $f(1)$ factors through $\Delta^n \hook \omega$. Now the claim is that the entire map $f$ factors through $\Delta^n$. This indeed holds, since we have shown in \zcref{speconomegaiscan} that the observational preorder on $\omega$ is pointwise. By monotonicity, for any $i \in \mbb I$ we have $fi \obsle f1$, which implies $fi$ belongs to $\Delta^n$ as well, since $\Delta^n \hook \omega$ is a downward closed. The same holds for cubes or simplices since they all have a top element. This way, $\omega^\mbb{I}$ classifies the pointwise order on $\omega$ since all simplices $\Delta^n$ satisfy Phoa's principle as shown in \zcref{thm:phoa-spectrum}.
\end{proof}

As another consequence, we can also show the following general result establishing a large family of orthogonality conditions satisfied by $\omega$:

\begin{theorem}[\AxiomNT, \AxiomSQCC]\label{thm:omegaortho}
  Let $f \colon X \to Y$ be a map where $X,Y$ are finite colimits of cubes or simplices. If each $\Delta^n$ is $f$-local, then so is $\omega$.
\end{theorem}
\begin{proof}
  Let $Y \cong \ct_{r}Y_r$ be a finite colimit with each $Y_r$ a simplex or a cube.
  \[ \omega^Y \cong \lt_r\big({\ct_{n\in\N}\Delta^n}\big)^{Y_r} \cong \lt_r\ct_{n\in\N}(\Delta^n)^{Y_r} \cong \ct_{n\in\N}\lt_r(\Delta^n)^{Y_r} \cong \ct_{n\in\N}(\Delta^n)^Y \]
  The second isomorphism holds by \zcref{lem:intervalcommuteomega}; the third holds since finite limits commute with sequential colimits. Thus, if each $\Delta^n$ is $f$-local, then so is $\omega$.
\end{proof}

\begin{remark}[Localisation classes closed under sequential colimits]
  Notice that in \zcref{lem:intervalcommuteomega} we do not need to use any form of choice principle to show the exponential $(-)^\mbb{I}$ commutes with the sequential colimit $\ct_{n\in\N}\Delta^n$, exactly because $\Delta^n \hook \omega$ is downward closed. However, for general sequential colimits, the same proof still goes through if we assume:
  \begin{axiom}[Projectivity of $\mbb{I}$]
    $\mbb I$ is a projective set.
  \end{axiom}
  Furthermore, the same holds for cubes and simplices since the sets satisfying the choice principle are closed under finite products and retracts. Assuming $\mbb{I}$ is projective, following the proof of \zcref{thm:omegaortho}, one can show more generally that any orthogonality class specified by maps between finite colimits of simplices or cubes are always \emph{closed under sequential colimits}. However, not in every model of \gls{QC} $\mbb{I}$ will be projective. Hence, we do not include this result in the main text, as we tend to keep our assumptions as minimalistic as possible.
\end{remark}

% \begin{figure}[ht]
%   \textbf{Quasi-coherence axioms:}
%   \PrintAxiomSQCI
%   \PrintAxiomSQCID
%   \PrintAxiomSQCP
%   \PrintAxiomSQCF
%   \PrintAxiomSQCC

%   \medskip
%   \textbf{Locality axioms:}
%   \PrintAxiomNT
%   \PrintAxiomL
%   \PrintAxiomCL
%   \PrintAxiomSL
%   \PrintAxiomOneCS
%   \vspace{6ex}

%   \caption{Summary of axioms considered in this paper. Unlike the usual axiom sets for synthetic domain theory~\citep{hyland_first_2006}, the goal of these axioms is not to modularly specify all the properties of the interval with minimal overlap but instead to identify a sequence of increasingly strong descriptions of intervals that may play a role in synthetic domain theory.}
% \end{figure}

\section{New models for synthetic domain theory}\label{sec:newmodelforsdt}

It is now instructive to discuss models for the axioms we have used in the previous sections. For simplicity, we now work externally under the assumption that the base topos $\Set$ is the category of sets under (a model of) \gls{ZFC}. Our axioms for synthetic domain theory can be organised into two classes:

\begin{enumerate}
  \item The \gls{QC} principles (\AxiomSQCF), (\AxiomSQCC);
  \item Various locality axioms discussed in \zcref{sec:locality}.
\end{enumerate}

By the developments in~\zcref{sec:qcforhorn}, we know that the strongest \gls{QC} principle (\AxiomSQCC) will be valid in the classifying $\omega_1$-topos of the $\omega_1$-Horn theory of $\sigma$-frames, and the various locality axioms will be valid in the corresponding subtopoi.

Let $\sFrm$ be the category of $\sigma$-frames (cf.~\zcref{exa:pshsigmaframe}). It is well-known the dual category of $\DL\fp \cong \sFrm\fp$ is the category $\Pos\fp$ of finite posets. In this case, we can also have a fairly explicit description of the larger dual category of $\sFrm\cp$. The first observation is that any countably presented $\sigma$-frame $A$ is indeed a \emph{frame}, i.e.\ it has \emph{all} joins and finite meets, which distributes with each other. This is clear for all finitely presented $\sigma$-frames, since they are simply finite distributive lattices. To see this for the countable case, consider the countably generated free $\sigma$-frame $\Sigma[\N]$:

\begin{lemma}\label{lem:cgfreesframe}
  We have an isomorphism of $\sigma$-frames
  \[ \Sigma[\N] \cong \Pos(P_f(\N),\Sigma)\text{,} \]
  where $P_f(\N)$ is the poset of finite subsets of $\N$.
\end{lemma}
\begin{proof}
  The free $\sigma$-frame can be generated by first freely adding finite meets to the discrete poset $\N$, and then freely adding all countable joins. The first step results in the poset $P_f(\N)\op$. Now since $P_f(\N)\op$ is countable, freely adding all countable joins is equivalently freely adding \emph{all} joins, which is achieved by the presheaf construction. This way,
  \[ \Sigma[\N] \cong \Pos((P_f(\N)\op)\op,\Sigma) \cong \Pos(P_f(\N),\Sigma). \qedhere \]
\end{proof}

\begin{corollary}\label{cor:dualsframe}
  There is a fully faithful embedding that preserves all countable colimits,
  \[ \sFrm\cp \hook \Frm\text{.} \]
  This is also fully faithful when post composed with $\ms{pt} \colon \Frm\op \to \Topp$, where $\Topp$ is the category of topological spaces.
\end{corollary}
\begin{proof}
  Any countably presented $\sigma$-frame will be isomorphic to one of the form $\Sigma[\N]/R$ for some countably generated congruence $R$. By \zcref{lem:cgfreesframe}, $\Sigma[\N]$ is a frame, so is any of its quotient. By Theorem 6.2.4 of \citet{makkai1977first}, such frames are indeed \emph{spatial}, hence they fully faithfully embed into the category of topological spaces via the functor $\pt$.
\end{proof}

This way, we can view $\sFrm\cp\op$ as a certain class of topological spaces. Notice that since all countably presented $\sigma$-frames will be quotients of $\Sigma[\N]$, their dual spaces will be a subspace of $\pt(\Sigma[\N])$, which we can compute quite easily:

\begin{lemma}
  The space of points of $\Sigma[\N]$ is Scott's graph model $G$, which is the countable product of the Sierpi\'nski space $G \cong \Sigma^\N$.
\end{lemma}
\begin{proof}
  Note $\pt$ takes colimits in $\Frm$ to limits in $\Topp$, since $\pt \colon \Loc \to \Topp$ is a right adjoint. Since $\Sigma[\N]$ is the countable coproduct of the free frame on one generator, we have
  \[ \pt(\Sigma[\N]) \cong \pt(\Sigma[\ms{x}])^\N \cong \Sigma^\N\text{,} \]
  Here $\pt(\Sigma[\ms{x}]) \cong \Sigma$ follows from a simple computation.
\end{proof}

Write $\wTop$ as the essential image of the fully faithful functor $\pt \colon \sFrm\cp\op \to \Topp$, whose objects will all be subspaces of $G$. We would then have the following diagram,
\[
\begin{tikzcd}
  \sFrm\fp\op \ar[d, hook] \ar[r, "\simeq"] & \Pos\fp \ar[d, hook] \\
  \sFrm\cp\op \ar[r, "\simeq"'] & \wTop
\end{tikzcd}
\]
where here the inclusion $\Pos\fp \hook \wTop$ simply takes each finite poset to its Alexandroff topological space. Hence, at the presheaf level, we have a local geometric morphism
\[ \psh(\wTop) \surj \psh(\Pos\fp)\text{.} \]
In this case, the representable presheaf on the Sierpinski space will again be a $\sigma$-frame $\mbb{I}$ in the presheaf $\psh(\wTop)$.

Below we discuss the corresponding admissible topologies on $\wTop$, modelling the various locality principles we have considered in \zcref{sec:locality}. We encourage the readers to notice the connection between the topologies we discuss below, and the developments in \zcref{sec:local}.

\begin{example}[\AxiomNT]
  To model (\AxiomNT), we would want the empty sieve on $\pt(\Sigma/0=1) \cong \emp$ to be a covering. Since $\emp$ is a strict initial object in $\wTop$, this is a subcanonical topology. Thus, this gives us a least topology $J_{\ms{NT}}$ that models (\AxiomNT), and we have
  \[ \sh(\wTop,J_{\ms{NT}}) \cong \psh(\wTop_+)\text{,} \]
  where $\wTop_+$ is the full subcategory of $\wTop$ excluding $\emp$. The induced local geometric morphism now is given by 
  \[ \Gamma \colon \psh(\wTop_+) \surj \psh(\Pos_{\mr{\omega,+}})\text{,} \]
  where, again $\Pos_{\mr{\omega,+}}$ is the full subcategory of $\Pos_{\mr{\omega}}$ consisting of non-empty posets, and $\psh(\Pos_{\mr{\omega,+}})$ is the classifying topos for distributive lattices that are non-trivial.
\end{example}

\begin{example}[\AxiomL]
  The additional axiom for (\AxiomL) besides (\AxiomNT) is that 
  \[ x \vee y = 1 \vdash x = 1 \vee y = 1\text{,} \] 
  This means the dual embeddings of the following quotients need to form a covering family,
  \[ \Sigma[\ms{y}] \cong \Sigma[\ms{x},\ms{y}]/\ms{x} \twoheadleftarrow \Sigma[\ms{x},\ms{y}]/\ms{x}\vee \ms{y} \surj \Sigma[\ms{x},\ms{y}]/\ms{y} \cong \Sigma[\ms{x}]\text{.} \]
  It is easy to see that $\pt(\Sigma[\ms{x},\ms{y}]/\ms{x} \vee \ms{y})$ is the space $\Lambda^2_1$ obtained by glueing the open points of two copies of Sierpinski spaces together, i.e.\ it is the following pushout:
  \[
  \begin{tikzcd}
    1 \ar[d, "1"'] \ar[r, "1"] & \Sigma \ar[d, "l"] \\ 
    \Sigma \ar[r, "r"'] & \Lambda^2_1
    \arrow["\lrcorner"{anchor=center, pos=0.125, rotate=180}, draw=none, from=2-2, to=1-1]    
  \end{tikzcd}
  \]
  The least topology $J_{\ms L}$ is thus generated by the empty covering on $\emp$, and the single covering family $\set{l,r \colon \Sigma \to \Lambda^2_1}$. Since we have the pushout above, this topology is again subcanonical. 
  
  In this case, it is not hard to give an explicit description of a covering family in $\Pos\fp$: a family of maps is a $\ms J_{\ms L}$-covering on $P \in \Pos\fp$ iff it contains a (finite) subfamily of \emph{upward closed subsets} $\set{P_i\subseteq P}_{i\in I}$, where $P \cong \bigcup_{i\in I}P_i$; see a similar calculation for commutative rings in \citet[VIII. 6]{maclane1992sheaves}. Following the terminology for algebraic geometry, this can be denoted as the \emph{Zariski topology}, and the local geometric morphism
  \[ \sh(\wTop,J_{\ms L}) \surj \sh(\Pos\fp,J_{\ms L}) \simeq \mb{Zar}(\mbb D)\text{,} \]
  is over the \emph{Zariski topos} $\mb{Zar}(\mbb D)$ for distributive lattices.
\end{example}

\begin{example}[\AxiomSL]
  An even stronger axiom is the linearity axiom, 
  \[ \top \vdash x \le y \vee y \le x\text{,} \]
  which requires the dual embeddings of the following quotients to form a covering family,
  \[ \Sigma[\ms{x},\ms{y}]/\ms{x} \le \ms{y} \twoheadleftarrow \Sigma[\ms{x},\ms{y}] \surj \Sigma[\ms{x},\ms{y}]/\ms{y} \le \ms{x}\text{.} \]
  We can compute that $\spec \Sigma[\ms{x},\ms{y}] \cong \Sigma^2$, and 
  \[ \pt(\Sigma[\ms{x},\ms{y}]/\ms{x} \le \ms{y}) \cong \Sigma_2 \cong \set{0 < 1 < 2}\text{.} \]
  Again we have a pushout diagram,
  \[
  \begin{tikzcd}
    \Sigma \ar[r, "\partial"] \ar[d, "\partial"'] & \Sigma_2 \ar[d, "l"] \\ 
    \Sigma_2 \ar[r, "r"'] & \Sigma^2
    \arrow["\lrcorner"{anchor=center, pos=0.125, rotate=180}, draw=none, from=2-2, to=1-1]
  \end{tikzcd}
  \]
  where $\partial \colon \Sigma \to \Sigma_2$ maps $\Sigma$ to the end points of $\Sigma_2$, and $l,r \colon \Sigma_2 \hook \Sigma^2$ is the two $\Sigma_2$-chains in $\Sigma^2$. The least topology $\ms J_{\ms{SL}}$ for the linearity axiom is thus generated by the empty covering on $\emp$, and $\set{l,r \colon \Sigma_2 \hook \Sigma^2}$. Similarly, given this pushout, $J_{\ms{SL}}$ will be subcanonical.

  The topos $\sh(\wTop,J_{\ms{SL}})$ is closely related to simplicial sets. This is due to Joyal's result that the category of simplicial sets is the classifying topos of strict bounded linear orders, and thus we have an equivalence 
  \[ \sh(\Pos\fp,J_{\ms{SL}}) \cong \psh(\Delta)\text{.} \]
  This makes $\sh(\wTop,J_{\ms{SL}}) \surj \psh(\Delta)$ a local topos over simplicial sets.
\end{example}

\begin{example}[\AxiomOneCS]\label{exm:model1T}
  Similarly to the case of (\AxiomSL), the 1-truncation axiom requires dual embeddings of the following quotient maps
  \[\begin{tikzcd}
    & {\Sigma[\ms{x},\ms{y}]/\ms{x}\le\ms{y}} \\
    {\Sigma[\ms{x},\ms{y}]/\ms{x}=0} & {\Sigma[\ms{x},\ms{y}]/\ms{x}=\ms{y}} & {\Sigma[\ms{x},\ms{y}]/\ms{y}=1}
    \arrow[two heads, from=1-2, to=2-1]
    \arrow[two heads, from=1-2, to=2-2]
    \arrow[two heads, from=1-2, to=2-3]
  \end{tikzcd}\]
  to be a covering, which translates to the fact that the three inclusions
  \[ d_0,d_1,d_2 \colon \Sigma \hook \Sigma_2 \]
  covers $\Sigma_2$. Again we have a colimit diagram as follows, 
  \[\begin{tikzcd}
    1 & 1 & 1 \\
    \Sigma & \Sigma & \Sigma \\
    & {\Sigma_2}
    \arrow["0"', from=1-1, to=2-1]
    \arrow["1"{description, pos=0.25}, from=1-1, to=2-2]
    \arrow["0"{description, pos=0.25}, from=1-2, to=2-1]
    \arrow["1"{description, pos=0.25}, from=1-2, to=2-3]
    \arrow["0"{description, pos=0.25}, from=1-3, to=2-2]
    \arrow["0", from=1-3, to=2-3]
    \arrow["{d_0}"', from=2-1, to=3-2]
    \arrow["{d_1}"{description}, from=2-2, to=3-2]
    \arrow["{d_2}", from=2-3, to=3-2]
  \end{tikzcd}\]
  which implies $J_{\ms{1cS}}$ is subcanonical.
  
  Similar to the case of (\AxiomSL), $\sh(\wTop,J_{\ms{1cS}})$ is closely related to the topos of \emph{truncated simplicial sets} $\psh(\Delta_{\le 1})$. This is again induced by the fact that the classifying topos for 1-coskeletal distributive lattices is given by 
  \[ \sh(\Pos\fp,J_{\ms{1cS}}) \cong \psh(\Delta_{\le 1})\text{.} \]
  Hence, this again gives us a local geometric morphism
  \[ \sh(\wTop,J_{\ms{1cS}}) \surj \psh(\Delta_{\le 1})\text{.} \]
\end{example}

\subsection{Connection with existing models}\label{subsec:compare}

Finally, it will be instructive to compare the sheaf models for synthetic domain theory constructed by \citet{FIORE1997151} with the models discussed in \zcref{sec:newmodelforsdt}. Recall the topos $\mc H$ constructed in \emph{loc.\ cit.}\ is the following sheaf topos,
\[ \mc H \coloneq \sh(\mb P,J_{\ms{can}})\text{,} \]
where $\mb P$ be the full subcategory of the category $\wCPO$ of $\omega$-cpos consisting of retracts of the Scott's graph model $G$, and $J_{\ms{can}}$ is the canonical topology on $\mb P$. An $\omega$-cpo is simply a poset having joins of countable chains, and morphisms between them are maps preserving joins of countable chains.

Here we observe that the inclusion $\mb P \hook \wTop$ is \emph{fully faithful}, because
\begin{align*}
  &\wCPO(G,G) \\
  &\quad\cong \Pos(P_f(\N),P(\N)) \\
  &\quad\cong \Pos(\N,\Pos(P_f(\N),\Sigma)) \\
  &\quad\cong \sFrm\cp(\Pos(P_f(\N),\Sigma),\Pos(P_f(\N),\Sigma)) \\ 
  &\quad\cong \wTop(G,G)
\end{align*}
The first isomorphism holds because $G$ is the \emph{free} $\omega$-cpo on $P_f(\N)$; the second holds by formal manipulation; the third is due to the fact that $\Pos(P_f(\N),\Sigma)$ is the free $\sigma$-frame over the discrete poset $\N$; the final one is given by the duality between countably presented $\sigma$-frames and $\wTop$. 

This leads us to the question of the exact relationship between $\mc H$ and the sheaf model $\sh(\wTop,J_{\mr{can}})$ we have constructed. A detailed investigation requires us to better understand the categorical properties of $\wTop$, and we leave this for future work.



% \section{Future directions}

% \subsection{Domain theory with QC}

% In this paper we have explained the axioms of synthetic domain theory using the \gls{QC} principle for distributive lattices and $\sigma$-frames. As we have seen from various examples (\zcref{prop:liftingofalgebra,prop:liftofsober}, \zcref{exm:2issober} etc.) the \gls{QC} principle moreover helps with the actual \emph{computation} for operations on domains. Thus, the natural step next is to further develop domain theory within this framework.

% \subsection{Relationship with Taylor's \emph{Abstract Stone Duality}}

% We would be remiss in failing to mention Taylor's framework of \emph{Abstract Stone Duality}~\citep{Taylor2011,TaylorP:insema,TaylorP:sobsc}, which is based on a similar duality between ``algebras'' and ``spaces''. One apparent difference in methodology is that in Taylor's case, the duality is of the form $\Sigma^{(-)} \dashv \Sigma^{(-)}$ and is moreover required to be monadic. In our setting, the spectrum of an algebra consists not in \emph{all} functions from the algebra into the dualising object but rather only the ones that preserve the structure of observational algebras—the \emph{primes} in Taylor's terminology.

% Our own adjunction $\opens\dashv\spec$ shall not be monadic, but a natural step toward understanding the relationship between abstract Stone duality and synthetic domain theory in the language of \gls{QC} might be to restrict $\opens\dashv\spec$ to an adjunction involving certain more well-behaved subuniverses of $\Set$. We leave this investigation for future work.


\chapter{Coherent arithmetic and synthetic computability}


\chapter{Decidable, atomless Boolean algebra and synthetic probability}\label{ch:BAandprobability}

%%%%%%%%%%%%%%%%%%%%%%%%%%%%%%%%%%%%%%%%%%%%%%%%%%%%%%%%%%%%%%%%%%%%%%%%%%%%%%%%
%% References:
%%
% If you include some work not referenced in the main text (e.g. using \nocite{}), consider changing "References" to "Bibliography".
%

% \renewcommand to change default "Bibliography" to "References"
\renewcommand{\bibname}{References}
\cleardoublepage
\phantomsection
\addcontentsline{toc}{chapter}{References}
\bibliographystyle{plainnat}
\bibliography{thesis}

\appendix

\chapter{Preliminaries}\label{ch:preliminary}

In this chapter we set up basic notation convention for this document, and collect various notions and constructions that will be useful later in the file.

\section{Notation and terminology convention}\label{sec:notation}

For any $n\in\N$, we use $\mb n$ to denote the finite cardinal containing exactly $n$ elements. In particular, $\mb 0$ and $\mb 1$ will be identified with the proposition \emph{falsum} and \emph{verum}, respectively.

For any set $X$, we write $\pss{X}$ for the proposition that $X$ is inhabited, viz.\ $\ex{x}{X}\top$.

We reserve symbols $\mc A,\mc B,\mc C$ for small categories. Fixed large categories are often denoted by bold face symbols, like $\Set$, $\Fin$, etc. We will reserve the sf font for 2-categories. For instance, $\Catt$ will denote the \emph{2-category} of categories, while $\Cat$ will be the \emph{1-category} of categories and functors. Other common 2-categories include $\Lex$, the 2-category of small lex categories, and $\Top$, the 2-category of topoi.

We use $\yon$ to denote the Yoneda embedding. More explicitly, we use the symbol for both the covariant and contravariant Yoneda embedding, but distinguishing them by subscription or superscription,
\[ \yon_- \colon \mc C \hook \psh(\mc C), \quad \yon^- \colon \mc C\op \hook [\mc C,\Set]. \]

In our constructive setting, by an \emph{equivalence} between categories we mean a functor $F \colon \mc A \to \mc B$ that is both \emph{fully faithful} and \emph{essentially surjective}. This is sometimes also referred to as a \emph{weak equivalence} between categories. For most practical usage, the notion of weak equivalence suffices, since one can show that all relevant categorical properties are preserved under weak equivalence. 


\begin{remark}[Comparison with homotopy type theory]
  This document has adopted the language of the old-fashioned Bishop-style constructive set theory, instead of currently more trending homotopy type theory (\gls{HoTT}), as described in the book~\cite{hottbook}. In particular, using the notion of \emph{univalent categories} developed in~\citet{AHRENS_KAPULKIN_SHULMAN_2015}, one can in fact show every weak equivalence between univalent categories admits an inverse. However, \ldots
\end{remark}

\section{Dominance}\label{sec:prelim-dominance}

We recall the notion of \emph{dominance} in the sense of~\citet{rosolini1986continuity}:

\begin{definition}[Dominance]\label{def:dominance}
  A dominance $\Sigma$ is a subuniverse of $\pp$ that is closed under dependent sums, i.e.\ $\mb 1\in\Sigma$, and for any $\varphi\in\Sigma$ together with a family $\psi \colon \varphi \to \Sigma$, the dependent sum $\sum_{x\in\varphi}\psi(x)$ also belongs to $\Sigma$.
\end{definition}

\section{Classical propositions}

Since we work in general in a constructive environment, propositions satisfying certain classical properties deserve additional attention:

\begin{definition}[Decidable and classical propositions]
  We say a proposition $p \in \pp$ is \emph{decidable}, if $p = 0 \vee p = 1$. We say $p$ is \emph{classical}, if $\neg\neg p = p$. 
\end{definition}

\begin{remark}[The subuniverse of decidable and classical propositions]
  Note that the set of decidable propositions is given by $\mb 2 = \set{\mb 0,\mb 1}$. We also denote the set of classical propositions as $\pp_{\dneg}$, and it is evident that there are inclusions $\mb 2 \subseteq \pp_{\dneg} \subseteq \pp$. These are both \emph{dominances} as in~\zcref{def:dominance}.
\end{remark}

\begin{definition}[Decidable and classical sets]\label{def:decidableset}
  For a set $X$, we say it is \emph{decidable} (resp.\ \emph{classical}) if its equalities are decidable (resp.\ classical), i.e.\ for any $x,y\in X$, the proposition $x = y$ is decidable (resp.\ classical).
\end{definition}

\begin{remark}[Caveat on complemented subtypes]\label{rem:complementedsubtype}
  Given a set $X$, it is common in the literature to refer to $S$ as a \emph{decidable subset} of $X$ if, as a function $S \colon X \to \pp$, it factors through the decidable propositions $2 \subseteq \pp$. However, this might be misunderstood as the set $S$ itself is decidable as in~\zcref{def:decidableset}. Hence, we will always refer to such a subset as a \emph{complemented} subset.
\end{remark}

\section{Stack semantics}\label{sec:stacksem}

Another important property for a set is whether it satisfies \emph{choice}. We introduce the following notion of \emph{projectivity} of a set:

\begin{definition}[Projectivity]
  A set $X$ is \emph{projective} if for any surjection $f \colon Y \surj X$ there exists a section.
\end{definition}

\begin{example}[Finite sets are projective]\label{exa:finiteprojective}
  For any $n\in\N$, the finite cardinal $\mb n$ is projective.
\end{example}

\begin{example}[Countable choice]\label{exa:countablechoice}
  The projectivity of $\N$ is simply \emph{choice principle} as discussed in~\zcref{sec:foundation}.
\end{example}

\begin{lemma}\label{lem:projdecsubset}
  If $X$ is projective, then any complemented subset $I$ of $X$ is also projective.
\end{lemma}
\begin{proof}
  This holds since if $I$ is complemented with complement $J$, then for any surjection $Y \surj I$ over $I$, one can construct a surjection $Y + J \surj I + J \cong X$, where now projectivity of $X$ induces a section. Restricting that section to $I$ provides a section of $Y \surj I$. 
\end{proof}

Recall the well-established notion of \emph{internal projectivity} in a topos; see e.g.~\citet[Exer. IV.16]{maclane1992sheaves}. An object $X\in\mc E$ is internally projective if and only if the functor $(-)^X \colon \mc E \to \mc E$ preserves epis.

\begin{proposition}[Projectivity in topoi]\label{prop:projectiveintopos}
  An object $X\in\mc E$ is projective in the stack semantics if and only if it is internally projective.
\end{proposition}
\begin{proof}
  For a proof see~\citet[Thm. 3.3]{nlab:internally_projective_object}.
\end{proof}

\begin{theorem}[Presheaf categories have countable choice]\label{thm:liftingofcountablechoice}
  For any small category $\mc C$, the natural numbers object $\N$ in $\psh(\mc C)$ is (internally) projective.
\end{theorem}
\begin{proof}
  Recall that the natural numbers object $\N$ in $\psh(\mc C)$ is the constant presheaf on $\N$, viz.\ $\N \cong \Delta\N \cong \coprod_{n\in\N}\mb 1$. For any $X\in\psh(\mc C)$ and $c\in\mc C$, the exponential is thus computed by 
  \begin{align*}
    &X^\N(c) \\
    &\quad \cong \mc E(\yon_c \times \N, X) \\
    &\quad \cong \mc E(\coprod_{n\in\N}\yon_c,X) \\
    &\quad \cong \prod_{n\in\N}\mc E(\yon_c,X) \\
    &\quad \cong X(c)^\N
  \end{align*}
  In particular, the exponential $(-)^\N$ is in fact computed \emph{pointwise}. Since colimits in $\psh(\mc C)$ are also computed pointwise, an epi in $\psh(\mc C)$ is simply a pointwise surjection. Then the claim holds if $(-)^\N$ preserves epis in $\Set$, which follows from countable choice (in $\Set$).
\end{proof}

%%%%%%%%%%%%%%%%%%%%%%%%%%%%%%%%%%%%%%%%%%%%%%%%%%%%%%%%%%%%%%%%%%%%%%%%%%%%%%%%
%% Index:
%%
\printthesisindex

\end{document}
